\documentclass[11pt]{article}

\usepackage[T1]{fontenc}
\usepackage[utf8]{inputenc}
\usepackage{lmodern}
\usepackage[margin=1in]{geometry}
\usepackage{amsmath,amssymb,amsthm,mathtools}
\usepackage{enumitem}
\usepackage{booktabs}
\usepackage{array}
\usepackage{longtable}
\usepackage{xcolor}
\usepackage{microtype}
\usepackage{etoolbox}
\usepackage{url}
\usepackage[colorlinks=true,linkcolor=blue!50!black,citecolor=green!35!black,urlcolor=blue!60!black]{hyperref}
\usepackage[capitalize,nameinlink]{cleveref}

% ---------------------------------------------------------------------------
% Theorem environments
% ---------------------------------------------------------------------------
\theoremstyle{plain}
\newtheorem{theorem}{Theorem}[section]
\newtheorem{lemma}[theorem]{Lemma}
\newtheorem{corollary}[theorem]{Corollary}
\newtheorem{proposition}[theorem]{Proposition}
\newtheorem{conjecture}[theorem]{Conjecture}
\newtheorem{question}[theorem]{Question}
\newtheorem{problem}[theorem]{Problem}
\theoremstyle{definition}
\newtheorem{definition}[theorem]{Definition}
\newtheorem{remark}[theorem]{Remark}
\newtheorem{example}[theorem]{Example}
\crefname{problem}{Problem}{Problems}
\crefname{question}{Question}{Questions}

% The Lean proof is the proof; the text gives a faithful outline of it.
\newenvironment{outline}[1][Outline of the formal proof]{\begin{proof}[#1]}{\end{proof}}
% Results that are not formalized get complete proofs, labelled as such.
\newenvironment{informalproof}{\begin{proof}[Proof (not formalized)]}{\end{proof}}

% ---------------------------------------------------------------------------
% Lean references
%   \lean{name}            a declaration, relative to the namespace BoundedNesting
%   \formal{file}{names}   the formal counterpart of a numbered result: a file relative to
%                          lean/BoundedNesting/ and a comma-separated list of declarations
%   \partformal{file}{names}{gap}
%                          a numbered result whose formal counterpart covers all but the
%                          part described in {gap}; the gap is proved or cited in the text
%   \notformalized         a numbered result that has no formal counterpart
%   \cited{where}          a result of an earlier paper used without proof: where it is proved,
%                          and whether it is Lean-verified there
% ---------------------------------------------------------------------------
% Lean names are typeset like paths, with line breaks allowed after "_", "." and "/".
\DeclareUrlCommand\leanname{\def\UrlFont{\ttfamily\upshape}}
\renewcommand{\UrlBreaks}{\do\_\do\.\do\/\do\-}
\newcommand{\lean}[1]{\leanname{#1}}
\newcommand{\leanfile}[1]{\leanname{#1}}
\newcommand{\leanlist}[1]{\begingroup\def\leansep{}%
  \renewcommand*{\do}[1]{\leansep\leanname{##1}\def\leansep{, }}\docsvlist{#1}\endgroup}
\newcommand{\formal}[2]{\par\noindent{\small\upshape\raggedright\textsf{Lean:} \leanlist{#2}\unskip
  \nobreak\hfil\penalty50\hskip1em\hbox{}\nobreak\hfill[\leanfile{#1}]\parfillskip=0pt\par}%
  \smallskip\noindent\ignorespaces}
\newcommand{\partformal}[3]{\par\noindent{\small\upshape\raggedright\textsf{Lean (partial):} \leanlist{#2}\unskip
  \nobreak\hfil\penalty50\hskip1em\hbox{}\nobreak\hfill[\leanfile{#1}]\parfillskip=0pt\par}%
  \par\noindent{\small\upshape\textsf{Not formalized:} #3}\par\smallskip\noindent\ignorespaces}
\newcommand{\notformalized}{\par\noindent{\small\upshape\textsf{Not formalized;} a complete proof is given
  below.}\par\smallskip\noindent\ignorespaces}
\newcommand{\cited}[1]{\par\noindent{\small\upshape\textsf{Cited:} #1}\par\smallskip\noindent\ignorespaces}

% ---------------------------------------------------------------------------
% Notation
% ---------------------------------------------------------------------------
\newcommand{\N}{\mathbb{N}}
\newcommand{\Z}{\mathbb{Z}}
\newcommand{\Q}{\mathbb{Q}}
\newcommand{\R}{\mathbb{R}}
\newcommand{\C}{\mathbb{C}}
\newcommand{\F}{\mathbb{F}}
\newcommand{\T}{\mathsf{T}}
\newcommand{\clog}[1]{\lceil \log_2 #1 \rceil}
\newcommand{\abs}[1]{\lvert #1 \rvert}
\newcommand{\card}[1]{\lvert #1 \rvert}
\newcommand{\eps}{\varepsilon}
\newcommand{\prev}{\mathrm{prev}}
\newcommand{\cost}{\mathrm{cost}}
\newcommand{\gn}{\mathrm{gn}}
\newcommand{\Sh}{\mathrm{Sh}}
\newcommand{\cl}{\mathrm{cl}}
\newcommand{\Aut}{\mathrm{Aut}}
\newcommand{\minrk}{\mathrm{minrk}}
\newcommand{\MAIS}{\mathrm{MAIS}}
\newcommand{\Crux}{\mathrm{Crux}}
\newcommand{\TCP}{\mathrm{TCP}}
\DeclareMathOperator{\rank}{rank}
\DeclareMathOperator{\dist}{dist}
\DeclareMathOperator{\poly}{poly}

% ---------------------------------------------------------------------------
% Additions for this paper. Lean declarations are relative to the namespace BoundedNesting, and the
% files named in \formal and \partformal lines are relative to lean/BoundedNesting/.
% ---------------------------------------------------------------------------
\newcommand{\desc}{\mathrm{desc}}
\newcommand{\lev}{\mathrm{lev}}
\newcommand{\Live}{\mathrm{Live}}
\newcommand{\DF}{D_{\mathrm F}}
\newcommand{\DM}{D_{\mathrm M}}
\newcommand{\Dlin}{\mathrm{D}_{\mathrm{lin}}}
\newcommand{\irk}{\operatorname{irk}}
\newcommand{\nullity}{\operatorname{nullity}}
\newcommand{\tpos}{\operatorname{tpos}}
\newcommand{\supp}{\operatorname{supp}}
\newcommand{\rowsp}{\operatorname{rowspace}}
\newcommand{\LinCrux}{\mathrm{LinCrux}}
\newcommand{\Sball}{\mathcal S^{\mathrm{ball}}}
\newcommand{\Str}{\mathcal S^{\mathrm{tr}}}
\newcommand{\DNP}{\mathrm{DNP}}
\newcommand{\sigb}{\mathrm{b}}
\newcommand{\sigf}{\mathrm{f}}


\title{Bounded Nesting and Depth Reduction\\ for Matrix Transposition on Tapes}
\author{Hans Brende \and Claude Opus 5.5 (Anthropic)}
\date{October 2026}

\begin{document}

\maketitle

\begin{abstract}
Can a deterministic multitape Turing machine transpose a $d\times d$ matrix with $\lceil\log_2d\rceil$-bit entries in time
linear in the input length? A negative answer would give the first superlinear lower bound for off-line integer multiplication
on such machines. The fourth paper of this series reduced the impossibility of oblivious linear-time transposers to depth
reduction of the time--space dependency graph by a cheap genie, the trajectory cut problem, plus a coding bound, the crux.
We prove depth reduction for every family of head walks of bounded nesting load: if at every time each of $k$ heads has at most
$\ell$ distinct turning points among its open nested hops of each side, then $\eps T$ edges cut the walk DAG to depth
$2^{O(k^3\ell^3/\eps)}$, whatever the homes, drift, scales and coupling of the heads (Theorem~F). Nested hops with distinct
turning points force revisits, so this covers every machine whose work heads enter each cell boundedly often, and in every
schedule the depth-robust part lives on a set of cells of vanishing capacity. With one arbitrary head and bounded heads at
boundedly many scales, depth reduction still holds (Theorem~M). On the deep-nesting pair, the sharpest open instance, a turn
counter serves every level at depth $2^{m+O(1)}$; on average over shifts, fine levels are blind to the other head for every
label that changes only at turns, while for block labels coarse levels see its cancellations. By Theorem~F, the crux alone, or
the all-code paired-distance question of \emph{Directed Recovery Structures in Tape Computations}, excludes oblivious
linear-time transposers whose work heads enter every cell boundedly often, and the latter also excludes such machines of low
adaptivity. For linear machines, linear depth reduction together with a linear crux excludes oblivious linear-time
transposers, and the linear crux can be stated exactly as relative rigidity of the transposition matrix (Theorem~$\Dlin$). On
the coding side, extending the lift identity of \emph{Covers and Local Generation in Linear Storage Codes} to asymmetric checks,
we show that lifting a binary linear code along a tower of double covers never raises its rate, and, using mod-$2$ Lück
approximation, we construct towers that never become bipartite and have girth tending to infinity, along which the lifted
Clebsch code has rate tending to zero. A scalar linear code on a triangle-free graph of maximum degree four, the kernel of a
matrix with unit diagonal, has rate at most one half when the exceptional vertices of that matrix are few or far apart
(Theorem~SH). Theorem~F with the structure of families, the owner lemma, the strip cut for one walk (Lemma~W), the path form of the lemma that cuts are
labels and the label arithmetic of the turn counter are formalized in Lean~4; all other results, Theorem~M among them, are
proved in full in the text but not formalized. The question remains open; we state the precise problems on which its oblivious case
rests.

\end{abstract}

\begin{quote}\small
\textbf{Contributions.} Hans Brende chose the problem, fixed the verification standard, and directed the research.
Claude Opus 5.5, a large language model developed by Anthropic, developed the mathematics, wrote the Lean
formalization, and wrote this paper under Hans Brende's direction. \Cref{sec:transposers} combines Theorem~F with the
paired-distance theorems and the coding question of \emph{Directed Recovery Structures in Tape Computations} by Hans Brende
and GPT-6 Astra (OpenAI)~\cite{DirectedRecovery2026}, and \cref{sec:crux} builds on the lift identity and the free-group
results of \emph{Covers and Local Generation in Linear Storage Codes}, by the same authors~\cite{Covers2026}.
\end{quote}

\tableofcontents
\clearpage

\section{Introduction}\label{sec:intro}

\subsection{The problem}

Can a deterministic multitape Turing machine transpose a $d\times d$ matrix with $b(d)=\max(1,\clog d)$-bit entries, written
row by row, in time linear in the input length $L=d^2b(d)$? This is the statement~(T) of the first paper of this
series~\cite{FirstPaper2026}. By a reduction of Harvey and van der Hoeven~\cite[Cor.~6.9]{HvdH2025}, a linear-time
multiplication machine would give a linear-time transposer, so a proof of $\neg$(T) would give the first superlinear lower
bound for off-line integer multiplication on multitape Turing machines, and a proof of~(T) would remove the data-movement
obstacle to linear-time multiplication. The question is open, and this paper does not decide it.

For oblivious machines the fourth paper of the series~\cite{FourthPaper2026} reduced $\neg$(T) to two statements of different
kinds (its Theorem~D, \cref{thm:D4} here). \emph{Depth reduction}: the time--space dependency DAG of every oblivious linear-time
run can be cut to bounded depth by a genie that reveals $\eps L$ bits. At the level of head walks this is the trajectory cut
problem TCP: deleting $\eps T$ edges of the walk DAG of $k$ walks of length~$T$ should leave every path with boundedly many long
hops. \emph{The crux}: index codes for the side information of the remaining bounded-depth cones need more than $\eps L$ bits.
The fourth paper proved TCP for home-based work heads, whose legs all begin or end near fixed homes, for every coupling of the
heads (its Theorem~B), and it identified as the sharpest open instance the \emph{deep-nesting pair}: two copies of a walk nested
at $m+1$ levels, shifted against each other in time, with a conjecture (A$'$) predicting depth-robustness as $m$ grows.
\emph{Directed Recovery Structures in Tape Computations}~\cite{DirectedRecovery2026} showed that the cone relations of every
linear-time schedule are paired-distance separated after a paid trimming, and isolated the crux as one machine-free question
\cite[Question~9.1]{DirectedRecovery2026}: do all storage codes on bounded-degree recovery relations with this separation have
rate at most $\frac12+o(1)$?

This paper is about the depth side, and about how it meets the coding side. Its main theorem replaces homes by
\emph{owners}, the turning points that hops are nested around, and proves TCP for all walks of bounded nesting load.

\subsection{Results}

\paragraph{Bounded nesting (\cref{sec:nesting}).} The \emph{owner} of a long hop of a walk is the first time inside it at which
the walk is farthest from the hop's cell. The \emph{owner lemma} (\cref{lem:owner}): if $\ell$ nested hops of one side have
pairwise distinct owners, the cell of the innermost hop is entered at least $(\ell+1)/2$ times. So walks that enter every cell
at most $V$ times have \emph{owner-load} at most $2V-1$: at every time, each head has at most $2V-1$ distinct owners among its open
hops of each side. \emph{Theorem~F} (\cref{thm:F}): for $k$ walks of length~$T$, coupled arbitrarily, and any set of their long
hops of owner-load at most~$\ell$, some $\eps T$ edges cut the DAG of these hops to depth $2^{O(k^3\ell^3/\eps)}$, independently
of homes, drift and scales. The label is the fourth paper's band label run on virtual heads, one per head, side and nesting
depth, with bands given by the heights of families of hops with a common owner. Consequently TCP holds for walks of bounded
entry (\cref{cor:entry-tcp}), and in every schedule the depth-robust part, if any, lives on fewer than $k\eps T$ cells that are
entered more than $\lceil1/\eps\rceil$ times, while the rest is cut to depth $2^{O(k^3/\eps^4)}$ (\cref{cor:capacity}). The
deep part must be allowed to relay information: for walks without pauses, forbidding it would already imply TCP for one head
fewer (\cref{prop:relays}).

\paragraph{One arbitrary head (\cref{sec:onehead}).} \emph{Lemma~W} (\cref{lem:W}) cuts a single walk with a strip label, anchored
to cells: $\eps T$ edges and depth less than $16/\eps^2+2/\eps$. \emph{Theorem~M} (\cref{thm:M}) combines it with Theorem~F:
one arbitrary walk, together with walks of owner-load at most~$\ell$ whose families have heights in $q$ dyadic windows of
width~$h$, is cut to a depth that depends only on $(k,\ell,q,h,\eps)$, not on the arbitrary walk or on~$T$; with one window the
depth is $\poly(k,\ell,1/\eps)$ times the square of the ratio of the largest to the smallest family height (\cref{cor:M1}). So a
depth-robust family needs two heads of unbounded owner-load, or one together with bounded heads at unboundedly many separated
scales (\cref{cor:where}). TCP for $k$ walks is, up to a factor~$2$ in~$k$, depth reduction for DAGs with a spine and $k$
non-crossing families of arcs (\cref{prop:stacks}); two families are planar (\cref{prop:sr2}), and three is the first open case.

\paragraph{The deep-nesting pair (\cref{sec:dnp}).} A \emph{turn counter}, the number of turns of either head so far modulo~$q$,
doubly serves every level of both heads at $\log_2q=m+O(\log(1/\theta))$ and cost at most
$\frac{11}{32}\theta n_{\mathrm h}/(m+1)$, at every shift (\cref{thm:counter,cor:counter-served}). Hence Conjecture~A$'$ needs
the constant $C\ge1$ (\cref{cor:C1}), and the pair is cut at depth $\lceil11\cdot2^{m+1}/((m+1)\eps)\rceil-1$, against the fourth
paper's $\lceil12/\eps\rceil^{m+1}-1$ (\cref{cor:counter-cut}): the open question is exponential against subexponential depth in
the nesting depth, with base at most~$2$. We prove that every level class of the pair has exactly $n_{\mathrm h}/(2(m+1))$ hops
(\cref{lem:dnp-classes}). For labels that change only at turns, fine classes are \emph{blind}: on average over shifts, a class with
$g$ levels below it sees a turn of the other head on at most a fraction $(g+3+\lceil\log_2(g+2)\rceil)/(m+1)$ of its hops
(\cref{cor:blind}).
Coarse classes are \emph{visible}: for the block labels that serve one head cheaply, a fraction of order $(m-k)/(m+1)$ of the
apex hops of every coarse class of the other head sees a cancellation at full size, an other-head increment of at most
$-Q-1$, on average over shifts (Theorem~V, \cref{thm:V}). Blindness leaves a safe form of the conjecture, A$''$, which still gives exponential depth (\cref{conj:App}).

\paragraph{Consequences for transposers (\cref{sec:transposers}).} \emph{Corollary~F}: the crux alone, at the depth given by
Theorem~F, excludes oblivious linear-time transposers whose work heads enter every cell boundedly often (\cref{cor:F}); such
heads may reverse unboundedly often, and we know no unconditional lower bound for such machines. \emph{Corollary~F$'$}: an affirmative answer to
Question~9.1 of \emph{Directed Recovery Structures} excludes them (\cref{cor:Fprime}), and \emph{Corollary~F$''$} extends this to
adaptive machines whose traces have entropy at most $\eta L$ for a fixed $\eta<\frac12$, provided every trace of probability at
least $2^{-\eta L}$ has bounded entry (\cref{cor:Fsecond}). Theorem~M gives the analogous
statement when one work head is unrestricted (\cref{cor:M-machine}).

\paragraph{Linear machines (\cref{sec:linear}).} For machines with a data path that is linear over a field, linear depth-TCP,
in which depth is replaced by linear depth and long paths may cancel, together with a linear crux excludes oblivious
linear-time transposers (Theorem~$\Dlin$, \cref{thm:Dlin}). The linear crux can be stated exactly as relative rigidity: the
transposition matrix must not be the sum of a matrix of rank below $\rho L$ and a matrix supported on the bounded-depth ball
pattern of the schedule (\cref{lem:rigidity}). Relative rigidity alone excludes no machine.

\paragraph{The crux in special linear settings (\cref{sec:crux}).} For binary linear codes lifted from a fixed code along a
tower of double covers, an exact recursion with a twist pairing extends the lift identity of \emph{Covers and Local
Generation}~\cite{Covers2026} from symmetric to asymmetric checks, so lifting never raises the rate. Along towers that become
bipartite every code has rate at most~$\frac12$ from that level on, and along towers pulled back injectively from chains with
trivial intersection in free pro-$2$ groups the rate tends to a universal rate at most~$\frac12$, by mod-$2$ Lück approximation
in the form of Jaikin-Zapirain and Souza; for the Clebsch code this limit is~$0$, as \emph{Covers and Local Generation} also
shows. We construct such towers that never become bipartite and have girth tending to infinity (\cref{thm:eqcrux,prop:eqtower}).
On triangle-free graphs~$G$ of maximum degree four, over every field, the kernel of a matrix~$M$ with unit diagonal and support
in~$G$ has dimension at most $\lfloor N/2\rfloor$ when the vertices at which $M$ violates the reciprocity condition of the fourth
paper's Theorem~S are at most two or pairwise at distance at least three (Theorem~SH, \cref{thm:SH}), and in general dimension at
most $\nu(G)+\card H-\rho_3(H)$, where $H$ is the set of these vertices (\cref{cor:SH}).

\paragraph{The construction side and open problems (\cref{sec:construction,sec:open}).} We list what a linear-time transposer
would have to do, among other things carry unbounded nesting and a deep part of vanishing capacity, and the problems on which the
oblivious case rests.

\subsection{How to read this paper: the formal proofs}

Six results are formalized in Lean, two of them in part: the owner lemma (\cref{lem:owner}), the structure of families
(\cref{lem:families}), Theorem~F (\cref{thm:F}) and Lemma~W (\cref{lem:W}) completely; the lemma that cuts are labels
(\cref{lem:labels}) in its path form, its consequence for walk DAGs being deduced in the text (the formal statements of
Theorem~F and Lemma~W include it for their labels); and the label arithmetic of the turn counter (\cref{thm:counter}). The
formal Theorem~F holds for any choice of owners strictly inside the hops, of which the owners of \cref{sec:prelim-walks} are
one. For the turn counter, part~(1), the hypotheses
of the Lean theorem on the cover, and the counts of turns and hops, which rest on the fourth paper's Lemma~5.4, are proved in
the text. Everything else is
proved in the text and not formalized; this includes Theorem~M, Corollaries F, F$'$, F$''$ and~M1, the capacity corollary, the stack form of TCP,
the corollaries of the turn counter, the class count, blindness, Theorem~V, Theorem~$\Dlin$, the equivariant crux and Theorem~SH.
Each of them comes with a complete proof, with all constants, and each numbered result states its status as described in
\cref{sec:prelim-lean}. Part~(f) of the equivariant crux follows from a lemma of \emph{Covers and Local Generation}, and a short script in
\leanfile{tools/} checks it independently by one exact rank computation over~$\F_2$. Results of earlier papers that our proofs use are cited with their status there. Lean-verified there: the first paper's
encoding lemma, pass bounds and coding-gap theorems; and the fourth paper's Lemma~Zq, Theorem~S with its matching form, and
rate bounds for bipartite graphs and for the Clebsch graph. With written proofs there: the fourth paper's Theorem~D,
Corollary~B, local recursion, cone lemma, long-edge conversion and Lemma~5.4; Proposition~9.2 of \emph{Directed Recovery
Structures}; and Lemma~6.1 and Theorem~6.2 of \emph{Covers and Local Generation}. Partly formalized there: the paired-distance
theorem and the paid trimming of \emph{Directed Recovery Structures} (its Theorem~3.3 and Corollary~3.8), whose extension to a
frozen adaptive trace is argued in prose (its Remark~3.9). Numerical observations appear only in remarks labelled as numerical,
heuristics are labelled as such, and claims with only a proof sketch appear only in remarks labelled as such; none of them is
used in proofs.

\subsection{Related work}

\paragraph{The series.} The first paper formulated (T), proved in Lean pass and motion bounds, and reduced $\neg$(T) to the Tape
Coding Conjecture~\cite{FirstPaper2026}; the second excluded, in its models, local coding schemes with low-treewidth interaction graphs
and affine limiting motion~\cite{SecondPaper2026}; the third related bit and word transposition~\cite{ThirdPaper2026}; the fourth proved
Theorems D and~B, curved-motion theorems and storage-code bounds~\cite{FourthPaper2026}. \emph{Directed Recovery Structures in
Tape Computations}~\cite{DirectedRecovery2026} proved the paired-distance structure of native cone relations and posed
Question~9.1, and \emph{Covers and Local Generation in Linear Storage Codes}~\cite{Covers2026} studied linear codes along covers.
Memos to the fourth paper proved the half bound for graphs of circular chromatic number below three~\cite{CircularMemo2026},
subcubic bounds beyond scalar recovery~\cite{SubcubicMemo2026}, and the equivalence of TCP with information-cheap depth
reduction~\cite{InfoDepthMemo2026}. This paper is the eighth of the series. Its companion, the seventh, \emph{Nested
Elimination and Half-Rate Bounds for Storage Codes}~\cite{SeventhPaper2026}, published together with it, proves on the coding side that every storage code on a triangle-free graph of maximum degree three,
and on several classes of directed relations without transitive triangles, has rate at most~$\frac12$, that the guessing number
equals the fractional matching number on such graphs, and that the broadcast rate of index coding there is the number of vertices
minus the fractional matching number.

\paragraph{Depth reduction and depth-robust graphs.} The route from small depth to lower bounds goes back to Valiant's depth
reduction for circuits~\cite{Valiant1977}, and the computation-graph method for multitape machines to Hopcroft, Paul and
Valiant~\cite{HPV1977}. Edge-deletion depth reduction was tied to multitape machines by Paul and
Reischuk~\cite{PaulReischuk1980}, and separators of $k$-page graphs, which contain tape computation graphs, by Galil, Kannan and
Szemerédi~\cite{GKS1989separators}, who also constructed $3$-pushdown graphs with nearly linear
separators~\cite{GKS1989pushdown}. Paul, Pippenger, Szemerédi and Trotter separated deterministic from nondeterministic linear
time with segregators of the computation graph~\cite{PPST1983}, and Santhanam related segregators of multi-pushdown graphs to
time and space~\cite{Santhanam2001}; depth reduction at linear cost deletes a fixed fraction of the edges but must leave
constant depth. Depth-robust graphs go back to a question of Stoß studied by Erdős, Graham and Szemerédi~\cite{EGS1975}, and their
pebbling cost was related to depth-robustness by Alwen, Blocki and Pietrzak~\cite{ABP2017}. Expanders with $O(1)$-pushdown and
three-page layouts exist~\cite{DvirWigderson2010,BourgainYehudayoff2013,DSW2016}. We found no work on reducing DAGs of bounded
stack number, or walk DAGs of bounded nesting, to bounded depth by deleting a linear fraction of the edges.

\paragraph{Coding.} The path sums of \cref{sec:linear} are the transfer matrices of algebraic network coding~\cite{KoetterMedard2003};
minrank is that of Haemers and of Bar-Yossef, Birk, Jayram and Kol~\cite{Haemers1979,BBJK2011}; the route from small depth to
index coding goes back to Valiant's depth reduction and Riis's guessing games~\cite{Valiant1977,Riis2007}, and was used in the
language of min-rank by Jukna and Schnitger~\cite{JuknaSchnitger2011}; for static data structures, lower bounds imply
rigidity~\cite{DGW2019}, and for systematic linear data structures the two are equivalent~\cite{NRR2020}; and the relation between
oblivious algorithms and network codes goes back to Adler, Harvey, Jain, Kleinberg and Lehman~\cite{AHJKL2006}. Coding gaps over routing in
undirected networks were found by Braverman and He~\cite{BravermanHe2026} and by Zhang, Li and Li~\cite{ZhangLiLi2026}. The towers of
\cref{sec:crux-towers} build on \emph{Covers and Local Generation}~\cite{Covers2026}, which proves the lift identity for
symmetric binary checks, the monotonicity of normalized nullity under refinements with $p$-group kernels, the half bound for
the inner rank of graph checks, and full inner rank for checks on graphs with a spanning $2$-factor. They also use Cohn's
theory of free ideal rings~\cite{Cohn2006}, the monotonicity argument of Bergeron, Linnell, Lück and
Sauer~\cite{BLLS2014,Kionke2017}, and the mod-$p$ approximation theorems of Jaikin-Zapirain and of Jaikin-Zapirain and
Souza~\cite{JaikinZapirain2024,JaikinSouza2026}, descendants of Lück's approximation theorem~\cite{Luck1994}.

\subsection{Organization}

\Cref{sec:prelim} recalls the model, walk DAGs, the decomposition of the fourth paper and the deep-nesting pair.
\Cref{sec:nesting} proves the owner lemma and Theorem~F, \cref{sec:onehead} Lemma~W and Theorem~M, and \cref{sec:dnp} the results on
the deep-nesting pair. \Cref{sec:transposers} derives the consequences for transposers, \cref{sec:linear} treats linear machines,
and \cref{sec:crux} proves the two cases of the crux. \Cref{sec:construction} collects what a construction would need, and
\cref{sec:open} the open problems. The appendices index the Lean declarations and explain how to build and check the library.

\section{Preliminaries}\label{sec:prelim}

\subsection{Machines and the statement (T)}\label{sec:prelim-machines}

We use the machine model of the first paper~\cite{FirstPaper2026}, as in the fourth~\cite[Section~2.1]{FourthPaper2026}:
a deterministic Turing machine with a finite state set, a read-only two-way input tape, $k\ge1$ two-way work tapes over a
finite alphabet, and a write-only one-way output tape that receives at most one bit per step. All heads start at
cell~$0$, head moves are $-1$, $0$ or~$+1$, and every step writes the cell under each work head. For $d\ge1$ let
$b(d)=\max(1,\clog d)$. The input is a $d\times d$ matrix with $b(d)$-bit entries written row by row, of length
$L=d^2b(d)$, and the machine must write its transpose in the same encoding. Statement~(T) says that some machine does
this within $c(L+1)$ steps, for a constant~$c$ and all~$d$; its formal version,
\lean{Transpose.LinearTimeTransposable}, is fixed in a frozen file of the first paper's library. Bit~$s$ of entry
$(i,j)$ sits at input position $(di+j)b+s$ and at output position $\tpos((di+j)b+s)=(dj+i)b+s$, so on input~$w$ the
output bit of a transposer at position~$\ell$ is $w_{\tpos(\ell)}$~\cite[Section~2.2 and Lemma~2.5]{FirstPaper2026}.

A machine is \emph{oblivious} at an input length if the positions of all heads at every time, the times at which it
appends output bits, and its halting time do not depend on the input of that length. The \emph{trace} of a run is the
sequence of all head moves together with the append times and the halting time.

\subsection{Walks, hops and owners}\label{sec:prelim-walks}

A \emph{walk} of \emph{length}~$T\ge1$ is a map $x\colon\{0,\dots,T-1\}\to\Z$ with $\abs{x(t+1)-x(t)}\le1$; a step with
$x(t+1)=x(t)$ is a \emph{pause}. The work heads of a machine trace walks $x_1,\dots,x_k$, one per work tape, and these may
be coupled arbitrarily. We follow the fourth paper~\cite[Sections~2.2 and~4.4]{FourthPaper2026}.

For a head~$h$ and a time~$u$, let $\prev_h(u)=\max\{t<u:x_h(t)=x_h(u)\}$ when this set is nonempty. The pair
$(\prev_h(u),u)$ is the \emph{hop} of head~$h$ ending at~$u$, at the \emph{cell} $x_h(u)$; a hop is determined by either
of its endpoints. A hop $(p,u)$ is \emph{long} if $u-p\ge2$. The hops of head~$h$ at a cell~$c$ partition the interval
from the first to the last visit of~$c$ into half-open intervals $(p,u]$, so:
\begin{enumerate}[label=(P\arabic*)]
\item for every time~$\tau$ and every cell~$c$, at most one hop of head~$h$ at~$c$ \emph{spans}~$\tau$, that
is, satisfies $p<\tau\le u$.
\end{enumerate}
Let $(p,u)$ be a long hop at cell~$c$. Then $x_h(t)\ne c$ for $p<t<u$, and since the walk moves by at most one cell per step,
either $x_h(t)>c$ for all $t\in(p,u)$ or $x_h(t)<c$ for all of them. The hop is \emph{upper} in the first case and
\emph{lower} in the second; this is its \emph{side}. The \emph{owner} of the hop is the first time $o\in(p,u)$ at which
$\abs{x_h(o)-c}$ is largest, and $\abs{x_h(o)-c}\ge1$ is its \emph{height}. For a walk without pauses, a \emph{turn} is a
time~$t$ with $x(t+1)-x(t)\ne x(t)-x(t-1)$, and the owner of every long hop is a turn: if the height is~$1$, then
$u=p+2$ and $o=p+1$; otherwise the walk is at distance one less than the height at $o-1$ and at $o+1$.

Two upper long hops of one walk never \emph{cross}: there are no upper long hops $(p,u)$, $(p',u')$ with
$p<p'<u<u'$, and the same holds for lower long hops~\cite[Section~5.1]{FourthPaper2026}. (If $p<p'<u<u'$, then
$x(p')>x(p)$ because $p'\in(p,u)$, while $u\in(p',u')$ and $x(u)=x(p)<x(p')$, contradicting $x>x(p')$ on $(p',u')$.
The argument uses only $x>x(p)$ on $(p,u)$, so it holds with pauses.) Hence two long hops of one walk on the same side
are either disjoint as open intervals or strictly nested, $p<p'<u'<u$; and long hops of one side whose open intervals
contain a common time form a chain under inclusion.

\begin{definition}[Entries]\label{def:entries}
\formal{Owner.lean}{entries}
An \emph{entry} of a cell~$c$ by a walk~$x$ is a time~$t$ with $x(t)=c$ and either $t=0$ or $x(t-1)\ne c$: the first time
of a maximal run of visits.
\end{definition}

\begin{definition}[Owner-load, multiplicity]\label{def:owner-load}
\partformal{TheoremF.lean}{OwnerLoadLE}{the multiplicity.}
Let $H'$ be a set of long hops of walks $x_1,\dots,x_k$ of length~$T$. It has \emph{owner-load at most~$\ell$} (per head
and side) if for every time~$t$, head~$h$ and side~$\sigma$, the hops of~$H'$ of head~$h$ on side~$\sigma$ whose open
intervals contain~$t$ have at most $\ell$ distinct owners. The \emph{multiplicity} $\mu(h,c)$ of a pair $(h,c)$ is the
number of times~$t$ with $x_h(t)=c$; the walk $x_h$ enters~$c$ at most $\mu(h,c)$ times.
\end{definition}

Since $\sum_c\mu(h,c)=T$ for each head, at most $kT/(V+1)$ pairs $(h,c)$ have multiplicity greater than~$V$, and the
same bound holds for pairs with more than $V$ entries.

\subsection{Walk DAGs, cuts and labels}\label{sec:prelim-dag}

The \emph{walk DAG} of $x_1,\dots,x_k$~\cite[Section~4.4]{FourthPaper2026} has vertex set $[0,T)=\{0,\dots,T-1\}$ and two
kinds of edges: \emph{state edges} $t\to t-1$ for $1\le t<T$, and \emph{hop edges} $u\to p$, one for every long hop
$(p,u)$ of any head. A pause joins the same two vertices as a state edge and is identified with it, and an edge is
determined by its endpoints: a pair that is a long hop of several heads is a single edge. For a set $H'$ of long hops,
$G(H')$ denotes the DAG of the state edges and of the edges $u\to p$ with $(p,u)\in H'$; for a set~$R$ of pairs
$(h,c)$, $G-R$ denotes $G(H')$ for the set $H'$ of long hops at pairs outside~$R$. A \emph{cut} is a set of edges; it has
\emph{depth at most~$D$} if every path that uses no edge of the cut contains at most $D$ hop edges.

A \emph{label} with $q$ values is a map $\Psi\colon[0,T)\to\{0,\dots,q-1\}$. A hop edge $u\to p$ is \emph{violated} if
$\Psi(u)\le\Psi(p)$, a time $t\ge1$ is a \emph{descent} if $\Psi(t)<\Psi(t-1)$, and $F_\Psi$ is the cut consisting of
the violated hop edges and of the state edges $t\to t-1$ at descents. Its size is the \emph{cost} of~$\Psi$.

\begin{lemma}[Cuts are labels]\label{lem:labels}
\partformal{Labels.lean}{depth_le_of_label}{the last sentence, the consequence for walk DAGs.}
Let $\Psi$ be a map from a set to $\{0,\dots,q-1\}$, and let $f(0),f(1),\dots,f(n)$ be a sequence of points of that set
with some of its steps $i<n$ marked. If $\Psi(f(i+1))\le\Psi(f(i))$ for every $i<n$, and $\Psi(f(i+1))<\Psi(f(i))$ for
every marked~$i$, then at most $q-1$ steps are marked. Consequently, for every label $\Psi$ with $q$ values of a
walk DAG, or of a sub-DAG $G(H')$, the cut $F_\Psi$ has depth at most $q-1$.
\end{lemma}

\begin{outline}[Outline of the formal proof of the statement about sequences]
By induction on~$n$, the number of marked steps plus $\Psi(f(n))$ is at most $\Psi(f(0))$
(\lean{card_marked_add_le}), and $\Psi(f(0))\le q-1$.
\end{outline}

\begin{proof}[Proof of the last sentence (not formalized)]
Let $t_0\to t_1\to\dots\to t_n$ be a path that avoids $F_\Psi$, and mark its hop edges. Along a state edge
$t_i\to t_{i+1}=t_i-1$, the time $t_i$ is not a descent, so $\Psi(t_{i+1})\le\Psi(t_i)$; along a hop edge, which is not
violated, $\Psi(t_{i+1})<\Psi(t_i)$. By the statement about sequences, the path has at most $q-1$ hop edges.
\end{proof}

The lemma is one direction of the fourth paper's lemma ``cuts are labels''~\cite[Lemma~5.2]{FourthPaper2026}, which also
proves the converse: for every~$D$, the smallest size of a cut of depth at most~$D$ equals the smallest cost of a label
with $D+1$ values. All cuts in this paper are built from labels.

\subsection{Depth reduction and the crux}\label{sec:prelim-tcp}

\begin{definition}[TCP and depth-robustness]\label{def:tcp}
\emph{TCP holds} for a family of $k$-tuples of walks if for every $\eps>0$ there is an integer~$D$ such that the walk DAG
of every member, of length~$T$, has a cut of at most $\eps T$ edges and depth at most~$D$. $\TCP(k)$ is the statement
that TCP holds for the family of all $k$-tuples of walks~\cite[Problem~5.1]{FourthPaper2026}. A family is
\emph{depth-robust} if TCP fails for it: for some $\eps>0$ no single~$D$ serves all members.
\end{definition}

For an oblivious machine~$M$ at length~$L$ with halting time~$T$, the \emph{dependency DAG}
$H_M$~\cite[Definition~4.1]{FourthPaper2026} has vertex set $[0,T)$, state edges $t\to t-1$, and cell edges
$t\to\prev_j(t)$ along the work-head walks. Every edge of $H_M$ joins the endpoints of an edge of the walk DAG of the
work heads: a cell edge of length~$1$ is a pause, and a cell edge of length at least~$2$ is a long hop. The input head
contributes no edges. The same definition applies to a single trace~$\tau$ of any machine at length~$L$, with the head walks,
append times and halting time of~$\tau$; we write $H_\tau$ for this \emph{dependency DAG of the trace}. All inputs that
follow~$\tau$ share this schedule, and the fourth paper's long-edge conversion~\cite[Lemma~4.21]{FourthPaper2026}, whose
proof uses only the DAG, applies to~$H_\tau$. A \emph{genie} is a set $\Gamma\subseteq[0,T)$ of revealed steps; it reveals the values
$v_s$, $s\in\Gamma$ (state after the step, symbols written, output bit), with $\beta\ge\beta(M)$ bits per value, at
\emph{cost} $\beta\card\Gamma$. A path $t_0\to\dots\to t_m$ of $H_M$ is $\Gamma$-free if $t_1,\dots,t_m\notin\Gamma$,
and $\Gamma$ has \emph{full depth at most~$D$} if every $\Gamma$-free path has at most $D$ edges. The cone lemma
\cite[Lemma~4.4 and Corollary~4.5]{FourthPaper2026}, with written proofs there, says that output bit~$\ell$, appended at step~$t_\ell$, is a
function of the revealed word $G_\Gamma(w)=(v_s(w))_{s\in\Gamma}$ and of the input bits at the positions $I_\Gamma(\ell)$
read at the ends of $\Gamma$-free paths from~$t_\ell$, and that $\card{I_\Gamma(\ell)}\le(k+1)^{D+1}$ if $\Gamma$ has
full depth at most~$D$.

The fourth paper calls \emph{depth-TCP} the statement that genies of cost $\eps L$ and bounded full depth exist for all
oblivious machines with $k$ work tapes running in time $C(L+1)$~\cite[Definition~4.7]{FourthPaper2026}. A
\emph{side-information system} at size~$d$ is a family $I=(I(\ell))_{\ell\in[0,L)}$ of subsets of $[0,L)$, and an
\emph{index code} for~$I$ with $m$ bits is a map $G\colon\{0,1\}^L\to\{0,1\}^m$ such that $G(w)$ and $w|_{I(\ell)}$
determine $w_{\tpos(\ell)}$. For a class $\mathcal S$ of systems and $\rho\ge0$, $\Crux(\mathcal S,\rho)$ says that
for every $\eta>0$ there is $d_0$ such that every index code for every system of $\mathcal S$ at a size $d\ge d_0$ has at
least $(\rho-\eta)L$ bits; $\mathcal S_{k,C,\beta,D}$ is the class of the systems $I_\Gamma$ of oblivious machines with $k$
work tapes, value width at most~$\beta$ and halting time at most $C(L+1)$ that append exactly $L$ bits, for genies
of full depth at most~$D$~\cite[Definition~4.8]{FourthPaper2026}. Every fibre of an index code for~$I$ is a binary storage
code for the relation $N_I(p)=I(\tpos(p))$, in which every symbol is a function of the symbols in its recovery set
\cite[Remark~4.9(b)]{FourthPaper2026}.

\begin{theorem}[Theorem D of the fourth paper]\label{thm:D4}
\cited{the fourth paper~\cite[Theorem~4.10]{FourthPaper2026}, with a complete written proof there; not formalized.}
If $\mathrm{depth}\text{-}\TCP(k,C,\beta,\eps,D)$ and $\Crux(\mathcal S_{k,C,\beta,D},\rho)$ hold with $0<\eps<\rho$,
then for every machine with $k$ work tapes and value width at most~$\beta$ there is $d_0$ such that for no $d\ge d_0$ is it
oblivious at length $L=d^2b(d)$ with halting time at most $C(L+1)$ and a transposer of all $d\times d$ matrices.
\end{theorem}

\paragraph{Why it matters.} Theorem~D splits a proof that no oblivious machine transposes in linear time into a
structural half, depth-TCP, and a coding half, the crux. \Cref{sec:nesting,sec:onehead,sec:dnp} work on the structural half through cuts of
walk DAGs, \cref{sec:linear,sec:crux} on the coding half, and \cref{sec:transposers} combines the two. An oblivious
linear-time transposer must violate, for its own parameters, one of the two hypotheses.

The structural hypothesis is a statement about walk DAGs. Every application below converts a cut of the walk DAG into a
genie, as in the proof of the fourth paper's Corollary~B.

\begin{lemma}[From cuts to genies]\label{lem:cut-genie}
\notformalized
Let $M$ have $k$ work tapes and value width at most~$\beta$, let $C\ge1$ and $0<\eps\le1$, and put
$P=\lceil4\beta C/\eps\rceil$. Suppose that at length~$L\ge1$ the run of $M$ along a trace of halting time
$T\le C(L+1)$ is given, and that the walk DAG of its work-head walks on $[0,T)$ has a cut of at most
$\eps T/(4\beta C)$ edges and depth less than~$D_1$. Then the dependency DAG of that trace has a genie of cost at most
$\eps L$ and full depth at most $D_1P-1$.
\end{lemma}

\begin{informalproof}
Let $F$ be the cut and $\Gamma_0$ the set of steps~$s$ at which some edge $t\to s$ of~$F$ ends, so
$\card{\Gamma_0}\le\card F$. Every edge of a $\Gamma_0$-free path of the dependency DAG ends outside~$\Gamma_0$, so its
endpoints are not those of an edge of~$F$; read in the walk DAG, the path avoids~$F$, and its edges of length at least~$2$
are long hops. Hence every $\Gamma_0$-free path has at most $D_1-1$ edges of length greater than~$1$. By the long-edge
conversion of the fourth paper~\cite[Lemma~4.21]{FourthPaper2026} (a written proof there), with $\mu=1<P$ and
$D'=D_1-1$, the genie
$\Gamma=\Gamma_0\cup\{t<T:t\equiv P-1\pmod P\}$ has full depth at most $D_1P-1$ and
$\card\Gamma\le\card F+\lfloor T/P\rfloor$. Its cost is at most
$\eps T/(4C)+\beta T/P\le\eps T/(2C)\le\eps(L+1)/2\le\eps L$.
\end{informalproof}

The lemma uses the trace only through its dependency DAG. When the machine is oblivious and transposes, the revealed word
of the genie is an index code with at most $\eps L$ bits for the system $I_\Gamma\in\mathcal S_{k,C,\beta,D_1P}$, by the
cone lemma and correctness, exactly as in the proof of Theorem~D.

\subsection{The deep-nesting pair}\label{sec:prelim-dnp}

We recall the instance that the fourth paper regards as the sharpest open case of $\TCP(2)$, with the parameter $g=1$
throughout. Let $s_0>m\ge0$ be integers, $\Lambda=2^{s_0}$ and $R_i=2^{s_0-i}$ for $0\le i\le m$, so $R_0=\Lambda$
and $R_m\ge2$. A \emph{level-$m$ stroke} of direction $\sigma\in\{-1,1\}$ is a run of $R_m$ steps equal to~$\sigma$. For
$i<m$, a \emph{level-$i$ stroke} of direction~$\sigma$ is a \emph{segment} of $R_i$ steps equal to~$\sigma$ followed by a
\emph{V} of level $i+1$: a level-$(i+1)$ stroke of direction~$-\sigma$ (the \emph{back stroke}) followed by a
level-$(i+1)$ stroke of direction~$\sigma$ (the \emph{forward stroke}). The walk $U=U(s_0,m,1)$ satisfies $U(0)=0$,
is periodic with period
\[
  P=4(m+1)\Lambda,
\]
and its steps over one period are those of four level-$0$ strokes of directions $1,-1,-1,1$
\cite[Definition~5.3]{FourthPaper2026}. A level-$i$ stroke has $n_iR_i$ steps, where $n_i=m-i+1$, and moves the head by
$\sigma R_i$; U has no pauses. The first time of a back stroke is the \emph{apex} time of its V, the first time of its
forward stroke the \emph{valley} time; the V is \emph{attached} at its \emph{base}, the cell where it starts and ends, and
its \emph{valley cell} is at distance $R_i$ from the base, for a V of level~$i$. The turns of~$U$ are exactly the apex and
valley times of all V's, each of the level of its V, and the first times of the second and fourth level-$0$ strokes of
every period, of level~$0$, at the \emph{tips} $\pm\Lambda$~\cite[Lemma~5.4]{FourthPaper2026}, with a written proof
there. A period contains $2^{i+1}$ V's of level~$i$ for $1\le i\le m$, as counted in the proof of the fourth paper's
Proposition~5.13, hence $2^{i+2}$ turns of level~$i$ and two tips; all turn times and turn positions are multiples
of~$R_m$~\cite[Remark~5.23]{FourthPaper2026}.

For an integer shift~$\xi$ the pair is $x_1=U$ and $x_2(t)=x_1(t+\xi)$; the \emph{deep-nesting pair}
$\DNP(s_0,m,1)$ of the fourth paper is the pair for a particular odd shift, the golden shift
\cite[Definition~5.9]{FourthPaper2026}, and we call the pair for any odd shift a deep-nesting pair. A pair $(p,u)$ is a hop
of head~$2$ if and only if $(p+\xi,u+\xi)$ is a hop of head~$1$. For $r\ge1$ the \emph{$r$-fold cover} is the pair on the
cyclic time axis $\Z_T$, $T=rP$; its hops are the residues of the hops on~$\Z$, every time is the earlier and the
later endpoint of exactly one hop of each head, and there are $n_{\mathrm h}=2T$ hops in all. Each head visits every cell of
$[-\Lambda,\Lambda]$ in every period, so every time lies in the half-open interval $(p,u]$ of exactly $2(2\Lambda+1)$ hops of
the two heads, one per head and cell. A \emph{label} on the cover is a map $G\colon\Z_T\to\{0,\dots,q-1\}$; a hop $(p,u)$
is violated if $G(u)\le G(p)$, a time~$t$ is a descent if $G(t)<G(t-1)$, and $\cost(G)$ is the number of violated hops of
both heads (counted per head) plus the number $\desc(G)$ of descents. The \emph{error} of a hop is
$e(p,u)=(G(u)-G(p)-1)\bmod q\in\{0,\dots,q-1\}$.

\begin{lemma}[Lemma Zq of the fourth paper]\label{lem:Zq4}
\cited{the fourth paper~\cite[Lemma~5.15]{FourthPaper2026}; Lean-verified there (\lean{MotionDepthCodes.lemma_Zq},
\lean{MotionDepthCodes.card_descents_eq}).}
On the $r$-fold cover, every label satisfies $q\cdot\cost(G)=n_{\mathrm h}+\sum_{\text{hops}}e+q\cdot\desc(G)$ and
$q\cdot\desc(G)=\sum_{t\in\Z_T}\chi_t$, where $\chi_t=(G(t)-G(t-1))\bmod q$.
\end{lemma}

The \emph{level} of a hop is the least level of a turn of its own head strictly inside it (inside any lift to~$\Z$); the
owner is such a turn, so levels are defined. The hops of head~$h$ of level~$i$ form the \emph{class} $(h,i)$. A label
\emph{$\theta$-serves} level~$i$ for head~$h$ if the class $(h,i)$ has total error at most
$\theta qn_{\mathrm h}/(2(m+1))$, and \emph{doubly $\theta$-serves} it if it $\theta$-serves it for both heads
\cite[Section~5.7]{FourthPaper2026}. By \cref{lem:dnp-classes} below, every class has exactly $n_{\mathrm h}/(2(m+1))$
hops, so this says that the mean error of the class is at most~$\theta q$.

\begin{conjecture}[Conjecture A$'$ of the fourth paper]\label{conj:Aprime}
There are $\theta_0>0$ and constants $C,C'$, and for all $s_0$ and~$m$ an odd shift~$\xi$, such that for all $r$ and~$q$
every label~$G$ with $\cost(G)\le n_{\mathrm h}/8$ on the $r$-fold cover of the pair $x_1=U(s_0,m,1)$,
$x_2=x_1(\cdot+\xi)$, has at most $C\log_2q+C'$ doubly $\theta_0$-served levels~\cite[Conjecture~5.25]{FourthPaper2026}.
\end{conjecture}

By the per-class split of Lemma~Zq~\cite[Corollary~5.16]{FourthPaper2026}, $\cost(G)$ is at least the total error of
all hops divided by~$q$, so a level that is not doubly $\theta_0$-served contributes more than
$\theta_0n_{\mathrm h}/(2(m+1))$ to $\cost(G)$; hence a label of cost at most $\eps n_{\mathrm h}$ with $\eps\le1/8$ and
$\eps<\theta_0/2$ doubly serves more than $(m+1)(1-2\eps/\theta_0)$ levels. Under
Conjecture~A$'$ this forces $q\ge\exp(c(1-2\eps/\theta_0)(m+1)-c')$ for constants $c,c'>0$, and the same statement on the
walk DAG of many periods would make the pairs depth-robust as $m\to\infty$, so that $\TCP(2)$ fails
\cite[Section~5.7]{FourthPaper2026}. The fourth paper's clock labels cut every pair at depth $\lceil12/\eps\rceil^{m+1}-1$
with $\eps T$ edges~\cite[Proposition~5.13]{FourthPaper2026}.

\subsection{Formal status of the results}\label{sec:prelim-lean}

Every numbered result names its formal status in one of four ways. A line
\begin{center}
{\small\textsf{Lean:} \texttt{declaration, \dots}\hfill[\texttt{file}]}
\end{center}
gives the Lean declarations of a fully formalized result. The declarations live in the namespace
\texttt{BoundedNesting} and the files in the directory \leanfile{lean/BoundedNesting/}. The text after such a statement,
labelled \emph{Outline of the formal proof}, is a guide to the formal proof; it is a convenience, not a proof. A line
beginning \textsf{Lean (partial)} marks a result whose formal counterpart covers all but an explicitly named part, which
the text proves; on a definition, it names the notions that have no Lean counterpart. The words \textsf{Not formalized} mark a result with no formal counterpart; it comes with a complete
proof, labelled \emph{Proof (not formalized)}. A line beginning \textsf{Cited} marks a result of an earlier paper of the
series that is used without proof, and says where it is proved and whether it is Lean-verified there. Where the English
statement simplifies, the Lean statement is authoritative. Results from the literature are cited. Numerical observations
appear only in remarks labelled as numerical, heuristics are labelled as such, and claims that have only a proof sketch are
stated in remarks beginning \emph{Proof sketch only}; none of them is used in proofs. Numbered definitions carry a status line
only when they have a Lean counterpart.

The library depends on Lean~4 and Mathlib (both v4.34.1) and on the Lean library of the fourth paper, which imports those
of the first two papers and the repository's shared library \texttt{lean-machine}, all through relative paths. It
contains no \texttt{sorry}, no \texttt{native\_decide} and no unsafe evaluation, and every declaration cited here depends
only on the axioms \texttt{propext}, \texttt{Classical.choice} and \texttt{Quot.sound} (\cref{app:verify}).
\Cref{app:index} lists every declaration of this library cited in the paper.

\section{Bounded nesting: the owner lemma and Theorem F}\label{sec:nesting}

The fourth paper's Theorem~B cuts the walk DAG of any coupling of home-based walks with a lexicographic label: one
digit per occupied group of scales, each counting legs measured from a fixed home~\cite[Theorem~4.20]{FourthPaper2026}.
Its hypothesis, linear leg area about fixed homes, fails when heads turn often, far from every home, on legs much shorter
than their distance from home. This section
replaces homes by owners. \Cref{lem:owner} shows that nested structures with distinct turning points force revisits, and
Theorem~F (\cref{thm:F}) proves depth reduction for every family of walks of bounded \emph{owner-load}
(\cref{def:owner-load}), whatever their homes, drift and scales. Its label is Theorem~B's band label run on \emph{virtual
heads}, one per head, side and nesting depth, with the band of a virtual head given by the height of its current family of
hops instead of a distance from home. \Cref{sec:nest-capacity} derives the consequences at walk level: TCP for walks of
bounded entry, and the confinement of depth-robustness to a set of cells of vanishing capacity.

\subsection{The owner lemma}\label{sec:nest-owner}

\begin{lemma}[Nesting with distinct owners costs visits]\label{lem:owner}
\formal{Owner.lean}{owner_entries, card_crossings_le}
Let $x$ be a walk, $\ell\ge1$, and let $(p_i,u_i)$, $0\le i<\ell$, be upper long hops of~$x$ that are strictly nested,
$p_i<p_{i+1}$ and $u_{i+1}<u_i$, with pairwise distinct owners. Then for every time $T'\ge u_0$, the cell of the innermost
hop $(p_{\ell-1},u_{\ell-1})$ has at least $(\ell+1)/2$ entries at times at most~$T'$. By the reflection $x\mapsto-x$, the
same holds for lower long hops.
\end{lemma}

\begin{outline}
The Lean statement concerns upper hops; the lower case follows by the reflection $x\mapsto-x$, which maps lower long hops
to upper ones with the same owners and does not change entries. In Lean a walk is a map $x\colon\N\to\Z$ with
$\abs{x(t+1)-x(t)}\le1$ (\lean{IsWalk}); a walk of length~$T$ is extended by staying at $x(T-1)$, which adds no entries
before time~$T$. Upper long hops and owners are the structures
\lean{IsUpperHop} and \lean{IsOwner}, and entries and crossings are \lean{entries} and \lean{crossings}. Let $c_i$ be the
cell of $(p_i,u_i)$ and $c=c_{\ell-1}$. Since $p_{i+1}\in(p_i,u_i)$, the cells increase inwards, $c_0<\dots<c_{\ell-1}=c$.
Every owner $o_i$ is a maximum of $x$ over $(p_i,u_i)\supseteq(p_{\ell-1},u_{\ell-1})$, where $x\ge c+1$, so
$x(o_i)\ge c+1$. For $i<\ell-1$, $o_i$ lies outside $[p_{i+1},u_{i+1}]$: at the endpoints $x=c_{i+1}\le c$, and if $o_i$
lay strictly inside $(p_{i+1},u_{i+1})$ it would also be the first maximum over that interval, that is $o_{i+1}$
(\lean{IsOwner.unique}), contrary to distinctness. So the walk crosses the edge between $c$ and $c+1$ at some step in
$(p_i,p_{i+1})$, between $o_i$ and $p_{i+1}$, or in $[u_{i+1},u_i)$, between $u_{i+1}$ and~$o_i$. The innermost hop
crosses it at the steps $p_{\ell-1}$ and $u_{\ell-1}-1$. These $\ell+1$ crossings occur at distinct steps, because the
intervals $(p_i,p_{i+1})\cup[u_{i+1},u_i)$ are pairwise disjoint and disjoint from $[p_{\ell-1},u_{\ell-1})$. Finally,
every crossing of the edge either enters~$c$ or leaves it, each maximal run of visits to~$c$ is entered once and left at
most once, and so the number of crossings is at most twice the number of entries (\lean{card_crossings_le}).
\end{outline}

\begin{corollary}[Bounded entry bounds the owner-load]\label{cor:owner-load}
\notformalized
Let $x_1,\dots,x_k$ be walks of length~$T$ and $H'$ a set of their long hops such that, for every hop of~$H'$ of head~$h$
at a cell~$c$, the walk $x_h$ enters~$c$ at most $V$ times. Then $H'$ has owner-load at most $2V-1$. In particular, if every
walk enters every cell at most $V$ times, the set of all long hops has owner-load at most $2V-1$.
\end{corollary}

\begin{informalproof}
Fix a time~$t$, a head~$h$ and a side, say upper. The hops of~$H'$ of head~$h$ on this side whose open intervals
contain~$t$ form a chain (\cref{sec:prelim-walks}). Choosing one hop for each of their $\ell$ distinct owners gives a strictly
nested chain of upper long hops with distinct owners, and by \cref{lem:owner} with $T'=T-1$ its innermost cell has at
least $(\ell+1)/2$ entries. This cell is the cell of a hop of~$H'$ of head~$h$, so $(\ell+1)/2\le V$. Lower hops are
symmetric.
\end{informalproof}

\subsection{Families and virtual heads}\label{sec:nest-families}

\begin{definition}[Families]\label{def:families}
\formal{Families.lean}{side, sw, fam, famP, famU, famd, famDepth, vcode, famData}
Let $H'$ be a set of long hops of walks $x_1,\dots,x_k$ of length~$T$. For a head~$h$, a side~$\sigma$ and a time~$o$,
the \emph{family} of $(h,\sigma,o)$ is the set of hops of~$H'$ of head~$h$ on side~$\sigma$ with owner~$o$, when it is
nonempty; every hop of~$H'$ lies in exactly one family. All hops of a family~$F$ contain its owner~$o_F$, so they form a
chain. The \emph{span} of~$F$ is its outermost hop $(P_F,U_F)$, at the cell~$c_F$; its \emph{height} is
$d_F=\abs{x_h(o_F)-c_F}\ge1$; its \emph{band} is $\lambda(d_F)$, where $\lambda(d)=\lfloor\log_2d\rfloor$, so
$2^{\lambda(d)}\le d<2^{\lambda(d)+1}$; and its \emph{depth} is the number of families of $(h,\sigma)$ whose span
contains the span of~$F$ as an interval, $F$ included. A \emph{virtual head} is a triple $v=(h,\sigma,i)$ such that some
family of $(h,\sigma)$ has depth~$i$; its families are those of $(h,\sigma)$ of depth~$i$. We write $N$ for the number of
virtual heads.
\end{definition}

\begin{lemma}[Structure of families]\label{lem:families}
\formal{Families.lean}{fam_cell_mem, fam_cell_injOn, card_fam_le, fam_long, famP_ne, fam_no_cross, famDepth_le, vcode_lt, vh_eq, famData_disj, famData_valid}
\formal{Clock.lean}{Fam.Valid.unique_span, Fam.Valid.sum_d_le, Fam.Valid.overlap}
Let $H'$ have owner-load at most~$\ell$.
\begin{enumerate}[label=(\alph*)]
\item The cells of the hops of a family~$F$ are distinct and lie between $c_F$, included, and $x_h(o_F)$, excluded, so
$\card F\le d_F$; moreover $U_F-P_F\ge2d_F$.
\item Distinct families of $(h,\sigma)$ have distinct spans, and their spans do not cross.
\item Every family has depth at most~$\ell$. Hence $N\le2k\ell$.
\item Two distinct families $F,F'$ of one virtual head satisfy $U_F\le P_{F'}$ or $U_{F'}\le P_F$.
\item For every virtual head~$v$ and time~$\tau$, at most one family~$F$ of~$v$ has $\tau\in(P_F,U_F]$, and
$\sum_{F\text{ of }v}d_F\le(T-1)/2$.
\item \emph{(Overlap.)} For every time~$\tau$ and integer $\theta\ge0$, at most $N(2^\theta-1)$ hops
$(p,u)\in H'$ with $p<\tau\le u$ lie in families of band less than~$\theta$.
\end{enumerate}
\end{lemma}

\begin{outline}
In Lean the walks are maps $x_h\colon\N\to\Z$ (\lean{IsWalk}), extended beyond~$T$ as in \cref{lem:owner}; the hops of~$H'$ are
triples $(h,p,u)$ such that $(p,u)$ is a long hop of~$x_h$ before~$T$ (\lean{IsHop}); and the owners may be any assignment of
times strictly inside the hops, the owners of \cref{sec:prelim-walks} being one (\lean{hopOwner_spec}). A hop is upper if the
walk steps up after its start (\lean{side}); seen from its side, that is, after the reflection $x\mapsto-x$ for lower hops
(\lean{sw}), it lies strictly above its cell (\lean{sw_above}), so all hops of a family are upper hops of one walk. Families,
spans, heights and depths are \lean{fam}, \lean{famP}, \lean{famU}, \lean{famd} and \lean{famDepth}; the span is the hop of
the family that starts first, and it also ends last (\lean{fam_nested}, \lean{fam_span}). With the heads numbered
$0,\dots,k-1$ and $\sigma=1$ for the upper side and $0$ for the lower one, a virtual head $(h,\sigma,i)$ is coded by the number
$(i-1)+(2h+\sigma)\ell<2k\ell$ (\lean{vcode}, \lean{vcode_lt}, \lean{vh_eq}). The formal family clock runs
on all $2k\ell$ codes, the unused ones being virtual heads without families, of band~$0$ at all times; so the formal~$N$ is
its bound~$2k\ell$. The data of all families form an instance \lean{famData} of abstract family data (\lean{Fam}), and
\lean{famData_valid} proves the hypotheses \lean{Fam.Valid} of the family clock: every hop lies in the span of its family with
the owner strictly inside, a hop that ends at the end of its span is the span, (d), (a), and spans end before~$T$. Parts (e)
and~(f) are proved from these hypotheses alone (\lean{Fam.Valid.unique_span}, \lean{Fam.Valid.sum_d_le},
\lean{Fam.Valid.overlap}).

(a) (\lean{fam_cell_mem}, \lean{fam_cell_injOn}, \lean{card_fam_le}, \lean{fam_long}.) Let $F$ be upper; the lower case is the same seen from its side. If $(p,u)\supsetneq(p',u')$ are hops of~$F$, then $p'\in(p,u)$, so
$x_h(p')>x_h(p)$: the cells increase strictly inwards. The innermost hop contains~$o_F$, where the walk is above its cell,
so all cells lie in $[c_F,x_h(o_F)-1]$, and $\card F\le d_F$. The walk moves by at most one cell per step and goes from
$c_F$ at $P_F$ to $x_h(o_F)$ at~$o_F$ and back to $c_F$ at~$U_F$, so $o_F-P_F\ge d_F$ and $U_F-o_F\ge d_F$.

(b) (\lean{famP_ne}, \lean{fam_no_cross}.) The span of~$F$ is a hop of~$F$, with owner~$o_F$; equal spans give equal owners
and hence equal families. Since a hop is determined by its start, distinct spans even start at distinct times. Spans are
long hops of one head on one side, which do not cross.

(c) (\lean{famDepth_le}.) Let $t\in(P_F,U_F)$, a nonempty interval by~(a). Every family~$F'$ whose span contains the span of~$F$ has $t$ in the
open interval of its span, a hop of~$H'$ of head~$h$ on side~$\sigma$ with owner~$o_{F'}$, and distinct such families have
distinct owners. So there are at most $\ell$ of them. A virtual head is then a triple $(h,\sigma,i)$ with $i\le\ell$.

(d) (\lean{famData_disj}, \lean{disj_aux}.) By~(b) we may assume $P_F<P_{F'}$. Suppose that $P_{F'}<U_F$. Since the spans do
not cross, $U_{F'}\le U_F$, and $U_{F'}\ne U_F$, since a hop is also determined by its end; so
$(P_{F'},U_{F'})\subsetneq(P_F,U_F)$. Every family whose span contains the span of~$F$ also contains the span of~$F'$, and
$F'$ contains its own span but not that of~$F$; so $F'$ would have larger depth than~$F$ (\lean{famDepth_lt}). Hence
$U_F\le P_{F'}$.

(e) (\lean{Fam.Valid.unique_span}, \lean{Fam.Valid.sum_d_le}.) The first claim is~(d). The half-open spans $(P_F,U_F]$ of the families of~$v$ are therefore disjoint subsets of
$\{1,\dots,T-1\}$, and by~(a) each has at least $2d_F$ elements.

(f) (\lean{Fam.Valid.overlap}.) A hop $(p,u)$ of a family~$F$ of a virtual head~$v$ with $p<\tau\le u$ satisfies $P_F\le p<\tau\le u\le U_F$. By~(e)
at most one family of~$v$ qualifies, and if its band is less than~$\theta$, then by~(a) it has at most
$d_F<2^{\lambda(d_F)+1}\le2^\theta$ hops. Summing over the $N$ virtual heads gives the bound.
\end{outline}

\begin{definition}[The family clock]\label{def:clock}
\formal{Clock.lean}{lam, grp, flr, Fam.band, Fam.group, Fam.occ, Fam.rk, Fam.cnt, Fam.digit, Fam.clock}
Let $H'$ have owner-load at most~$\ell$ and $N\ge1$ virtual heads, numbered $1,\dots,N$.
\begin{itemize}[leftmargin=2em]
\item The \emph{band} of a virtual head~$v$ at time~$t$ is $b_v(t)=\lambda(d_F)$ for the family~$F$ of~$v$ with the
largest $P_F\le t$, and $b_v(t)=0$ if no family of~$v$ starts at or before~$t$. It changes only at the start times
$P_F\ge1$ of the families of~$v$, the \emph{jumps} of~$v$.
\item The \emph{events} are the owners $o_F$ of all families of all virtual heads, one per family, with \emph{scale}
$\lambda(d_F)$.
\item Fix integers $a\ge1$ and $K\ge2$, an \emph{offset} $r\in\{0,\dots,a-1\}$ and $\omega\in\{0,\dots,K-1\}^N$. The
\emph{group} of~$v$ at~$t$ is $G_v(t)=\lfloor(b_v(t)+r)/a\rfloor$, and the \emph{floor} of a group $G\ge0$ is
$\theta(G)=\max(aG-r,0)$; a band $b$ lies in group~$G$ exactly when $\theta(G)\le b<\theta(G+1)$, and the floors at
least~$1$ are the integers congruent to $-r$ modulo~$a$.
\item The \emph{occupied groups} at~$t$ are the distinct values among $G_1(t),\dots,G_N(t)$, listed from the top as
$G_1(t)>\dots>G_{m(t)}(t)$, with floors $\theta_i(t)=\theta(G_i(t))$; here $m(t)\le N$.
\item The \emph{counts} are $n_i(t)=\#\{\text{events at times at most }t\text{ of scale at least }\theta_i(t)\}$.
\item The \emph{family clock} is the label
\[
  \Psi(t)=\sum_{i=1}^{m(t)}\bigl((n_i(t)+\omega_i)\bmod K\bigr)\,K^{N-i}\;\in\;[0,K^N).
\]
\end{itemize}
\end{definition}

So $\Psi(t)$ is the integer whose base-$K$ digits, most significant first, are $\psi_i(t)=(n_i(t)+\omega_i)\bmod K$ for
$i\le m(t)$ and $0$ for $m(t)<i\le N$, and labels compare as these digit strings compare lexicographically. Two bands
$b<b'$ lie in different groups exactly when some floor $\theta\ge1$ satisfies $b<\theta\le b'$; for a
uniformly random offset each integer $\theta\ge1$ is a floor with probability $1/a$. This is the band label of the fourth
paper's Theorem~B~\cite[Definition~4.19]{FourthPaper2026}, with virtual heads in place of heads, family bands in place of
distances from home, and owners in place of leg starts.

\subsection{Theorem F}\label{sec:nest-F}

\begin{theorem}[Theorem F: depth reduction at bounded owner-load]\label{thm:F}
\formal{TheoremF.lean}{theoremF_owner, theoremF, DF}
Let $k,\ell\ge1$, $\eps\in(0,1]$ and $T\ge1$, let $x_1,\dots,x_k$ be walks of length~$T$, coupled arbitrarily, and let $H'$
be a set of their long hops of owner-load at most~$\ell$. Then the DAG $G(H')$ has a cut of at most $\eps T$ edges and depth
less than
\[
  \DF(k,\ell,\eps)=2^{\,2k\ell\left(\lceil64k^2\ell^2/\eps\rceil+5+\lceil\log_2(k^2\ell^2/\eps)\rceil\right)}
  =2^{O(k^3\ell^3/\eps)}.
\]
The bound does not depend on~$T$, on the coupling, on homes or on the scales of the walks, and pauses are allowed.
\end{theorem}

\begin{outline}
The theorem for the owners of \cref{sec:prelim-walks} (\lean{hopOwner}) is \lean{theoremF_owner}. It is an instance of
\lean{theoremF}, which holds for every assignment of owners strictly inside the hops and for every $\eps>0$. Owner-load is
\lean{OwnerLoadLE}, and the hop edges of $G(H')$ are the pairs of \lean{hopPairs}. The cut consists of the state edges
$t\to t-1$ at the times of a subset of $[1,T)$ and of a set of hop pairs, and the depth bound concerns every path
$f(0)\to\dots\to f(n)$ that avoids the cut, as in \lean{strip_cut_depth}. The family clock is defined for abstract family data
(\lean{Fam.clock}) and applied to \lean{famData}, whose hypotheses \lean{Fam.Valid} hold by \cref{lem:families}, with
$N=2k\ell$ as explained there; digit positions and the coordinates of~$\omega$ are numbered from~$0$, so that position $i-1$
carries the digit~$\psi_i$. The expectation is formalized as a count: summed over the $a$ offsets and the $K^N$ choices
of~$\omega$, the violated hops and the descents number at most $K^N\cdot4N^2T+aK^{N-1}2^{a+1}N^2T$
(\lean{Fam.Valid.cost_sum_le}), which is $aK^N$ times the right side of \eqref{eq:F-expect}, and some $(r,\omega)$ does at most
as well as the average (\lean{Fam.Valid.exists_cheap}). The bounds below are lemmas of that count.

Let $r$ and~$\omega$ be uniformly random and independent, and let $F_\Psi$ be the cut of the family clock in $G(H')$
(violated hop edges and state edges at descents). By \cref{lem:labels} it has depth at most $K^N-1$; in Lean, for a label
with $N$ digits in $[0,K)$ compared lexicographically (\lean{Fam.depth_le_of_lex}, \lean{Fam.digit_lt}). We show
\begin{equation}\label{eq:F-expect}
  \mathbb E\card{F_\Psi}\;\le\;\frac{4N^2T}a+\frac{2^{a+1}N^2T}K.
\end{equation}
We bound the number of violated hops of~$H'$, counted per hop, which is at least the number of violated hop edges, and the
number of descents. Throughout, $\eta=(p,u)$ is a hop of~$H'$ of the virtual head~$w$, in the family~$F$, and
$G_\eta=\lfloor(\lambda(d_F)+r)/a\rfloor$ is the group of its band.

\emph{Consistency} (\lean{Fam.Valid.band_of_mem}, \lean{Fam.Valid.band_at_end}). We have $G_w(p)=G_\eta$, and $G_w(u)=G_\eta$
unless $\eta$ is the span of~$F$ and the next family $F'$ of~$w$ starts at $P_{F'}=U_F$. Indeed, $P_F\le p<u\le U_F$, and
by \cref{lem:families}(d) no family of~$w$ starts in $(P_F,U_F)$, so $b_w=\lambda(d_F)$ on $[P_F,U_F)$; the hops of~$F$
other than its span have $u<U_F$, by strict nesting. Call $\eta$ \emph{inconsistent} if $G_w(u)\ne G_\eta$.

\emph{The jumps that change a group} (\lean{Fam.gjumps}, \lean{Fam.Valid.sum_gjumps}). A jump of~$v$ from band~$b$ to
band~$b'$ changes the group of~$v$ with probability at most $\abs{b-b'}/a\le(b+b')/a$ (\lean{card_grp_ne_le}). At the jump that starts a
family~$F''$, the new band is $\lambda(d_{F''})$ and the old band is that of the previous family of~$v$, or the initial
band~$0$; each family occurs at most once as a new band and at most once as an old band (\lean{Fam.Valid.sum_old_le}). So,
since $\lambda(d)<d$, $\sum_{\text{jumps of }v}(b+b')\le2\sum_F\lambda(d_F)\le T$ (\lean{Fam.Valid.sum_jumps_lam}), and
the jumps that change a group number at most $NT/a$ in expectation.

\emph{($\alpha$) Inconsistent hops} (\lean{Fam.incons}, \lean{Fam.Valid.card_incons_le}). We count them all as violated. By
consistency, an inconsistent hop is the span of its family~$F$, and the next family $F'$ of~$w$ starts at $U_F$ with a jump
that changes the group of~$w$; distinct inconsistent hops give distinct such jumps. So their expected number is at most
$NT/a$.

\emph{($\beta$) Consistent hops whose higher groups differ} (\lean{Fam.higher}, \lean{Fam.Valid.exists_charge},
\lean{Fam.Valid.card_higher_le}, \lean{Fam.Valid.sum_higher_le}). Let $\mathcal U(t)$ be the set of occupied groups above
$G_\eta$ at time~$t$, and suppose $\eta$ is consistent and $\mathcal U(p)\ne\mathcal U(u)$; we count it as violated. Some group
$G'>G_\eta$ is occupied at exactly one of the times $p,u$, by a virtual head~$v$ that is not in~$G'$ at the other time;
$v\ne w$, since $w$ is in $G_\eta$ at both times. So $G_v(p)\ne G_v(u)$ and $\max(G_v(p),G_v(u))>G_\eta$. The group of~$v$
changes only at its jumps, and so some jump $\tau\in(p,u]$ of~$v$ moves it between two groups whose larger one, $G^\circ$,
exceeds~$G_\eta$ (\lean{exists_change}: going back from~$u$, the first step at which the group of~$v$ changes between two
groups the larger of which exceeds~$G_\eta$). Let the jump move the band of~$v$ from~$b$ to~$b'$. The floor
$\theta^\circ=\theta(G^\circ)$ satisfies $\min(b,b')<\theta^\circ\le\max(b,b')$, and since
$G^\circ\ge G_\eta+1$, $\lambda(d_F)<\theta(G_\eta+1)\le\theta^\circ$. So $\eta$ spans~$\tau$ and its family has band
less than~$\theta^\circ$. Charge $\eta$ to the jump $(v,\tau)$. By the overlap bound of \cref{lem:families}, at most
$N(2^{\theta^\circ}-1)$ hops are charged to one jump. Given~$r$, $\theta^\circ$ is the largest floor in
$(\min(b,b'),\max(b,b')]$, and a jump that changes no group receives no charge. Over the $a$ offsets at which the jump
changes a group, these floors are distinct integers in $(\min(b,b'),\max(b,b')]$, since the floors at least~$1$ for the
offset~$r$ are the integers congruent to $-r$ modulo~$a$ (\lean{sum_floor_pow_le}); so the expected charge of a jump is at
most $(N/a)2^{M+1}$, where $M=\max(b,b')$ and $2^M\le2^b+2^{b'}$. As above, each family occurs at most once as a new band
and at most once as an old band, so $\sum_{\text{jumps of }v}2^M\le1+2\sum_F2^{\lambda(d_F)}\le1+2\sum_Fd_F\le T$
(\lean{Fam.Valid.sum_jumps_pow}). The expected number of such hops is therefore at most $(2N/a)\sum_vT=2N^2T/a$.

\emph{($\gamma$) The remaining hops} (\lean{Fam.calm}, \lean{Fam.Valid.card_viol_calm_le}, \lean{Fam.Valid.sum_calm_growth_le},
\lean{Fam.Valid.sum_viol_calm_le}). Now $\eta$ is consistent and $\mathcal U(p)=\mathcal U(u)$. Then $G_\eta$ has the same
position $j=\card{\mathcal U(p)}+1$ in the lists of occupied groups at $p$ and at~$u$, and the groups above it, hence
their floors $\theta_i$, $i\le j$, agree. Put $\Delta_i=n_i(u)-n_i(p)$ for $i\le j$; then $\Delta_i\ge0$, because the
floor is the same at both times. Moreover $\Delta_j\ge1$: the owner $o_F\in(p,u)$ is an event of scale
$\lambda(d_F)\ge\theta(G_\eta)=\theta_j$. Say that digit~$i$ \emph{wraps} if
$\lfloor(n_i(p)+\omega_i)/K\rfloor\ne\lfloor(n_i(u)+\omega_i)/K\rfloor$. If no digit $i\le j$ wraps, then
$\psi_i(u)=\psi_i(p)+\Delta_i$ for $i\le j$, the first position $i\le j$ with $\Delta_i>0$ carries a larger digit at~$u$,
all earlier digits agree, and $\Psi(u)>\Psi(p)$ (\lean{Fam.clock_lt_of_no_wrap}). Given~$r$, digit~$i$ wraps with probability
at most $\Delta_i/K$ (\lean{Fam.card_wrap_le}, \lean{Fam.card_wrap_union}), so $\eta$ is violated with probability at most
$\sum_{i\le j}\Delta_i/K$.

We bound $\sum_\eta\sum_{i\le j(\eta)}\Delta_i(\eta)$ over these hops, for every~$r$, by exchanging the order of
summation. $\Delta_i(\eta)$ is the number of events $o'\in(p,u]$ of scale at least~$\theta_i$. Fix an event~$o'$ of
scale~$s$. A pair $(\eta,i)$ counts~$o'$ only if $\eta$ spans~$o'$ and $\theta(G_\eta)\le\theta_i\le s$; the floors
$\theta_i$, $i\le j(\eta)$, are distinct, since $\theta$ is strictly increasing on groups. If $\theta(G_\eta)\le\varphi$,
then $\lambda(d_{F(\eta)})<\theta(G_\eta+1)\le\theta(G_\eta)+a\le\varphi+a$. So, by the overlap bound with $\theta=\varphi+a$,
the pairs
counting~$o'$ number at most
\[
  \sum_{\varphi}N\bigl(2^{\varphi+a}-1\bigr)\;<\;N2^{a+1}2^s,
\]
the sum over the floors $\varphi\le s$, which are distinct integers in $[0,s]$. Each family is the family of one event, and
$\sum_{o'}2^{s(o')}=\sum_v\sum_F2^{\lambda(d_F)}\le\sum_v\sum_Fd_F\le NT/2$. Hence these hops contribute at most
$N2^{a+1}\cdot NT/(2K)=2^aN^2T/K$ in expectation.

\emph{Descents} (\lean{Fam.still}, \lean{Fam.card_desc_still_le}, \lean{Fam.sum_desc_still_le}, \lean{Fam.gchange},
\lean{Fam.Valid.card_gchange_le}). Consider the state edge $t\to t-1$. If no virtual head changes group at step $t-1\to t$, the
occupied groups and their floors are the same at $t-1$ and~$t$, and $n_i(t)\ge n_i(t-1)$ for every~$i$; as in~($\gamma$), a
descent needs a wrap (\lean{Fam.clock_le_of_no_wrap}), which given~$r$ has probability at most $\sum_i(n_i(t)-n_i(t-1))/K$.
Summed over~$t$ this is at most $N$ times the number of events divided by~$K$, and there are at most
$\sum_v\sum_Fd_F\le NT/2$ events; so these descents number at most $N^2T/(2K)$ in expectation. Otherwise some virtual head
has a jump at~$t$ that changes its group, and we count the edge as cut; these edges number at most $NT/a$ in expectation,
by the bound on the jumps that change a group.

\emph{Total} (\lean{Fam.violated_subset}, \lean{Fam.descents_subset}, \lean{Fam.Valid.cost_sum_le}). Adding the four bounds,
$\mathbb E\card{F_\Psi}\le(2NT+2N^2T)/a+(2^a+1/2)N^2T/K$, which gives \eqref{eq:F-expect}. (The formal count bounds the two
terms with~$K$ together: by $N2^{a+1}\sum_Fd_F/K$ and by $N/K$ per family, at most $2^{a+1}N^2T/K$ in all.)

\emph{Parameters} (\lean{theoremF}, \lean{cost_le_eps}, \lean{ceil_le_two_pow_clog}). By \cref{lem:families}(c), $N\le2k\ell$,
so the right side of \eqref{eq:F-expect} is at most $16k^2\ell^2T/a+2^{a+3}k^2\ell^2T/K$. Take
$a=\lceil64k^2\ell^2/\eps\rceil$ and $K=2^{a+4}\lceil k^2\ell^2/\eps\rceil$.
Then the first term is at most $\eps T/4$ and the second at most $\eps T/2$, and some choice of $(r,\omega)$ gives a cut of
at most $\eps T$ edges. Its depth is at most $K^N-1<K^{2k\ell}$, and since $\lceil x\rceil\le2^{\lceil\log_2x\rceil}$
for $x>0$,
\[
  K^{2k\ell}\le2^{2k\ell\left(\lceil64k^2\ell^2/\eps\rceil+4+\lceil\log_2(k^2\ell^2/\eps)\rceil\right)}<\DF(k,\ell,\eps).
\]
In \lean{DF} the ceiling of the logarithm is taken in~$\N$, which agrees with the integer ceiling for $\eps\le1$, since
$k^2\ell^2/\eps\ge1$. Pauses are state edges, and the argument uses only that hops are long and that owners lie strictly
inside their hops.
\end{outline}

\paragraph{Why it matters.} Theorem~F proves the walk form of the structural hypothesis of Theorem~D for every family of
bounded owner-load, for every coupling, and with no reference to positions: neither homes nor leg area enter. Together
with the owner lemma it shows that a depth-robust family of walks, if one exists, must contain heads with unboundedly
many simultaneously open nested structures with distinct turning points, hence, by \cref{lem:owner}, cells entered
unboundedly often. This is the structure of the deep-nesting pair (\cref{sec:prelim-dnp}), which Theorem~F does not
cut at a depth independent of its nesting depth. Theorem~F is consistent with the fourth paper's nesting barrier for band
rules~\cite[Remarks~5.6 and~5.7]{FourthPaper2026}, which concerns rules whose load does not grow with the nesting depth;
the load of the family clock is $N\le2k\ell$.

\subsection{Bounded entry and the capacity of depth}\label{sec:nest-capacity}

\begin{corollary}[Bounded entry implies TCP]\label{cor:entry-tcp}
\notformalized
Let $V\ge1$. TCP holds for every family of $k$-tuples of walks in which every walk enters every cell at most $V$ times,
in particular for walks of multiplicity at most~$V$, with $D(\eps)=\DF(k,2V-1,\eps)$ for $\eps\in(0,1]$.
\end{corollary}

\begin{informalproof}
By \cref{cor:owner-load} the set of all long hops has owner-load at most $2V-1$, and $G(H')$ is the walk DAG; apply
\cref{thm:F}.
\end{informalproof}

The corollary is incomparable with Theorem~B: home-based sweeps $0\to n\to0$, repeated, enter cells near home unboundedly
often, while walks made of V's of varying sizes that drift away from every home enter every cell at most three times
but have superlinear leg area about every cell.

\begin{corollary}[The capacity of depth]\label{cor:capacity}
\notformalized
Let $x_1,\dots,x_k$ be walks of length~$T$, let $\eps\in(0,1]$ and $V=\lceil1/\eps\rceil$, and let $R$ be the set of pairs
$(h,c)$ such that $x_h$ enters~$c$ more than $V$ times. Then $\card R\le kT/(V+1)<k\eps T$; at every time~$\tau$ at most
$\card R$ hops at pairs of~$R$ span~$\tau$; and $G-R$ has a cut of at most $\eps T$ edges and depth less than
$\DF(k,2\lceil1/\eps\rceil-1,\eps)=2^{O(k^3/\eps^4)}$.
\end{corollary}

\begin{informalproof}
Each walk has at most $T$ entries in all, which bounds $\card R$. By (P1), at most one hop at each pair spans a
given time. The DAG $G-R$ is $G(H')$ for the set $H'$ of long hops at pairs entered at most $V$ times, which by
\cref{cor:owner-load} has owner-load at most $2V-1\le3/\eps$; apply \cref{thm:F}.
\end{informalproof}

\paragraph{Why it matters.} Every schedule of linear length splits, with no property of transposition used, into a
\emph{shallow} part, the hops at boundedly entered cells, which a cut of $\eps T$ edges reduces to bounded depth, and a
\emph{scratch}~$R$ of fewer than $k\eps T$ cells, which carries fewer than $k\eps T$ symbols across every time cut.
Paths that avoid the cut may still alternate between at most $\DF$ shallow hops and hops through the scratch. So
depth-robustness of a linear-time schedule can only live in a set of cells of vanishing capacity that its heads visit
over and over. On the deep-nesting family the trade-off between $\card R/T$ and the depth is of order $1/\log D$, and no
better under the walk-DAG form of Conjecture~A$'$ (\cref{prop:dnp-capacity}). The next proposition shows that the
alternation cannot be forbidden at no cost: for walks without pauses, the statement $\mathrm{RF}(k+1)$ that forbids it
implies $\TCP(k)$ in genie form, while $\TCP(k+1)$ in genie form implies $\mathrm{RF}(k+1)$ with $R=\emptyset$.

\begin{proposition}[Relays must be allowed]\label{prop:relays}
\notformalized
Restrict to walks without pauses, and let $\mathrm{RF}(k)$ be the statement: for every $\eps>0$ there are $D(\eps)$ and
$\delta(\eps)$, with $\delta(\eps)\to0$ as $\eps\to0$, such that for all $k$ walks of length~$T$ there are a set~$R$ of pairs
with $\card R\le\delta(\eps)T$ and a set $\Gamma\subseteq[0,T)$ with $\card\Gamma\le\eps T$ such that, for every path
$t_0\to\dots\to t_n$ of the walk DAG with $t_1,\dots,t_n\notin\Gamma$ and every $i$ with $D(\eps)<i\le n$, if
$t_{i-1}\to t_i$ is a hop edge then every hop that it carries lies at a pair of~$R$. Then $\mathrm{RF}(k+1)$ implies
$\TCP(k)$ in genie form: for every $\eps'>0$ there is~$D'$ such that every $k$ walks of length~$T$ have a set
$\Gamma'\subseteq[0,T)$ with $\card{\Gamma'}\le\eps'T$ such that every path whose vertices after the first avoid~$\Gamma'$
has at most $D'$ edges.
\end{proposition}

\begin{informalproof}
For $\eps'\ge1$ take $\Gamma'=[0,T)$. Let $0<\eps'<1$, let $M$ be the least odd integer with $M\ge4/\eps'$, and choose
$\eps\le\eps'/8$ with
$(M-1)\delta(\eps)\le\eps'/8$. Let $x_1,\dots,x_k$ be walks of length $T\ge16/\eps'$. Define an extra walk~$z$ on
$[0,T)$ blockwise: for $t\in[jM,(j+1)M)$, $z(t)=j+((t-jM)\bmod2)$. Inside block~$j$ it alternates between $j$ and $j+1$,
ending at~$j$ because $M$ is odd, and the next block starts at $j+1$; so $z$ has no pauses. Cell~$j$ is visited at most
$M$ times, $z$ has at most $T/M+2$ first visits, and every time that is not a first visit of~$z$ is the later endpoint of a
long hop of~$z$.

Apply $\mathrm{RF}(k+1)$ to $(x_1,\dots,x_k,z)$ at~$\eps$, obtaining $R$ and~$\Gamma$. Let $\Delta$ be the set of times~$s$
at which some path of length greater than $D(\eps)$, with all vertices after the first outside~$\Gamma$, ends. For such a
path ending at~$s$, if $s$ were not a first visit of~$z$, the pair $(z,z(s))$ were not in~$R$ and $\prev_z(s)$ were not
in~$\Gamma$, then appending the hop edge $s\to\prev_z(s)$ would give a path that violates $\mathrm{RF}$ at its last edge.
So every $s\in\Delta$ is a first visit of~$z$ (at most $T/M+2$ times), or the later endpoint of a hop of~$z$ at a pair
of~$R$ (at most $(M-1)\card R\le(M-1)\delta(\eps)T$ times, since each pair of~$z$ carries at most $M-1$ hops), or has
$\prev_z(s)\in\Gamma$ (at most $\card\Gamma\le\eps T$ times, since $\prev_z$ is injective).

Put $\Gamma'=\Gamma\cup\Delta$. A path of length greater than $D(\eps)$ whose vertices after the first avoid~$\Gamma'$
would avoid~$\Gamma$, so its last vertex would lie in~$\Delta\subseteq\Gamma'$, a contradiction. Hence such paths have at
most $D(\eps)$ edges in the walk DAG of $(x_1,\dots,x_k,z)$, and therefore in that of $(x_1,\dots,x_k)$, which has the same
vertices and fewer edges. Finally
$\card{\Gamma'}\le2\eps T+T/M+2+(M-1)\delta(\eps)T\le(\eps'/4+\eps'/4+\eps'/8)T+2\le\eps'T$, using $T\ge16/\eps'$. For
$T<16/\eps'$ every path has fewer than $16/\eps'$ edges, so $D'=\max(D(\eps),\lceil16/\eps'\rceil)$ works.
\end{informalproof}

The conversion is the reason the scratch of \cref{cor:capacity} must be allowed to relay: a busy extra head of bounded
multiplicity and space about $T/M$ turns any relay-free certificate into a genuine genie for the other heads.

\begin{remark}[Depth-robust schedules carry capacity]\label{rem:abp}
Proof sketch only. Alwen, Blocki and Pietrzak showed that a DAG that keeps a path with $d$ vertices after the deletion of
any $e$ vertices has cumulative pebbling cost greater than about $ed$~\cite{ABP2017}; graphs of this kind go back to Erdős,
Graham and Szemerédi~\cite{EGS1975}. Pebbling the walk DAG of $k$ walks in time order keeps, at step~$t$, pebbles on the
current and the previous step and on the earlier endpoints of the hops that span~$t$, at most $\Live(t)+2$ pebbles, where
$\Live(t)$ is the number of hops spanning~$t$. Hence, if a family is depth-robust at $(\eps,D)$ in genie form, then
$\sum_t\Live(t)$ is at least about $\eps TD-2T$: the mean time-cut capacity of a depth-robust schedule grows linearly with
the depth that it forces. \Cref{cor:capacity} bounds the capacity of the deep part from above; on the deep-nesting pairs
both bounds are compatible, with capacity about $T/(m+1)$ and, under Conjecture~A$'$, depth exponential in~$m$.
\end{remark}

\section{One arbitrary head}\label{sec:onehead}

Theorem~F needs every head to have bounded owner-load. One head of arbitrary owner-load is harmless on its own: the walk
DAG of a single walk is planar, and $\TCP(1)$ holds by separators~\cite[Section~5.1]{FourthPaper2026}.
\Cref{lem:W} gives $\TCP(1)$ in the form of a label that is anchored to cells rather than to times, and Theorem~M
(\cref{thm:M}) combines it with the family clock: one arbitrary walk plus walks of bounded owner-load whose families use
boundedly many scales are depth-reducible. \Cref{sec:one-counter} collects what Theorems F and~M, and a stack
formulation of TCP, say about where a depth-robust family must live.

\subsection{Lemma W: a strip label for one walk}\label{sec:one-W}

Fix integers $w\ge2$ and $r\ge1$, a \emph{line residue} $\rho\in\{0,\dots,w-1\}$ and an \emph{offset}
$\omega\in\{0,\dots,r-1\}$. The \emph{lines} are the cells congruent to~$\rho$ modulo~$w$, and the \emph{strips} are the
sets of $w-1$ consecutive cells strictly between consecutive lines. For a walk~$x$ of length~$T$, the \emph{line times} are
the times~$t$ at which $x(t)$ is a line. For a time~$t$ that is not a line time, let $S(t)$ be the strip of~$x(t)$ and
$c(t)$ the number of times $t'\le t$ with $x(t')\in S(t)$; the \emph{strip label} is
$\Psi(t)=(c(t)+\omega)\bmod r$. The \emph{strip cut} consists of the edges of the walk DAG with an endpoint at a line time,
the hop edges $u\to p$ with $\Psi(u)\le\Psi(p)$, and the state edges $t\to t-1$ with $\Psi(t)<\Psi(t-1)$. (The value of
the label at line times does not matter, since every edge at a line time is cut.)

\begin{lemma}[Lemma W: the strip cut for one walk]\label{lem:W}
\formal{OneWalk.lean}{strip_cut, strip_cut_depth}
Let $x$ be a walk of length~$T$ and let $w\ge2$ and $r\ge1$ be integers. For every line residue and offset, every path of
the walk DAG of~$x$ that avoids the strip cut contains at most $r-1$ hop edges. For some line residue and offset, the
strip cut has at most $4T/w+wT/r$ edges.
\end{lemma}

\begin{outline}
In Lean a walk is a map $\N\to\Z$ (\lean{IsWalk}); a walk of length~$T$ is extended by staying at $x(T-1)$, and since
hops, line times, strip counts and cut edges are taken at times below~$T$, the extension changes none of them. Hops are
the pairs of consecutive visits of a cell at distance at least two in time (\lean{IsHop}, \lean{hops});
lines, strips, the strip count and the strip label are \lean{IsLine}, \lean{strip}, \lean{stripCount} and
\lean{stripLabel}, and the cut consists of \lean{cutStates} and \lean{cutHops}. \emph{Depth.} Along a path that avoids the
cut, the label does not increase along state edges and decreases strictly along hop edges, so \cref{lem:labels} bounds the
number of hop edges by $r-1$. \emph{Size.} Some residue has at most $T/w$ line times, by averaging over the $w$ residues
(\lean{exists_line_residue}). A \emph{wrap} is a time $\tau$, not a line time, with $c(\tau)+\omega\equiv0\pmod r$; for
some offset there are at most $T/r$ wraps, again by averaging (\lean{exists_offset}). A hop at a cell that is not a line has
both endpoints in one strip~$S$, and $c$ rises by at least one from $p$ to~$u$, at the visit~$u$; so if the hop is cut, the
count of~$S$ passes a multiple of~$r$ at a visit to~$S$ in $(p,u]$, a wrap of~$S$. At each wrap at most $w-1$ hops of
the strip span it, one per cell, by~(P1) (\lean{hop_unique}). Two adjacent cells that are not lines lie in one strip
(\lean{strip_eq_of_adj}), so a state edge between non-line times is cut only at a wrap. Hence at most $2L+W$ state edges
and at most $L+(w-1)W$ hop edges are cut, where $L$ and $W$ are the numbers of line times and of wraps
(\lean{card_cutStates_le}, \lean{card_cutHops_le}), and $3L+wW\le3T/w+wT/r$.
\end{outline}

For $0<\eps\le4$, with $w=\lceil8/\eps\rceil\ge2$ and $r=\lceil2w/\eps\rceil$, the strip cut has at most $\eps T$ edges
and depth at most $r-1<2w/\eps\le16/\eps^2+2/\eps$. The label is local to the strip of the current cell: the two endpoints of a hop read
the counter of the same strip. So the label is not a band rule in the sense of the fourth paper's nesting
barrier~\cite[Definition~5.5]{FourthPaper2026}, and that barrier does not apply to it.

\subsection{Theorem M}\label{sec:one-M}

Let $x_0$ be any walk of length~$T$; its long hops are called \emph{chords}. Let $H'$ be a set of long hops of further
walks $x_1,\dots,x_k$ of length~$T$, of owner-load at most~$\ell$, with families, bands and events as in
\cref{sec:nest-families}, and put $N=2k\ell$, an upper bound for the number of virtual heads (\cref{lem:families}).
Let $G$ be the DAG on $[0,T)$ with the state edges, all chords and the hop edges of~$H'$; an edge is determined by its
endpoints.

\begin{theorem}[Theorem M: one arbitrary head]\label{thm:M}
\notformalized
Let $q\ge1$ and $h\ge0$ be integers and $0\le\beta_1<\beta_2<\dots<\beta_q$ integers, and suppose that every family~$F$
of~$H'$ is assigned a window $i(F)\in\{1,\dots,q\}$ with $\beta_{i(F)}\le\lambda(d_F)\le\beta_{i(F)}+h$. Let
$\eps\in(0,1]$ and put
\[
  w=2^{\lceil\log_2(8(4+2k)/\eps)\rceil},\qquad K_0=\lceil8(w+1)/\eps\rceil,\qquad
  K=\Bigl\lceil\frac8\eps\Bigl(q+qN+\frac{64q^2N^22^h}\eps+qwN+N^22^{h+1}\Bigr)\Bigr\rceil.
\]
Then $G$ has a cut of at most $\eps T$ edges and depth less than $\DM(k,\ell,q,h,\eps)=K_0K^{2q}$. The bound does not
depend on~$T$, on $x_0$ or on the $\beta_i$, and pauses are allowed.
\end{theorem}

For $q=0$, that is, $H'=\emptyset$, \cref{lem:W} applies. Windows may touch or overlap; any cover of a finite set of bands
by $q$ windows of width~$h$ can moreover be replaced by the greedy cover, in which each window starts at the smallest band
not yet covered, with at most $q$ windows of the same width.

\begin{informalproof}
\emph{Widths and lines.} Put $W_0=w$ and, for $1\le i\le q$, let $W_i$ be the least power of~$2$ that is at least
$\max(\mathrm{lo}_i,w)$, where $\mathrm{lo}_i=64qN2^{\beta_i+h}/\eps$. The $\mathrm{lo}_i$ increase with~$i$, so
$w=W_0\mid W_1\mid\dots\mid W_q$, and $W_i\le2\max(\mathrm{lo}_i,w)\le2(\mathrm{lo}_i+w)$. Draw $\rho$ uniformly from
$\{0,\dots,W_q-1\}$. The \emph{level-$i$ lines} are the cells congruent to~$\rho$ modulo~$W_i$; every level-$i$ line is a
level-$i'$ line for $i'\le i$. A \emph{level-$i$ strip} is the set of cells strictly between two consecutive level-$i$
lines, so it has $W_i-1$ cells and every level-$i$ strip lies inside a level-$(i+1)$ strip. Let
$Z=\{t:x_0(t)\text{ is a level-}0\text{ line}\}$, and for $t\notin Z$ let $\sigma_i(t)$ be the level-$i$ strip of $x_0(t)$.

\emph{Counters and label.} The level-$0$ counter of a level-$0$ strip~$S$ is $c^0_S(t)=\#\{t'\le t:x_0(t')\in S\}$, and for
$i\ge1$ the level-$i$ counter of a level-$i$ strip~$S$ counts \emph{entries} into~$S$,
$c^i_S(t)=\#\{t'\le t:x_0(t')\in S\text{ and }x_0(t'-1)\text{ is a level-}i\text{ line}\}$. For the walks of~$H'$, let
$n_i(t)$ be the number of events at times at most~$t$ of scale at least~$\beta_i$. Choose independent uniform offsets
$\omega_0\in[0,K_0)$ and $\omega_1,\dots,\omega_q,\omega'_1,\dots,\omega'_q\in[0,K)$, and put
$d_0(t)=(c^0_{\sigma_0(t)}(t)+\omega_0)\bmod K_0$, $d_i(t)=(c^i_{\sigma_i(t)}(t)+\omega_i)\bmod K$ and
$e_i(t)=(n_i(t)+\omega'_i)\bmod K$ for $1\le i\le q$. The label $\Psi(t)$, for $t\notin Z$, is the mixed-radix number with
the \emph{static} digit order, most significant first,
\[
  d_q,\ e_q,\ d_{q-1},\ e_{q-1},\ \dots,\ d_1,\ e_1,\ d_0,
\]
so it has $K_0K^{2q}$ values; put $\Psi(t)=0$ for $t\in Z$. The cut~$F$ consists of all edges with an endpoint in~$Z$,
the hop edges $u\to p$ (chords or hops of~$H'$) with $\Psi(u)\le\Psi(p)$, and the state edges at descents. Every edge
at~$Z$ is cut, so the value of $\Psi$ on~$Z$ is never compared, and by \cref{lem:labels} $F$ has depth less than
$K_0K^{2q}$. A digit \emph{wraps} at a step if its count passes a value congruent to minus its offset modulo its radix;
a digit whose count rises by~$\Delta$ at a step wraps there with probability at most $\Delta$ divided by its radix.

\emph{The deciding level of a chord.} Let $\gamma=(p,u)$ be a chord at a cell~$c$ that is not a level-$0$ line (chords at
lines have both endpoints in~$Z$). Its \emph{excursion} $E(\gamma)=\{x_0(t):p<t<u\}$ is a nonempty interval of cells next
to~$c$, on one side of~$c$. Let $\lev(\gamma)$ be the largest $i\in\{0,\dots,q\}$ such that $E(\gamma)$ contains a level-$i$
line, and $0$ if it contains none. If $\lev(\gamma)=i<q$, then $E(\gamma)$ contains no level-$(i+1)$ line, so
$E(\gamma)\cup\{c\}\subseteq\sigma_{i+1}(c)$ (\emph{confinement}).

\emph{(M-a) State edges.} Let $t,t-1\notin Z$. Then $x_0(t-1)$ and $x_0(t)$ are equal or adjacent cells that are not lines,
so they lie in the same strip at every level. No counter of level $i\ge1$ changes at~$t$, since an entry at~$t$ needs
$x_0(t-1)$ on a line; the level-$0$ counter rises by one, and $n_i$ by the number of events at~$t$ of scale at least
$\beta_i$. So $\Psi(t)<\Psi(t-1)$ only if some digit wraps at~$t$.

\emph{(M-b) Chords.} Let $\gamma=(p,u)$ be a chord at a cell~$c$ that is not a line, with $i=\lev(\gamma)$. (i) At $p$ and
at~$u$ every head-$0$ digit reads the counter of the same strip $\sigma_{i'}(c)$, and counters do not decrease. (ii) The
level-$i$ counter of $\sigma_i(c)$ rises strictly in $(p,u]$: for $i=0$ by the visit at~$u$; for $i\ge1$, $E(\gamma)$
contains the nearest level-$i$ line~$y_*$ to~$c$ on the side of the excursion, and if $t_1$ is the last time in
$(p,u)$ with $x_0(t_1)=y_*$, then $x_0(t_1+1)$ is the neighbour of $y_*$ on the side of~$c$, since the walk reaches~$c$
at~$u$ without returning to~$y_*$; this neighbour is a cell of $\sigma_i(c)$, so $t_1+1\in(p,u]$ is an entry into
$\sigma_i(c)$. (iii) If $i<q$, then by confinement $x_0$ visits no
level-$(i+1)$ line at the times $p,\dots,u-1$, so no counter of a level greater than~$i$ changes in $(p,u]$. (iv) The event
counts do not decrease. Hence, if no digit at a position at or above~$d_i$ wraps in $(p,u]$, the most significant digit
whose count changed has risen without wrapping, and $\Psi(u)>\Psi(p)$. By~(iii), $\gamma$ is violated only if
\textup{(c1)} the level-$i$ counter of $\sigma_i(c)$ wraps at a time in $(p,u]$, or \textup{(c2)} some $e_{i'}$ with
$i'>i$ wraps in $(p,u]$.

\emph{(M-c) Hops of~$H'$.} Let $\eta=(p,u)\in H'$ lie in a family~$F$ with window $i=i(F)$, and $p,u\notin Z$. If $x_0$
visits no level-$i$ line during $(p,u)$, it stays in one level-$i$ strip on $[p,u]$, hence in one level-$i'$ strip for every
$i'\ge i$, and enters none of them in $(p,u]$; so the digits $d_{i'}$, $i'\ge i$, do not change. The counts
$n_{i'}$ do not decrease, and $n_i$ rises, because the owner $o_F\in(p,u)$ is an event of scale
$\lambda(d_F)\ge\beta_i$. So $\eta$ is violated only if \textup{(d1)} $x_0$ visits a level-$i$ line during $(p,u)$, or
\textup{(d2)} some $e_{i'}$ with $i'\ge i$ wraps in $(p,u]$.

\emph{Cost.} All expectations are over $\rho$ and the offsets; each violated edge is charged to one of the causes above,
and the classes may overlap, which only over-counts.
\begin{enumerate}[label=(\alph*),leftmargin=2em]
\item \emph{Edges at~$Z$.} A vertex has at most $2$ state edges, $2$ chords and $2k$ hops of~$H'$. Since
$\rho\bmod w$ is uniform, $\mathbb E\card Z=T/w$, and these edges cost at most $(4+2k)T/w\le\eps T/8$.
\item \emph{Descents.} By (M-a), they cost at most $T/K_0+\sum_i\#\{\text{events of scale}\ge\beta_i\}/K\le T/K_0+qNT/(2K)$,
because there are at most $\sum_v\sum_Fd_F\le NT/2$ events (\cref{lem:families}(e)).
\item \emph{(c1).} A wrap at time~$\tau$ of the level-$i$ counter of a strip~$S$ is charged only by chords at cells of~$S$
that span~$\tau$, at most $W_i-1$ of them by~(P1). At level~$0$ there are at most $T/K_0$ wraps in expectation,
costing at most $(w-1)T/K_0$. At level $i\ge1$ every entry follows a visit to a level-$i$ line, there are $T/W_i$ such
visits in expectation, and each entry is a wrap with probability at most $1/K$; so each level $i\ge1$ costs at most $T/K$.
\item \emph{(c2).} Let $e_{i'}$ wrap at~$\tau$. The chords charged to it have level $i<i'$ and span~$\tau$, so by
confinement $x_0(\tau)\in E(\gamma)\cup\{c\}\subseteq\sigma_{i+1}(c)\subseteq\sigma_{i'}(c)$: all of them lie at cells of the
level-$i'$ strip containing $x_0(\tau)$, and by~(P1) there are at most $W_{i'}-1$ of them. Events of scale at least
$\beta_{i'}$ are owners of families of height at least $2^{\beta_{i'}}$, so by \cref{lem:families}(e) there are at most
$NT/2^{\beta_{i'}+1}$ of them. The cost is at most
$\sum_{i'}W_{i'}NT/(2^{\beta_{i'}+1}K)\le\sum_{i'}(\mathrm{lo}_{i'}+w)NT/(2^{\beta_{i'}}K)\le
q(64qN^22^h/\eps+wN)T/K$.
\item \emph{(d1).} Charge $\eta$ to its first visit $\tau\in(p,u)$ to a level-$i$ line. The hops of window~$i$ that
span~$\tau$ lie in families of band at most $\beta_i+h$, so by the overlap bound of \cref{lem:families} there are fewer than
$N2^{\beta_i+h+1}$ of them. There are $T/W_i\le T/\mathrm{lo}_i$ visits to level-$i$ lines in expectation, so the cost is
at most $\sum_iN2^{\beta_i+h+1}T/\mathrm{lo}_i=\eps T/32$.
\item \emph{(d2).} An event~$o'$ of scale~$s$ raises the counts $n_{i'}$ with $\beta_{i'}\le s$, and a wrap of $e_{i'}$ at
$o'$ exposes the hops of windows at most~$i'$ that span~$o'$, whose bands are at most $\beta_{i'}+h$: fewer than
$N2^{\beta_{i'}+h+1}$ of them. Since the $\beta_{i'}$ are distinct integers, $\sum_{\beta_{i'}\le s}2^{\beta_{i'}}<2^{s+1}$,
so the cost is at most $\sum_{o'}N2^{h+2}2^{s(o')}/K\le N^22^{h+1}T/K$, using
$\sum_{o'}2^{s(o')}\le\sum_v\sum_Fd_F\le NT/2$.
\end{enumerate}
\emph{Total.} The edges at~$Z$ cost at most $\eps T/8$; the level-$0$ terms, $T/K_0+(w-1)T/K_0=wT/K_0$, at most
$\eps T/8$; all terms with denominator~$K$ together at most
$(qN+q+64q^2N^22^h/\eps+qwN+N^22^{h+1})T/K\le\eps T/8$; and (d1) at most $\eps T/32$. The sum is less than $\eps T/2$, and
some choice of $\rho$ and the offsets attains at most the expectation.
\end{informalproof}

The chords never see the digits of the bounded walks except through wraps, because the digit positions are static and the
event counts are global. A chord of a strip that $x_0$ has left is decided at the coarsest level whose line it crosses, by
the re-entry into its own strip there; above that level it is confined, which bounds its exposure in~(c2) by a strip
width rather than by the number of live chords. The hops of the bounded walks compare only head-$0$ digits at or above
their window, whose lines are far apart compared with their heights.

\begin{corollary}[Corollary M1: scale-ratio form]\label{cor:M1}
\notformalized
If all family heights of~$H'$ lie in $[d_{\min},d_{\max}]$, then $G$ has a cut of at most $\eps T$ edges and depth
\[
  O\bigl(k\,(N^2d_{\max}/d_{\min}+kN)^2/\eps^6\bigr)=\poly(k,\ell,1/\eps)\cdot(d_{\max}/d_{\min})^2 .
\]
In particular: (i) bounded walks whose families have height at most $d_{\max}$, together with any $x_0$, are cut at depth
$\poly(k,\ell,1/\eps)\,d_{\max}^2$; (ii) if all families except at most two per virtual head have their band in one window of
width~$h$, the depth is bounded in terms of $k$, $\ell$, $h$ and~$\eps$ alone; this covers zigzags of any amplitude,
even growing with~$T$.
\end{corollary}

\begin{informalproof}
Take $q=1$, $\beta_1=\lambda(d_{\min})$ and $h=\lambda(d_{\max})-\lambda(d_{\min})$, so that
$2^h\le2d_{\max}/d_{\min}$; then $K_0=O(k/\eps^2)$ and $K=O((N^22^h+kN)/\eps^2)$, and \cref{thm:M} gives depth less
than $K_0K^2$. For~(i), $d_{\min}\ge1$. For~(ii), the exceptional families are those truncated at the ends of $[0,T)$; with
them the bands are covered by at most $1+2N$ windows of width~$h$, and \cref{thm:M} applies with $q\le1+2N$. A zigzag
between two cells at distance~$A$ has owner-load one per side, and all its families have height~$A$ except the truncated
ones, so $h=0$.
\end{informalproof}

\paragraph{Why it matters.} Theorem~M allows one head to be anything: of unbounded owner-load, nested at every scale,
drifting, deep everywhere. So a depth-robust family must contain either two heads of unbounded owner-load, as the
deep-nesting pair does, or one such head together with bounded heads whose families occupy unboundedly many separated
scales (\cref{cor:where}). For machines this becomes \cref{cor:M-machine}.

\subsection{Where a counterexample to TCP must live}\label{sec:one-counter}

\begin{corollary}[Where depth-robustness lives]\label{cor:where}
\notformalized
Let $\mathcal W$ be a depth-robust family of $k$-tuples of walks. Then the owner-load of the set of all long hops is
unbounded on~$\mathcal W$. If moreover all walks but one have owner-load at most a fixed~$\ell$ on~$\mathcal W$, then for
every fixed~$h$ the least number of windows of width~$h$ covering the family bands of the bounded walks is unbounded
on~$\mathcal W$; in particular the ratio of their largest to their smallest family height is unbounded.
\end{corollary}

\begin{informalproof}
If the owner-load were bounded, \cref{thm:F} would give a single depth bound for each~$\eps$. If, for some~$h$, the bands
of the bounded walks of every member were covered by at most $q$ windows of width~$h$, \cref{thm:M}, applied with the
unbounded walk as~$x_0$, would do the same. A bounded ratio $d_{\max}/d_{\min}$ gives such a cover with $q=1$.
\end{informalproof}

The structure of walk DAGs that these theorems exploit is a stack structure. Call an \emph{$s$-stack DAG} a DAG on
$[0,N)$, with edges determined by their endpoints, consisting of the spine edges $t\to t-1$ and of \emph{arcs} $u\to p$ with
$u-p\ge2$, the arcs being partitioned into $s$ families; in each family no two arcs cross (there are no arcs $(p,u)$,
$(p',u')$ with $p<p'<u<u'$), and every vertex is the left endpoint of at most one arc of each family and the right endpoint
of at most one arc of each family. A cut is a set of edges, and its depth is the largest number of arcs on a path that avoids it. Let $\mathrm{SR}(s)$ be
the statement that for every $\eps>0$ there is $D$ such that every $s$-stack DAG on $[0,N)$ has a cut of at most $\eps N$
edges and depth at most~$D$.

\begin{proposition}[TCP is depth reduction for stack DAGs]\label{prop:stacks}
\notformalized
\begin{enumerate}[label=(\alph*)]
\item The walk DAG of $k$ walks, pauses allowed, is a $2k$-stack DAG.
\item Every $s$-stack DAG $H$ on $[0,N)$ is realized by $s$ walks of length $2N$ without pauses: every arc $(p,u)$ of
family~$i$ is the upper long hop $(2p,2u)$ of walk~$i$. Consequently a cut of size~$c$ and depth~$D$ of their walk DAG yields
a cut of~$H$ of size at most~$c$ and depth at most~$D$.
\end{enumerate}
Hence $\TCP(k)$ for walks without pauses implies $\mathrm{SR}(k)$, with the depth of $\TCP(k)$ at $\eps/2$, and
$\mathrm{SR}(2k)$ implies $\TCP(k)$ for walks with pauses; in particular the forms of TCP with and without pauses are
equivalent up to a factor~$2$ in the number of walks.
\end{proposition}

\begin{informalproof}
(a) The upper long hops of one walk do not cross, nor do its lower long hops (\cref{sec:prelim-walks}). Each time is the left
endpoint of at most one long hop of each walk and the right endpoint of at most one. A pair of times that is a long hop of
several walks is a single edge, kept in one of its families, and pauses are spine edges.

(b) For $0\le t\le N$ let $d_i(t)$ be the number of arcs of family~$i$ that contain~$t$ strictly inside, so
$d_i(0)=d_i(N)=0$. By the degree condition, $\abs{d_i(t+1)-d_i(t)}\le1$: $d_i(t+1)-d_i(t)$ is $1$ if $t$ is the left
endpoint of an arc and $0$ otherwise, minus $1$ if $t+1$ is the right endpoint of an arc and $0$ otherwise. The two
endpoints of an arc $(p,u)$ have the same value of~$d_i$: an arc $(p',u')$ of the family that contains~$p$ strictly inside
has $u'>u$, since $u'<u$ would cross and $u'=u$ would violate the degree condition, so it also contains~$u$, and
conversely. Every time strictly inside the arc has a larger value. Put $x_i(2t)=2d_i(t)$ and
$x_i(2t+1)=\min(x_i(2t),x_i(2t+2))+1$ for $t<N$. Consecutive even values differ by $0$ or~$\pm2$, so $x_i$ is a walk
without pauses. At $2p$ and $2u$ it is at $2d_i(p)$, and strictly between them it is at least $2d_i(p)+1$, so $(2p,2u)$
is an upper long hop of~$x_i$. A spine edge $t\to t-1$ maps to the state edges $2t\to2t-1\to2t-2$ and an arc $u\to p$ to the
hop edge $2u\to2p$; distinct edges of~$H$ have disjoint images. Given a cut~$F$ of the walk DAG, cut an edge of~$H$ if some
edge of its image lies in~$F$. This cut has at most $\card F$ edges, and a path of~$H$ that avoids it maps to a path that
avoids~$F$ with as many long hops as the path has arcs.

For the consequences, a $k$-stack DAG on $[0,N)$, padded with empty families, is realized by $k$ walks of length $2N$, and a
cut of $\eps N=(\eps/2)(2N)$ edges of their walk DAG pulls back; conversely~(a) applies to $k$ walks.
\end{informalproof}

In this language, Theorem~F says that $s$-stack DAGs are depth-reducible when the arcs open at any time in any family come
from boundedly many distinct owners, the owner of an arc being that of its hop in~(b). Forgetting the orientation of the
edges, $s$-stack DAGs are essentially the $s$-pushdown graphs of Galil, Kannan and Szemerédi~\cite{GKS1989separators}. Part~(a) of \cref{prop:stacks} is the
classical fact that a tape is two pushdown stores, so that the computation graphs of multitape machines, the objects of the
computation-graph method of Hopcroft, Paul and Valiant~\cite{HPV1977}, are $O(1)$-pushdown graphs; part~(b) is the converse,
a pushdown store being a tape whose head sits at the height of the stack. The stack view of tape computations was also used
by Paul, Pippenger, Szemerédi and Trotter~\cite{PPST1983} and by Santhanam~\cite{Santhanam2001}. $\mathrm{SR}(s)$ is the
depth analogue of the separator question of Galil, Kannan and Szemerédi, who showed that sublinear separators for
$3$-pushdown graphs, sublinear separators for $k$-page graphs, and subquadratic simulations of nondeterministic two-tape machines
by nondeterministic one-tape machines stand or fall together~\cite{GKS1989separators}. We found no work on reducing DAGs of bounded stack number, or walk DAGs
of bounded nesting, to bounded depth by deleting a linear fraction of the edges.

\begin{proposition}[Two stacks are planar]\label{prop:sr2}
\notformalized
$\mathrm{SR}(2)$ holds, with $D(\eps)=O(\eps^{-2})$.
\end{proposition}

\begin{informalproof}
Draw the vertices of a $2$-stack DAG in order on a line, the spine edges on the line, the arcs of one family above it and
those of the other below it. Arcs of one family do not cross, and arcs that share an endpoint can be drawn without
crossing, so the underlying graph is planar; its maximum degree is at most $2+2\cdot2=6$. By the planar separator theorem of
Lipton and Tarjan~\cite{LiptonTarjan1979}, applied recursively as in~\cite{LiptonTarjan1980}, some set of $O(N/\sqrt M)$
vertices leaves components of at most $M$ vertices. Cutting the at most six edges at each of these vertices cuts
$O(N/\sqrt M)$ edges, and a path that avoids the cut uses no vertex of the separator, except possibly as a path with no edge,
and so stays inside one component and has fewer than $M$ arcs. Taking $M$ of order $\eps^{-2}$ gives the claim.
\end{informalproof}

\paragraph{Why they matter.} \Cref{prop:sr2} recovers $\TCP(1)$, and \cref{prop:stacks} places the open part of TCP at
three or four stacks: $\mathrm{SR}(3)$ is the first open case, and if $\TCP(2)$ fails, as Conjecture~A$'$ with its walk-DAG
analogue predicts for the deep-nesting pairs, then $\mathrm{SR}(4)$ fails. From three stacks on, separators no longer help.
Galil, Kannan and Szemerédi constructed $3$-pushdown graphs whose separators have $\Omega(n/\log^{(s)}n)$ vertices, with
$\log^{(s)}$ the $s$-fold iterated logarithm, for every fixed~$s$~\cite{GKS1989pushdown}. Dvir and Wigderson obtained
$O(1)$-pushdown expanders from monotone expanders, answering a question of Galil, Kannan and
Szemerédi~\cite{DvirWigderson2010}, and Dujmović, Sidiropoulos and Wood built bounded-degree expanders with three-page book
embeddings from the constant-degree monotone expanders of Bourgain and Yehudayoff~\cite{DSW2016,BourgainYehudayoff2013};
each page of that embedding is a non-crossing matching, so with a spine added these expanders are $3$-stack DAGs.
Depth-robustness, however, needs arcs that span many scales, and within
boundedly many non-crossing families such arcs must nest; Theorem~F says that the nesting must then come from unboundedly
many distinct owners.

\begin{remark}[A candidate counterexample]\label{rem:umix-candidate}
Heuristic. Let $x_0=U(s_0,m,1)$, a nested hierarchy, and let $x_1$ be a sequential one: a walk made of V's, each a
segment of $2^{j+1}$ steps followed by a back stroke and a forward stroke of $2^j$ steps, with the exponents~$j$ running
through $s_m,\dots,s_0$ with frequencies proportional to $2^{-j}$, so that every scale gets the same time, as every level
of~$U$ does. Then $x_1$ has owner-load one per side and enters every cell at most three times. Theorem~F cuts $x_1$ alone
at depth $2^{O(1/\eps)}$ and \cref{lem:W} cuts $x_0$ alone at depth $O(\eps^{-2})$, while Theorem~M cuts the pair at depth
$O(4^m)\poly(1/\eps)$, with one window of width $h=m$. If the pair needs depth $2^{\Omega(m)}$ at some fixed~$\eps$, then
$\TCP(2)$ fails already for one deep head and one head of owner-load one. The fourth paper's forced-shape heuristics for the
deep-nesting pair apply with the small V's of~$x_1$ in the role of the second hierarchy; no lower-bound technique is
available, and packing certificates cannot certify depth-robustness~\cite[Proposition~5.14]{FourthPaper2026}.
\end{remark}

\section{The deep-nesting pair}\label{sec:dnp}

The deep-nesting pairs of \cref{sec:prelim-dnp} have two heads of unbounded owner-load as $m$ grows, so they escape
Theorems F and~M, and they remain the sharpest open instance of $\TCP(2)$. This section proves three things about them.
A single turn counter serves every level of both heads at $\log_2q=m+O(\log(1/\theta))$ and cuts the pair at depth
$2^{m+O(1)}/\eps$ (\cref{sec:dnp-counter}); so Conjecture~A$'$ can hold only with $C\ge1$, and the depth question is
exponential against subexponential in~$m$. Fine levels are \emph{blind}: on average over shifts, a class with $g$ levels
below it sees a turn of the other head on at most a fraction $(g+3+\lceil\log_2(g+2)\rceil)/(m+1)$ of its hops, for every
label that changes only at turns (\cref{sec:dnp-blind}). Coarse levels are \emph{visible}: for the block labels that serve one head
cheaply, a fixed fraction of the apex hops of every coarse class of the other head sees a cancellation at full size, an
other-head increment of at most $-Q-1$, on average over shifts (\cref{sec:dnp-visible}). Blindness leaves a weaker form of Conjecture~A$'$, which still gives
exponential depth (\cref{conj:App}).

Throughout, $x_1=U(s_0,m,1)$ and $x_2=x_1(\cdot+\xi)$ for an integer shift~$\xi$, with $s_0>m$, $\Lambda=2^{s_0}$,
$R_i=2^{s_0-i}$, $n_i=m-i+1$ and $P=4(m+1)\Lambda$, as in \cref{sec:prelim-dnp}.

\subsection{Turns and classes}\label{sec:dnp-classes}

Let $B$ be a V of level $k\ge1$ of a head, attached at its base~$b$ to a stroke of direction~$\sigma$, with apex~$\alpha$,
valley~$\nu$ and valley cell $v=b-\sigma R_k$. Its back stroke is a segment from $b$ to~$v$ followed, if $k<m$, by the V
$C_{k+1}$ of level $k+1$ attached at~$v$; for $k+1<l\le m$ let $C_l$ be the V of level~$l$ attached at the end of the segment
of the forward stroke of $C_{l-1}$. By \cite[Lemma~5.4(a)]{FourthPaper2026}, which has a written proof there, $C_l$ is attached at~$v$, has valley cell
$v+\sigma R_l$, stays between $v$ and $v+\sigma R_l$, and ends at time~$\nu$, since $C_{k+1}$ ends the back stroke of~$B$ and
$C_{l+1}$ ends the forward stroke of~$C_l$; its valley time is $\nu-n_lR_l$. We call $C_{k+1},\dots,C_m$ the \emph{chain}
of~$B$, and put $R_{m+1}=0$.

\begin{lemma}[Turns and classes]\label{lem:dnp-classes}
\notformalized
\begin{enumerate}[label=(\alph*)]
\item Each head has $2^{m+3}-6$ turns per period, and the pair has $N_P=2^{m+4}-12$.
\item For $1\le k\le m$, the hops of level~$k$ of a head are exactly the following $2R_k$ hops for each V~$B$ of level~$k$ of
that head, in the notation above: the \emph{apex hops} $(\alpha-d,\alpha+d)$, $1\le d\le R_k$, and the \emph{valley hops} at
the cells $v+\sigma x$, $1\le x\le R_k$, whose later endpoint is $\nu+x$. For $d<R_k$ the only turn of the head in the
half-open interval $(\alpha-d,\alpha+d]$ of an apex hop is~$\alpha$. For $k\le l\le m$ and $R_{l+1}<x<R_l$, the valley hop at
$v+\sigma x$ is
\[
  \bigl(\nu-(m-l)R_l-x,\ \nu+x\bigr),
\]
and the turns of the head in its half-open interval are $\nu$ and the turns of $C_{l+1}$ (only $\nu$ if $l=m$).
\item On the $r$-fold cover each head has exactly $4r\Lambda$ hops of each level $i\in\{0,\dots,m\}$; so every class has
$n_{\mathrm h}/(2(m+1))$ hops, and the hops of~(b) other than those listed with $d<R_k$ or $R_{l+1}<x<R_l$, the
\emph{boundary hops}, are a fraction $(m-k+2)/(2R_k)$ of the class.
\end{enumerate}
\end{lemma}

\begin{informalproof}
(a) The turns are the apex and valley times of the V's, two per V, and two tips per period; a period contains $2^{i+1}$ V's
of level~$i$ (\cref{sec:prelim-dnp}). So each head has $2+\sum_{i=1}^m2^{i+2}=2^{m+3}-6$ turns per period.

(b) The V~$B$ is attached at the end of the segment of a stroke of level $k-1$, which has $R_{k-1}=2R_k$ steps in
direction~$\sigma$ and ends at~$b$ at time~$\alpha$; the back stroke starts with a segment of $R_k$ steps in
direction~$-\sigma$. So for $1\le d\le R_k$ the cell $b-\sigma d$ is visited at $\alpha-d$ and at $\alpha+d$ and not in
between, and for $d<R_k$ the only turn in $(\alpha-d,\alpha+d]$ is~$\alpha$. For the valley hops, the forward stroke of~$B$
starts with a segment from~$v$ to~$b$, so $v+\sigma x$ is visited at $\nu+x$ for $1\le x\le R_k$, and before that time,
after~$\nu$, it is not visited. Let $R_{l+1}<x<R_l$. If $l=k$, the chain stays between $v$ and $v+\sigma R_{k+1}$, so the
last visit before~$\nu$ is by the back segment, at time $\alpha+R_k-x=\nu-(m-k)R_k-x$, using $\nu=\alpha+n_kR_k$. If
$l>k$, the forward segment of~$C_l$ runs from $v+\sigma R_l$ to~$v$ starting at time $\nu-n_lR_l$, and afterwards $C_{l+1}$
stays between $v$ and $v+\sigma R_{l+1}$; so the last visit is at $\nu-n_lR_l+R_l-x=\nu-(m-l)R_l-x$, after the valley
time of~$C_l$. In both cases the turns in the half-open interval are those of $C_{l+1}$ and~$\nu$, because $\nu+x$ lies
before the apex of the V of the forward stroke, at $\nu+R_k$.

The listed hops have level exactly~$k$. The only turn strictly inside an apex hop is~$\alpha$, since the times before~$\alpha$
in the hop belong to the parent's segment and those after it to the back segment. A valley hop at $v+\sigma x$ has its
earlier endpoint at or after~$\alpha$, because the back segment visits its cell at time $\alpha+R_k-x$, and its later
endpoint is $\nu+x\le\nu+R_k$; so its interior lies in the span of~$B$, whose turns have level at least~$k$, and it contains
the turn~$\nu$ of level~$k$. Each listed hop contains exactly one turn of level~$k$, the apex or the valley of~$B$, so it
determines~$B$ and its type; within~$B$, the apex hops are distinguished by~$d$ and the valley hops by their cells. So the
listed hops are distinct, and every head has at least $2^{k+1}\cdot2R_k=4\Lambda$ hops of level~$k$ per period, for
$1\le k\le m$.

(c) A hop has level~$0$ exactly when it contains a tip strictly inside. At a tip~$\tau$ the head is at $\pm\Lambda$, and for
each of the other $2\Lambda$ cells of $[-\Lambda,\Lambda]$ exactly one hop of the head at that cell has $\tau$ in its
half-open interval, hence in its open interval. No hop contains two tips: between consecutive tips the head sweeps all of
$[-\Lambda,\Lambda]$, the second and third level-$0$ strokes of a period sweeping from $\Lambda$ to $-\Lambda$ and the fourth
and first from $-\Lambda$ to $\Lambda$, and a hop at~$\pm\Lambda$ would contain the tip at its own cell. So each head has
exactly $4\Lambda$ hops of level~$0$ per period. With~(b), every level has at least $4\Lambda$ hops per period and head, and
there are exactly $P=4(m+1)\Lambda$ hops per period and head, one ending at every time. Hence every level has exactly
$4\Lambda$, the list of~(b) is complete, and on the $r$-fold cover every class has $4r\Lambda=n_{\mathrm h}/(2(m+1))$ hops.
The boundary hops are $m-k+2$ of the $2R_k$ hops of each V: $d=R_k$ and $x=R_l$ for $k\le l\le m$.
\end{informalproof}

The fourth paper defined service with the normalization $n_{\mathrm h}/(2(m+1))$ per class on the strength of numerical
counts~\cite[Section~5.7]{FourthPaper2026}; \cref{lem:dnp-classes}(c) proves that every class has exactly this size. All
turn times of head~$1$ are multiples of~$R_m$, and those of head~$2$ are congruent to $-\xi$ modulo~$R_m$; so for odd
$\xi$ no time is a turn of both heads.

\subsection{The turn counter}\label{sec:dnp-counter}

\begin{theorem}[Turn counter]\label{thm:counter}
\partformal{TurnCounter.lean}{turnCounter_cost}{part~(1), the hypotheses of the Lean theorem on the cover, and the counts
of \cref{lem:dnp-classes}(a) and \cref{sec:prelim-dnp}, which rest on the fourth paper's Lemma~5.4 and its count of V's.}
Let $\xi$ be any shift, $q\ge3$, and $r$ a multiple of $q/\gcd(q,N_P)$. On the $r$-fold cover let $\mathcal N(t)$, for
$t\in[0,T)$, be the number of turns of either head at times in $[0,t]$, and $G(t)=\mathcal N(t)\bmod q$. Then:
\begin{enumerate}[label=(\arabic*)]
\item every hop $(p,u)$ has error $e=(\kappa-1)\bmod q\le\kappa-1$, where $\kappa\ge1$ is the number of turns of either
head in the lifted interval $(p,u]$;
\item $\desc(G)=rN_P/q$ and $q\cdot\cost(G)\le rN_P(4\Lambda+3)$.
\end{enumerate}
\end{theorem}

\begin{outline}[Outline of the formal proof of the label arithmetic]
The Lean theorem \lean{turnCounter_cost} concerns a finite set of times with a cyclic predecessor bijection and a finite set
of hops arranged in closed chains, as in Lemma~Zq (\lean{MotionDepthCodes.lemma_Zq}). It assumes a label $G$ with values in
$[0,q)$ that advances, modulo~$q$, by $\chi_t\in[0,q)$ at every time~$t$ and by $\kappa_h\ge1$ along every hop~$h$, and it
concludes $q\cdot\cost(G)\le\sum_h\kappa_h+\sum_t\chi_t$ and $q\cdot\desc(G)=\sum_t\chi_t$. Its proof: Lemma~Zq writes
$q\cdot\cost(G)$ as the number of hops plus the sum of the reduced errors plus $\sum_t\chi_t$, and each reduced error
$(G(u)-G(p)-1)\bmod q=(\kappa_h-1)\bmod q$ is at most $\kappa_h-1$.
\end{outline}

\begin{proof}[Proof of~(1), of the hypotheses and of the counts (not formalized)]
The times form the cycle $\Z_T$ under $t\mapsto t-1$, and the hops of each head at each cell form closed chains, each hop
followed by the hop that starts at its later endpoint, since every time is the earlier and the later endpoint of exactly one
hop of each head (\cref{sec:prelim-dnp}). Extend $\mathcal N$ to~$\Z$ by $\mathcal N(t+T)=\mathcal N(t)+rN_P$. Since $r$ is a multiple of $q/\gcd(q,N_P)$, $q$
divides $rN_P$, so $G\equiv\mathcal N\pmod q$ on all of~$\Z$. Hence $G$ advances at each time~$t$ by the number $\chi_t$
of turns at~$t$, at most $2<q$, and along each hop by the number~$\kappa$ of turns in its lifted half-open interval, which
is at least~$1$ because the owner of the hop is a turn strictly inside it (\cref{sec:prelim-walks}). This gives~(1), since
$e\equiv\kappa-1\pmod q$ and $\kappa-1\ge0$. Each turn is counted once in $\sum_t\chi_t$, so $\sum_t\chi_t=rN_P$ by
\cref{lem:dnp-classes}(a). Each turn lies in the half-open intervals of exactly $2(2\Lambda+1)$ hops (\cref{sec:prelim-dnp}),
so $\sum_h\kappa_h=rN_P\cdot2(2\Lambda+1)$. The Lean theorem then gives $\desc(G)=rN_P/q$ and
$q\cdot\cost(G)\le rN_P(4\Lambda+2)+rN_P$.
\end{proof}

\begin{corollary}[The counter serves every level]\label{cor:counter-served}
\notformalized
In the setting of \cref{thm:counter}, let $0<\theta<1$. If $q\ge N_P(1+1/(2\Lambda))/\theta$, for instance if
$q\ge20\cdot2^m/\theta$, then every level is doubly $\theta$-served. If moreover $q\ge32\cdot2^m/\theta$ and
$\theta\le1/3$, then $\cost(G)\le\frac{11}{32}\,\theta\,n_{\mathrm h}/(m+1)\le n_{\mathrm h}/8$.
\end{corollary}

\begin{informalproof}
By \cref{thm:counter}(1), the total error of all hops is at most $\sum_h(\kappa_h-1)<rN_P(4\Lambda+2)$, so every class has
total error below this. It is at most $\theta qn_{\mathrm h}/(2(m+1))=4\theta qr\Lambda$ as soon as
$q\ge N_P(4\Lambda+2)/(4\theta\Lambda)=N_P(1+1/(2\Lambda))/\theta$; and $N_P<2^{m+4}$, $\Lambda\ge2$. By
\cref{thm:counter}(2), $\cost(G)/n_{\mathrm h}\le N_P(4\Lambda+3)/(8(m+1)\Lambda q)$, which for $q\ge32\cdot2^m/\theta$ is
at most $\theta(4\Lambda+3)/(16(m+1)\Lambda)\le\frac{11}{32}\theta/(m+1)$, since $(4\Lambda+3)/\Lambda\le11/2$.
\end{informalproof}

\begin{corollary}[Conjecture A$'$ needs $C\ge1$]\label{cor:C1}
\notformalized
If Conjecture~A$'$ holds with constants $(\theta_0,C,C')$, then $C\ge1$.
\end{corollary}

\begin{informalproof}
Put $\theta=\min(\theta_0,1/3)$ and $q=2^{m+5+\lceil\log_2(1/\theta)\rceil}\ge32\cdot2^m/\theta$. For every~$m$, every
$s_0>m$ and every shift, in particular the one provided by the conjecture, \cref{cor:counter-served} gives a label on a cover of
cost at most $n_{\mathrm h}/8$ that doubly $\theta_0$-serves all $m+1$ levels. Hence
$m+1\le C(m+5+\lceil\log_2(1/\theta)\rceil)+C'$ for all~$m$, and letting $m\to\infty$ gives $C\ge1$.
\end{informalproof}

\begin{corollary}[Cutting the pair with a counter]\label{cor:counter-cut}
\notformalized
For every shift, every $\eps>0$ and every multiple~$T$ of~$P$, the walk DAG of the pair on $[0,T)$ has a cut of at most
$\eps T$ edges and depth at most $\lceil11\cdot2^{m+1}/((m+1)\eps)\rceil-1$. For $m\ge1$ and $\eps\le1$ this is smaller
than the depth $\lceil12/\eps\rceil^{m+1}-1$ of the fourth paper's clock labels.
\end{corollary}

\begin{informalproof}
Let $\mathcal N(t)$ be the number of turns of either head at times in $[0,t]$, $\Psi(t)=\mathcal N(t)\bmod q$, and call $t$
a \emph{wrap time} if $\lfloor\mathcal N(t)/q\rfloor>\lfloor\mathcal N(t-1)/q\rfloor$. A hop $(p,u)$ with no wrap time
in $(p,u]$ is not violated, since then $\Psi(u)-\Psi(p)=\mathcal N(u)-\mathcal N(p)\ge1$ by its owner. Descents occur only
at wrap times. There are at most $\mathcal N(T-1)/q\le(T/P)N_P/q$ wrap times, and each lies in the half-open intervals of
at most $2(2\Lambda+1)$ hops. So $\cost(\Psi)\le(T/P)N_P(4\Lambda+3)/q$, which is at most $\eps T$ for
$q\ge N_P(4\Lambda+3)/(4(m+1)\Lambda\eps)$; this holds for $q=\lceil11\cdot2^{m+1}/((m+1)\eps)\rceil$, because
$N_P<2^{m+4}$ and $(4\Lambda+3)/(4\Lambda)\le11/8$. \Cref{lem:labels} gives depth at most $q-1$. For the comparison,
$11\cdot2^{m+1}/(m+1)\le12^{m+1}/6$ for $m\ge1$ (at $m=1$ it reads $22\le24$, and from $m$ to $m+1$ the left side less
than doubles while the right side grows twelvefold), so with $x=12^{m+1}/\eps\ge144$ we get
$\lceil11\cdot2^{m+1}/((m+1)\eps)\rceil\le\lceil x/6\rceil<x\le(12/\eps)^{m+1}\le\lceil12/\eps\rceil^{m+1}$ for
$\eps\le1$.
\end{informalproof}

\paragraph{Why they matter.} The counter shows what Conjecture~A$'$ must look like: all $m+1$ levels are doubly served
at $\log_2q=m+O(\log(1/\theta))$, within the cost budget, at every shift. So the conjecture, if true, is sharp up to the
constant~$C$. The least depth $D(\eps)$ of a cut of $\eps T$ edges of the walk DAG satisfies
$\log_2D(\eps)\le m+\log_2(1/\eps)+O(1)$ unconditionally, by \cref{cor:counter-cut}, while Conjecture~A$'$ would give $\log_2D(\eps)\ge c(m+1)-c'$ for small~$\eps$ (\cref{cor:wrap}), with $c\le1$ by the counter. So the open question is
exponential against subexponential depth in the nesting depth~$m$, and the base of an exponential is at most~$2$, not the
$c/\eps$ of the clock labels.

\begin{remark}[Numerical: the evidence for Conjecture A$'$]\label{rem:counter-evidence}
The fourth paper quotes, as evidence for Conjecture~A$'$, the numbers $1$, $1$, $2$ and $2$ or~$3$ of doubly
$0.1$-served levels of optimized labels at $m=3$, $s_0=6$ and $q=8,16,32,64$~\cite[Remark~5.26]{FourthPaper2026}. On the
walk DAG of eight periods at the golden shift, the turn counter doubly serves $1$, $1$, $2$ and $3$ levels at the same~$q$,
and at $m=6$, $s_0=10$ it serves $\log_2q-3$ levels for $\log_2q=7,8,9$. The quoted evidence is therefore essentially
what the counter does; it is consistent with Conjecture~A$'$ with $C\approx1$, but it cannot discriminate. (Not all of
these labels meet the cost budget; the comparison is of served counts against~$q$.)
\end{remark}

The counter is a special case of a general criterion.

\begin{proposition}[The counter criterion]\label{prop:counter-criterion}
\notformalized
Let $x_1,\dots,x_k$ be walks of length~$T$ without pauses, with $\mathcal N$ turns in all and at most $H$ hops spanning any
time. Then for every $\eps>0$ the walk DAG has a cut of at most $\eps T$ edges and depth at most
$\max(0,\lceil\mathcal N(H+1)/(\eps T)\rceil-1)$. Hence TCP holds for every family with $(\mathcal N/T)\cdot H$ bounded.
\end{proposition}

\begin{informalproof}
Let $\mathcal N(t)$ count the turns of all walks at times in $[0,t]$, each walk separately, and $\Psi(t)=\mathcal N(t)\bmod q$
with $q=\max(1,\lceil\mathcal N(H+1)/(\eps T)\rceil)$. As in \cref{cor:counter-cut}, a hop with no wrap time in its
half-open interval contains its owner, a turn, and is not violated; there are at most $\lfloor\mathcal N/q\rfloor$ wrap
times, each a descent at most once and inside at most $H$ hops. So the cost is at most
$\lfloor\mathcal N/q\rfloor(H+1)\le\eps T$, and the depth is at most $q-1$.
\end{informalproof}

The criterion is incomparable with the fourth paper's Theorem~B: two zigzags of fixed width whose centres drift
monotonically satisfy it but are not home-based. It fails on the deep-nesting pairs as $m\to\infty$, where
$(\mathcal N/T)\cdot H\approx2^{m+4}(4\Lambda+2)/P\approx2^{m+4}/(m+1)$.

The counter also pins down the capacity trade-off of \cref{cor:capacity} on the deep-nesting family. For two walks write
$\mathrm{Cap}(D,\delta,\eps)$ for the statement: for all pairs of walks of length~$T$ there are a set~$R$ of pairs with
$\card R\le\delta T$ and a cut of $G-R$ with at most $\eps T$ edges and depth at most~$D$.

\begin{proposition}[Capacity on the deep-nesting family]\label{prop:dnp-capacity}
\notformalized
\begin{enumerate}[label=(\alph*)]
\item Let $0<\eps<1/2$, $\delta(\eps)=2/\log_2(1/\eps)$ and $D(\eps)=\lceil11/\eps^2\rceil-1$. For every~$m$, every
$s_0>m$, every shift and every multiple $T$ of~$P$, the pair on $[0,T)$ has a set~$R$ of pairs with
$\card R\le\delta(\eps)T$ and a cut of $G-R$ with at most $\eps T$ edges and depth at most $D(\eps)$.
\item Suppose there are $\eps_0>0$, an integer $N_0\ge1$ and a function~$D_*$ such that for every $m\ge0$ there are $s_0>m$
and an odd shift for which the walk DAG of the pair on $[0,N_0P)$ has no cut of at most $\eps_0N_0P$ edges and depth less
than $D_*(m)$ (a walk-DAG form of Conjecture~A$'$, which follows from Conjecture~A$'$ by \cref{cor:wrap}). If $\mathrm{Cap}(D(\eps),\delta(\eps),\eps)$ holds for all~$\eps$,
then $D(\eps)\ge D_*(m)$ whenever $\eps\le\eps_0/2$ and $4(m+1)N_0\delta(\eps)\le\eps_0$. In particular, if
$D_*(m)\ge2^{c(m+1)-c'}$, then $\delta(\eps)\ge c\eps_0/(4N_0(\log_2D(\eps)+c+c'))$ whenever $\eps\le\eps_0/2$ and
$\delta(\eps)\le\eps_0/(4N_0)$.
\end{enumerate}
\end{proposition}

\begin{informalproof}
(a) If $m+1\ge2/\delta(\eps)$, let $R$ be the set of all pairs visited; then $\card R=2(2\Lambda+1)$ and
$\card R/T\le(4\Lambda+2)/(4(m+1)\Lambda)\le(5/4)/(m+1)\le\frac58\delta(\eps)$, every hop is at a pair of~$R$, and the empty
cut of $G-R$ has depth~$0$. If $m+1<2/\delta(\eps)=\log_2(1/\eps)$, let $R=\emptyset$; \cref{cor:counter-cut} gives depth at
most $\lceil11\cdot2^{m+1}/((m+1)\eps)\rceil-1\le\lceil11/\eps^2\rceil-1$, since $2^{m+1}<1/\eps$.

(b) Take the instance for~$m$ and $T=N_0P$, and let $R$ and a cut~$F$ of $G-R$ be given by the capacity statement. Each pair
is visited at most $2(m+1)N_0$ times in $[0,T)$, so it carries fewer than $2(m+1)N_0$ hops. Let $F_0$ be the set of hop edges
of~$G$ all of whose hops lie at pairs of~$R$; then $\card{F_0}\le2(m+1)N_0\card R\le\eps_0T/2$, and every edge of~$G$
outside~$F_0$ is an edge of $G-R$. So $F\cup F_0$ is a cut of~$G$ with at most $\eps T+\eps_0T/2\le\eps_0T$ edges and depth
at most $D(\eps)$, and the hypothesis gives $D(\eps)\ge D_*(m)$. For the second statement put $y=\eps_0/(4N_0\delta(\eps))\ge1$
and $m+1=\lfloor y\rfloor$; then $\log_2D(\eps)\ge c\lfloor y\rfloor-c'\ge c(y-1)-c'$, that is,
$y\le(\log_2D(\eps)+c+c')/c$.
\end{informalproof}

So on the deep-nesting family, under the walk-DAG form of Conjecture~A$'$ assumed in~(b), and so under Conjecture~A$'$
itself (\cref{cor:wrap}), the scratch must have size about
$T/\log D$ when the rest is cut at depth~$D$, and (a) achieves this: there $\log_2D(\eps)\approx2\log_2(1/\eps)+3.5$ and $\delta(\eps)\approx4/\log_2D(\eps)$.

\subsection{Labels that change only at turns}\label{sec:dnp-turnlabels}

Let $\xi$ be odd. A label~$G$ on the cover, or on the walk DAG of a window, is \emph{turn-supported} if there are integers
$J_h(\tau)$ at the turn times~$\tau$ of each head~$h$ such that $G(t)-G(t-1)\equiv J_h(t)\pmod q$ when $t$ is a turn of
head~$h$, and $G(t)\equiv G(t-1)$ when $t$ is not a turn. This class contains the counter ($J\equiv1$), the fourth paper's
clock labels, and every function of the vector of turn counts. The \emph{potential} of head~$h$ is
$\Phi_h(t)=\sum J_h(\tau)$ over the turns $\tau\le t$ of head~$h$ (on a window, or on~$\Z$ for periodic jumps, normalized
by $\Phi_h(0)=0$), so $G\equiv\Phi_1+\Phi_2+\mathrm{const}\pmod q$. For a hop of head~$h'$, the \emph{other-head
increment} is $F=\Phi_h(u)-\Phi_h(p)$, where $h\ne h'$.

For a V of head~$h$ let $M(V)$ be the sum of the jumps of head~$h$ at the turns in the span of~$V$, with $M(\emptyset)=0$;
let $V^\sigb$ and $V^\sigf$ be the V's attached at the ends of the segments of its back and forward strokes (none if $V$ has
level~$m$), and put
\[
  A_V=J_h(\alpha_V)-1,\qquad E_V=J_h(\nu_V)+M(V^\sigb)-1,\qquad c_\sigb(V)=1+A_V,\qquad c_\sigf(V)=2+A_V+E_V .
\]

\begin{lemma}[Path sums]\label{lem:pathsums}
\notformalized
Let $t$ be a time in a level-$0$ stroke~$S_0$ of head~$h$, and let $V_1\supset V_2\supset\dots\supset V_j$ be the V's whose
spans contain~$t$, with $V_i$ of level~$i$ and $V_1$ the V of~$S_0$. Then
\[
  \Phi_h(t)-\Phi_h(\alpha_{V_1}-1)=\sum_{i=1}^jc_{\delta_i}(V_i),
\]
where $\delta_i=\sigb$ if $t$ lies in the back stroke of~$V_i$ and $\delta_i=\sigf$ otherwise; for $t$ in the segment of~$S_0$,
$j=0$. Moreover $M(V)=c_\sigf(V)+M(V^\sigf)$.
\end{lemma}

\begin{informalproof}
By induction on the depth inside one V. The apex $\alpha_V$ is the first time of~$V$. For~$t$ in the segment of the back stroke,
the only turn in $[\alpha_V,t]$ is $\alpha_V$, so the increment over $(\alpha_V-1,t]$ is $J_h(\alpha_V)=c_\sigb(V)$; for $t$
in $V^\sigb$ it is $c_\sigb(V)$ plus the increment from $\alpha_{V^\sigb}-1$ to~$t$. For $t$ in the segment of the forward
stroke the increment is $J_h(\alpha_V)+M(V^\sigb)+J_h(\nu_V)=c_\sigf(V)$, by the definition of $E_V$, and for $t$ in
$V^\sigf$ it is $c_\sigf(V)$ plus the increment from $\alpha_{V^\sigf}-1$ to~$t$. At level~$m$ there are no inner V's and
the same computation applies with $M(\emptyset)=0$. The mass identity is the case where $t$ is the last time of~$V$. The
segment of~$S_0$ contains no turn after its first time.
\end{informalproof}

\begin{lemma}[Band decomposition]\label{lem:band}
\notformalized
Let $G$ be turn-supported, and let $B$ be a V of head~$h'$ of level $k\ge1$, with chain $C_{k+1},\dots,C_m$.
\begin{enumerate}[label=(\alph*)]
\item For $1\le d<R_k$ the apex hop $(\alpha-d,\alpha+d)$ has error $(A_B+F)\bmod q$.
\item For $k\le l\le m$ and $R_{l+1}<x<R_l$, the valley hop at $v+\sigma x$ has error
\[
  \Bigl(E_B-\sum_{i=k+1}^{l}\bigl(A_{C_i}+E_{C_i}+2\bigr)+F\Bigr)\bmod q .
\]
\end{enumerate}
Here $F$ is the other-head increment of the hop.
\end{lemma}

\begin{informalproof}
The error is $(G(u)-G(p)-1)\bmod q$, and $G(u)-G(p)\equiv\Phi_{h'}(u)-\Phi_{h'}(p)+F$. (a) By \cref{lem:dnp-classes}(b) the
only turn of head~$h'$ in $(\alpha-d,\alpha+d]$ is~$\alpha$, with jump $1+A_B$. (b) By \cref{lem:dnp-classes}(b) the
increment of $\Phi_{h'}$ is $M(C_{l+1})+J_{h'}(\nu)$, and $J_{h'}(\nu)=E_B+1-M(C_{k+1})$. Since $C_{i+1}=C_i^\sigf$, the mass
identity of \cref{lem:pathsums} gives $M(C_{k+1})-M(C_{l+1})=\sum_{i=k+1}^l(A_{C_i}+E_{C_i}+2)$.
\end{informalproof}

When $F=0$ on the hops concerned (the \emph{single-head regime}), $A_B$ is the apex error of~$B$ and $c_\sigb-1$,
$c_\sigf-c_\sigb-1$ are errors of $B$'s own hops. So a V whose own errors lie in $[0,E]$ contributes at least $1$ and~$2$ to
the path sums of \cref{lem:pathsums}: the potential of a head that serves many of its levels rises along every path, and it
can come back down only through V's with an apex or band error near~$q$. Serving many levels cheaply needs cancellation,
and the next labels cancel in blocks.

\begin{definition}[Block labels]\label{def:block}
Let $a\ge1$ and $K\ge0$ be integers. A level $l\in[1,m]$ is a \emph{canceller} if $m-l\equiv a\pmod{a+1}$ and
\emph{served} otherwise. Define integers $A_l,E_l$ from level~$m$ upwards: for a served level~$l$, $A_l=K$ and
$E_l=K+\sum_{i=l+1}^{c-1}(A_i+E_i+2)$, where $c$ is the first canceller below~$l$, or $m+1$ if there is none; for a
canceller~$c$, $A_c=0$ and $E_c=-Q_c-2$ with $Q_c=\sum_{i=c+1}^{c'-1}(A_i+E_i+2)$, where $c'$ is the next canceller
below~$c$, or $m+1$. The \emph{block label} $B(a,K)$ on head~$h$ has the jumps $J_h(\alpha_V)=1+A_{\lev V}$,
$J_h(\nu_V)=E_{\lev V}+1-M(V^\sigb)$, computed from level~$m$ upwards, and $J_h=1$ at the tips; so $A_V=A_{\lev V}$ and
$E_V=E_{\lev V}$.
\end{definition}

\begin{lemma}[Block masses]\label{lem:block}
\notformalized
In $B(a,K)$, a served level with $j$ served levels below it in its block has $A+E+2=2^j(2K+2)$. Every canceller has $a$
served levels directly below it, and $Q_c=Q:=(2^a-1)(2K+2)$. Hence $c_\sigb(V)=1+K$ and $c_\sigf(V)=2^j(2K+2)$ at a served
V, and $c_\sigb(V)=1$ and $c_\sigf(V)=-Q$ at a canceller V: the path sums depend only on the level. In the single-head
regime, and if $K+Q<q$, every apex hop with $d<R_k$ of a served V has error~$K$ and every valley hop of
\cref{lem:band}(b) of a served V has error in $[K,K+Q]$.
\end{lemma}

\begin{informalproof}
Write $T(j)$ for $A+E+2$ at a served level with $j$ served levels below it in its block. Then $T(0)=2K+2$ and
$T(j)=2K+2+\sum_{j'<j}T(j')$, so $T(j)=2^j(2K+2)$. A canceller~$c$ satisfies $m-c\equiv a\pmod{a+1}$ and $m-c\ge0$, so
$m-c\ge a$ and the levels $c+1,\dots,c+a$ are served and below~$c$ in one block; hence
$Q_c=\sum_{j<a}T(j)=Q$. The values of $c_\sigb$ and $c_\sigf$ follow. For the errors, write
$S(i_1,i_2)=\sum_{i=i_1}^{i_2}(A_i+E_i+2)$. A served level~$k$ has $E_k=K+S(k+1,c-1)$ with $c$ the first canceller below~$k$.
For $k\le l<c$ the band error of \cref{lem:band}(b) is $K+S(l+1,c-1)\in[K,K+Q)$. For $c\le l<c'$, with $c'$ the next
canceller, it is $K-S(c,l)=K+Q-S(c+1,l)=K+S(l+1,c'-1)\in[K,K+Q]$, because $A_c+E_c+2=-Q$ and $S(c+1,c'-1)=Q$; and the
pattern repeats below~$c'$. All values lie in $[0,q)$, so no reduction occurs.
\end{informalproof}

So, in the single-head regime, one head carrying $B(a,K)$ serves the $a/(a+1)$ fraction of its levels that are served,
with errors in $[K,K+Q]$, and its violated own hops are essentially the valley hops of the canceller classes. Heuristically,
its cost per own hop is then about $1/(2(a+1))$, which for $a=4$ is $1/10$, inside the budget $n_{\mathrm h}/8$ for
every~$m$ at fixed~$q$ (see also \cref{rem:single-head}). So no argument that uses only one head's own structure can be
expected to prove Conjecture~A$'$. At the aligned shift the second head inherits the first head's errors.

\begin{remark}[Numerical: the single-head cost of block labels]\label{rem:single-head}
Exact evaluations on covers gave costs per own hop of $0.05$ to $0.09$ for block labels with $a=4$, $q=256$ and
$4\le m\le10$.
\end{remark}

\begin{proposition}[The aligned shift]\label{prop:aligned}
\notformalized
Let $\xi=1$, put the same jumps on both heads, $J_2(t)=J_1(t+1)$, and for a hop $(a',b')$ of head~$1$ write
$I(a',b')=\Phi_1(b')-\Phi_1(a')$. The hops of head~$2$ are the pairs $(a'-1,b'-1)$ for the hops $(a',b')$ of head~$1$, and
\begin{align*}
  G(b'-1)-G(a'-1)&\equiv2I(a',b')-J_1(b')[b'\text{ a turn}]+J_1(a')[a'\text{ a turn}],\\
  G(b')-G(a')&\equiv2I(a',b')+J_1(b'+1)[b'+1\text{ a turn}]-J_1(a'+1)[a'+1\text{ a turn}].
\end{align*}
So every hop whose shifted endpoints are not turns has error $(1+2e)\bmod q$, where $e=(I-1)\bmod q$ is the single-head
error of the corresponding hop of head~$1$; each turn affects at most four hops.
\end{proposition}

\begin{informalproof}
Head~$2$ is at a cell at time~$t$ exactly when head~$1$ is there at time $t+1$, so its hops are the shifted hops of head~$1$
and its turns are the turns of head~$1$ minus one. Hence $\Phi_2(t)=\Phi_1(t+1)$ and $G(t)\equiv\Phi_1(t)+\Phi_1(t+1)$ up
to a constant, and both identities follow by expanding $\Phi_1(t\pm1)-\Phi_1(t)$. Each time is the earlier and the later
endpoint of exactly one hop of each head.
\end{informalproof}

\begin{remark}[Numerical: Conjecture A$'$ needs a misaligned shift]\label{rem:aligned}
With $B(4,0)$ on both heads at the aligned shift and $q=256$, exact evaluations on covers gave costs $0.080$, $0.0996$ and
$0.0935$ times~$n_{\mathrm h}$ at $m=6$, $9$ and~$10$, inside the budget $n_{\mathrm h}/8$, with $6$ of $7$, $8$ of $10$ and $9$
of $11$ levels doubly $0.1$-served; \cref{prop:aligned} explains this, since at the aligned shift both heads see their
single-head errors, doubled, on all but a vanishing fraction of hops. This indicates that Conjecture~A$'$ cannot hold with the
aligned shift, so that its choice of a misaligned shift is essential. At misaligned shifts
(the golden shift, and a shift whose dyadic phases alternate near $1/3$ and $2/3$ at every level) the same labels
collapse to the finest few levels at costs $0.22$ to $0.24$ times~$n_{\mathrm h}$.
\end{remark}

\subsection{Blindness of fine classes}\label{sec:dnp-blind}

\begin{lemma}[Turn-free runs]\label{lem:runs}
\notformalized
Let $m\ge1$. Over one period, the gaps between consecutive turns of~$U$ have the following lengths: for each level $1\le i\le m$,
$2^{i+2}$ gaps of length~$R_i$, except that two of the gaps of level~$m$ have length $R_m+R_0$ instead; and two gaps of
length~$R_0$.
\end{lemma}

\begin{informalproof}
We find the gap that follows each turn. After an apex of level $i<m$ comes the back segment, $R_i$ steps, and then the apex
of the V attached at its end; after an apex of level~$m$, a plain back stroke of $R_m$ steps and then the valley. After a
valley of level $i<m$ comes the forward segment, $R_i$ steps, and then an apex. After a tip comes the segment of a level-$0$
stroke, $R_0$ steps, and then an apex. After a valley~$\nu$ of level~$m$, the forward stroke takes $R_m$ steps and the V
ends at $\nu+R_m$, together with every stroke of which it is the last part. If some enclosing stroke is a back stroke, the
deepest such ends at $\nu+R_m$ with the valley of its V, a turn. If not, the level-$0$ stroke ends at $\nu+R_m$; this happens
for one level-$m$ V per level-$0$ stroke. If the next level-$0$ stroke is the second or fourth of its period, $\nu+R_m$ is a
tip; otherwise the directions continue (from the second to the third stroke, and from the fourth to the first), the next
segment adds $R_0$ steps, and the next turn is at $\nu+R_m+R_0$; this happens twice per period. Each level~$i$ has
$2^{i+2}$ turns per period, and the lengths add up to $4m\Lambda+4\Lambda=P$.
\end{informalproof}

\begin{proposition}[Blindness]\label{prop:blind}
\notformalized
Let $m\ge1$, let $G$ be turn-supported, let $(p,u)$ be a hop of head~$h'$ of length $\ell=u-p$, and let the shift~$\xi$ be uniformly
random in~$\Z_P$. If $(p,u]$ contains no turn of the other head, then $F=0$. The probability of this event is
\[
  \frac1P\sum_{\text{gaps }L}(L-\ell)^+\;\ge\;\frac1{m+1}\sum_{i=0}^m\Bigl(1-\frac\ell{R_i}\Bigr)^+,
\]
the sum over the gaps of \cref{lem:runs}. For $\xi$ uniform over the odd residues the bound decreases by at most
$2/((m+1)R_m)$.
\end{proposition}

\begin{informalproof}
$F$ is the sum of the other head's jumps over the window, an empty sum. In the other head's own time, where it is~$U$, the
window $(p,u]$ is a window $(a',a'+\ell]$ with $a'$ uniform in~$\Z_P$: a hop of head~$2$, given as a hop $(a,b)$ of~$U$ in the
own time of head~$2$, occupies the window $(a-\xi,b-\xi]$ of head~$1$; and for a hop $(a,b)$ of head~$1$, the turns of head~$2$
are those of~$U$ shifted by~$-\xi$, so the window is $(a+\xi,b+\xi]$ in the time of~$U$. This window contains no turn exactly when
$\tau_j\le a'$ and $a'+\ell<\tau_{j+1}$ for consecutive turns $\tau_j<\tau_{j+1}$, which gives $(L_j-\ell)^+$ values
of~$a'$ for the gap of length $L_j=\tau_{j+1}-\tau_j$. Insert \cref{lem:runs}, using
$(x+y-\ell)^+\ge(x-\ell)^++(y-\ell)^+$ for $x,y,\ell\ge0$ for the two long gaps, and $2^{i+2}/P=1/((m+1)R_i)$,
$4/P=1/((m+1)R_0)$. For odd shifts, each interval of admissible values of~$a'$ contains at least $((L_j-\ell)^+-1)/2$
values of each parity, there are fewer than $2^{m+3}$ gaps, and $2^{m+3}/P=2/((m+1)R_m)$.
\end{informalproof}

\begin{corollary}[Fine classes are blind]\label{cor:blind}
\notformalized
Let $1\le k\le m$, $g=m-k$, and let $\xi$ be uniform in~$\Z_P$. On average over~$\xi$, at least a fraction
\[
  \frac{k-2-\lceil\log_2(g+2)\rceil}{m+1}
\]
of the hops of each class $(h',k)$ contain no turn of the other head. Consequently, for every turn-supported label, the
expected fraction of the hops of the class with $F\ne0$ is at most
\[
  \frac{g+3+\lceil\log_2(g+2)\rceil}{m+1},
\]
which is $O(\log m/m)$ for the classes with $g=O(\log m)$.
\end{corollary}

\begin{informalproof}
By \cref{lem:dnp-classes}(b), the hops of the class are, for each V of level~$k$, its apex hops $(\alpha-d,\alpha+d)$,
$1\le d\le R_k$, and its valley hops at the cells $v+\sigma x$, $1\le x\le R_k$; the boundary hops among them obey the
same length bounds as the others. An apex hop has length $2d\le2R_k=R_{k-1}$, so by \cref{prop:blind} it is blind with
probability at least $\frac1{m+1}\sum_{i=0}^{k-1}(1-2^{i-k+1})\ge\frac{k-2}{m+1}$. For $k\le l\le m$ and
$R_{l+1}<x\le R_l$, the valley hop at $v+\sigma x$ is $(\nu-(m-l)R_l-x,\nu+x)$: for $x<R_l$ by \cref{lem:dnp-classes}(b),
and for the boundary hop $x=R_l$ because the cell $v+\sigma R_l$ is visited at the valley time $\nu-n_lR_l$ of~$C_l$ (at the
apex $\alpha=\nu-n_kR_k$ if $l=k$), left at the next step, and not visited again before $\nu+R_l$. So its length is at most
$(m-l+2)R_l\le R_{l-a}$, where $a=\lceil\log_2(m-l+2)\rceil$; if $l\ge a$ it is blind with probability at least
$\frac1{m+1}\sum_{i=0}^{l-a}(1-2^{i-l+a})\ge\frac{l-a-1}{m+1}\ge\frac{k-1-\lceil\log_2(g+2)\rceil}{m+1}$, using $l\ge k$
and $m-l\le g$; if $l<a$ the claimed bound is negative. The second statement follows, because $F\ne0$ needs a turn of the
other head in the window.
\end{informalproof}

\paragraph{Why it matters.} By Markov's inequality, for each class with $g=O(\log m)$ the set of shifts at which a fixed
fraction $c>0$ of the class has $F\ne0$ has density $O(\log m/(cm))$: for all large~$m$, at all but a vanishing density of
shifts, these classes see the other head on less than a fraction~$c$ of their hops, for every label that changes only at
turns. On their blind hops the error is the single-head error, so a head can serve them as if alone, for instance by a block
label. Heuristically, then, at typical shifts no argument that charges errors to the other head's increments can bound the
number of doubly served fine levels by $C\log_2q+C'+O(\log m)$, and if Conjecture~A$'$ holds at a typical shift, it holds
because serving the coarse levels is expensive. This is not a proof: the corollary bounds an average over shifts, and the
conjecture may choose an atypical shift. The measurements of \cref{rem:blind-numerics} agree with this picture.

\begin{remark}[Numerical: block pairs at a misaligned shift]\label{rem:blind-numerics}
With $B(4,0)$ on both heads at a misaligned shift and $q=256$, the number of doubly $0.1$-served levels was
$4,4,4,5,5,8,8$ at $m=10,15,20,25,30,40,50$, all in the blind fine region, at costs of $0.23$ to $0.29$ times
$n_{\mathrm h}$, above the budget; the value at $m=10$ is exact on a cover, the others are sampled.
\end{remark}

\subsection{Visibility of coarse classes}\label{sec:dnp-visible}

The coarse classes see the other head's cancellations. In the other head's own time an apex hop of half-width~$d$ is a
window $(c_0-d,c_0+d]$ around a centre~$c_0$, and its other-head increment is $F(d)=\Phi_h(c_0+d)-\Phi_h(c_0-d)$.

\begin{lemma}[Straddle involution]\label{lem:straddle}
\notformalized
Let head~$h$ carry $B(a,K)$, let $c_0$ be an integer, $k\ge1$, and let $V$ be a V of head~$h$ of a canceller level~$\kappa$,
with apex~$\alpha_V$, valley~$\nu_V$ and stroke length $\Pi_\kappa=\nu_V-\alpha_V=n_\kappa R_\kappa$. Let $\mathcal D$ be the
set of $d\in[1,R_k-1]$ with $c_0-d\in[\alpha_V,\nu_V)$, $c_0+d\in[\nu_V,\nu_V+\Pi_\kappa)$ and $\Pi_\kappa-d\in[1,R_k-1]$.
Then $d\mapsto\Pi_\kappa-d$ is an involution of~$\mathcal D$, and $F(d)+F(\Pi_\kappa-d)=-2Q-2$ for every $d\in\mathcal D$.
Hence at least half of the elements of~$\mathcal D$ have $F(d)\le-Q-1$.
\end{lemma}

\begin{informalproof}
For $d\in\mathcal D$ put $o_p=c_0-d-\alpha_V$ and $o_u=c_0+d-\nu_V$, both in $[0,\Pi_\kappa)$. For $d'=\Pi_\kappa-d$ we get
$c_0-d'-\alpha_V=o_u$ and $c_0+d'-\nu_V=o_p$, so $d'\in\mathcal D$ and the map swaps $o_p$ and~$o_u$. By \cref{lem:pathsums},
for a time~$x$ in the back stroke of~$V$, $\Phi_h(x)=C_V+c_\sigb(V)+Y(x-\alpha_V)$, and for $x$ in its forward stroke,
$\Phi_h(x)=C_V+c_\sigf(V)+Y(x-\nu_V)$. Here $C_V$ collects $\Phi_h(\alpha_{V_1}-1)$ and the contributions of the V's
containing~$V$, the same for all times in~$V$, and $Y(o)$ is the sum of $c_\delta(W)$ over the V's~$W$ inside a stroke of
level~$\kappa$ that contain its time at offset~$o$; by \cref{lem:block} these contributions depend only on levels, so $Y$ is
the same function for both strokes of~$V$. Hence $F(d)=c_\sigf(V)-c_\sigb(V)+Y(o_u)-Y(o_p)$ and
$F(\Pi_\kappa-d)=c_\sigf(V)-c_\sigb(V)+Y(o_p)-Y(o_u)$, and their sum is $2(c_\sigf(V)-c_\sigb(V))=2(-Q-1)$ by
\cref{lem:block}.
\end{informalproof}

\begin{theorem}[Theorem V: visibility of coarse classes on average over shifts]\label{thm:V}
\notformalized
Let head~$h$ carry $B(a,K)$ and let $\xi$ be uniform in~$\Z_P$. For every class $(h',k)$ of the other head with
$1\le k\le m$ and $g=m-k\ge\lceil\log_2(g+1)\rceil+a$, the expected fraction of its apex hops with $d<R_k$ that have
$F\le-Q-1$ is at least
\[
  \frac{g-\lceil\log_2(g+1)\rceil-2a}{2^{a+4}(m+1)}.
\]
\end{theorem}

\begin{informalproof}
Let $j_1=\min\{j\ge0:g+1-j\le2^j\}$; then $1\le j_1\le\lceil\log_2(g+1)\rceil$, since $g\ge1$. Let $\kappa=k+j$ be the first
canceller level with $j\ge j_1$; cancellers recur every $a+1$ levels, so $j\le j_1+a$, and $\kappa\le m$ by the hypothesis
on~$g$. Since $j\ge j_1$, $\Pi_\kappa=n_\kappa R_\kappa=(g+1-j)R_k2^{-j}\le R_k$, and $\Pi_\kappa$ is even.

Fix a V of level~$k$ of head~$h'$ and its apex hops. As $\xi$ is uniform, the centre~$c_0$ of their windows in the own time of
head~$h$ is uniform modulo~$P$. A window of half-width~$d$ \emph{straddles} a V of level~$\kappa$ of head~$h$ if its left end
lies in the back stroke and its right end in the forward stroke; then $d<\Pi_\kappa$, and $c_0$ lies in the span of that V,
so for fixed~$c_0$ at most one V of level~$\kappa$ is straddled, and for $d\le\Pi_\kappa-1$ the partner
$\Pi_\kappa-d$ lies in $[1,R_k-1]$. So, for fixed~$c_0$, the straddling half-widths form the set~$\mathcal D$ of
\cref{lem:straddle} for the straddled V. For a fixed~$d$, the centres that straddle a given V of level~$\kappa$ form an interval
of length $2\min(d,\Pi_\kappa-d)$, and there are $2^{\kappa+1}$ such V's per period; so
\[
  \mathbb E\card{\mathcal D}=\sum_{d=1}^{\Pi_\kappa-1}\frac{2^{\kappa+2}\min(d,\Pi_\kappa-d)}P
  =\frac{2^{\kappa+2}}P\cdot\frac{\Pi_\kappa^2}4=\frac{2^\kappa\Pi_\kappa^2}P .
\]
By \cref{lem:straddle}, at least $\card{\mathcal D}/2$ of the $R_k-1$ apex hops with $d<R_k$ have $F\le-Q-1$, so the expected
fraction is at least
\[
  \frac{2^{\kappa-1}\Pi_\kappa^2}{PR_k}=\frac{n_\kappa^2R_\kappa}{8(m+1)R_k}=\frac{(g+1-j)^22^{-j}}{8(m+1)},
\]
using $2^\kappa=\Lambda/R_\kappa$ and $P=4(m+1)\Lambda$. By the minimality of~$j_1$, $g+2-j_1>2^{j_1-1}$, so with
$G_1=g+2-j_1$ we have $2^j\le2^{j_1+a}<2^{a+1}G_1$ and $(g+1-j)^2\ge(G_1-1-a)^2\ge G_1^2-2(a+1)G_1$, where
$G_1-1-a=g+1-j_1-a\ge1$. Hence $(g+1-j)^22^{-j}>(G_1-2(a+1))/2^{a+1}=(g-j_1-2a)/2^{a+1}$, and dividing by $8(m+1)$ and using
$j_1\le\lceil\log_2(g+1)\rceil$ gives the bound.
\end{informalproof}

\paragraph{Why it matters.} Blindness bounds the visible fraction of a class from above by
$(g+3+\lceil\log_2(g+2)\rceil)/(m+1)$, and Theorem~V bounds it from below, for block labels, by about $g/(2^{a+4}(m+1))$:
both are of order $(m-k)/(m+1)$ up to logarithmic factors and constants depending on~$a$. So the coarse classes, a positive
fraction of all levels, see the other head's cancellations at full size on a positive fraction of their apex hops, on
average over shifts (see also \cref{rem:visible-numerics}). The next proposition says what this would give for block pairs,
but it needs more than Theorem~V provides. First, Conjecture~A$'$ fixes one shift for all classes at once. Second, the
proposition needs the cancellation in a window modulo~$q$, $-q/2<F\le-Q_h/2$, while Theorem~V gives only the integer event
$F\le-Q-1$, with no lower bound on~$F$: by the proof of \cref{lem:straddle}, $F=-Q-1+Y(o_u)-Y(o_p)$, and $Y$ sums
contributions of up to $m-\kappa$ nested levels, so $\abs F$ may grow linearly in~$m$, and the event is guaranteed to imply the
window only for $q$ of order $(m-\kappa)Q$ or larger.

\begin{remark}[Numerical: visible fractions]\label{rem:visible-numerics}
Measured fractions of apex hops with $F\le-Q-1$ were ten to a hundred times larger than the bound of Theorem~V.
\end{remark}

\begin{proposition}[Visibility at a fixed shift bounds the served levels of block pairs]\label{prop:Vfixed}
\notformalized
Let the heads carry $B(a_1,K_1)$ and $B(a_2,K_2)$ with $a_1,a_2\ge1$, at a fixed odd shift, on the $r$-fold cover with $r$
a multiple of~$q$, so that the two potentials define a label there modulo~$q$, and let $Q_1,Q_2$ be their block masses. Let $0<\beta<1$ and $\theta_0>0$, and suppose that for each head~$h'$ and each class
$(h',k)$ with $m-k\ge\beta(m+1)$, a fraction at least $c(\beta)>8\theta_0$ of its apex hops with $d<R_k$ have
$-q/2<F\le-Q_h/2$, where $h$ is the other head. Then at most $\beta(m+1)+2$ levels are doubly $\theta_0$-served.
\end{proposition}

\begin{informalproof}
Since $Q_h=(2^{a_h}-1)(2K_h+2)\ge2K_h+2$, the inequalities $K_1\ge Q_2/2$ and $K_2\ge Q_1/2$ cannot both hold: they would
give $Q_1\ge2K_1+2\ge Q_2+2$ and $Q_2\ge Q_1+2$. So some head~$h'$ has $K_{h'}<Q_h/2$ for the other head~$h$. Let $k\ge1$
with $m-k\ge\beta(m+1)$. By \cref{lem:band}(a) an apex hop with $d<R_k$ of a V of level~$k$ of head~$h'$ has error
$(A+F)\bmod q$ with $A\in\{0,K_{h'}\}$; on the hops of the hypothesis, $-q/2<A+F<0$, so the error lies in $[q/2,q)$. The
apex hops with $d<R_k$ are a fraction $(R_k-1)/(2R_k)\ge1/4$ of the class, so the class has total error at least
$c(\beta)\cdot\frac14\cdot\frac q2$ times its size, more than $\theta_0q$ times its size, and level~$k$ is not
$\theta_0$-served for head~$h'$. The levels $k\ge1$ with $m-k<\beta(m+1)$, together with level~$0$, number at most
$\beta(m+1)+2$.
\end{informalproof}

\subsection{A safe form of Conjecture A\texorpdfstring{$'$}{'}}\label{sec:dnp-App}

Blindness suggests, heuristically, that cheap labels may serve a number of fine levels that grows with~$m$; in the
measurements of \cref{rem:blind-numerics} the block pairs serve between $4$ and $8$ levels for $10\le m\le50$, above the
budget. The following form of the conjecture allows up to a fixed fraction of the levels to be served and keeps its
consequence.

\begin{conjecture}[Conjecture A$''$]\label{conj:App}
There are $\theta_0>0$, $\eta<1$ and constants $C,C'$ such that for all $s_0$ and~$m$ some odd shift~$\xi$ has the
following property: for all $r$ and~$q$, every label of cost at most $n_{\mathrm h}/8$ on the $r$-fold cover of the pair
$x_1=U(s_0,m,1)$, $x_2=x_1(\cdot+\xi)$, has at most $C\log_2q+C'+\eta(m+1)$ doubly $\theta_0$-served levels.
\end{conjecture}

If Conjecture~A$''$ holds and a label of cost at most $\eps n_{\mathrm h}$ with $\eps\le1/8$ and $\eps<(1-\eta)\theta_0/2$
exists on such a cover, then, as in \cref{sec:prelim-dnp}, it doubly serves more than $(m+1)(1-2\eps/\theta_0)$ levels, so
\[
  \log_2q\;>\;\frac{(1-\eta-2\eps/\theta_0)(m+1)-C'}{C}.
\]
By \cref{cor:wrap} this makes the pairs depth-robust as $m\to\infty$, and $\TCP(2)$ would fail; only the admissible
range of~$\eps$ shrinks, from $\theta_0/2$ to $(1-\eta)\theta_0/2$. \Cref{prop:Vfixed} proves Conjecture~A$''$ for pairs of
block labels, with $\eta=\beta+O(1/m)$ and no logarithmic term, from visibility in a window modulo~$q$ at a fixed shift.
Theorem~V does not supply this hypothesis, even on average over shifts: it gives the integer event $F\le-Q-1$ with no lower
bound on~$F$, hence the window only for $q$ of order $(m-\kappa)Q$, and only on average over shifts. Two steps remain open: a
bound on $\abs F$ on the straddling hops, or a version of Theorem~V that counts the window directly; and, at one fixed shift
for all coarse classes, equidistribution of a class's apex centres relative to the other head's canceller valleys of level
about $k+\log_2(m-k)$, a mixed-radix carry problem that we could not solve.

\subsection{The walk-DAG forms of the conjectures}\label{sec:dnp-wrap}

Conjectures~A$'$ and~A$''$ concern labels on the cyclic covers, where the chains of hops close up and Lemma~Zq is exact,
while $\TCP(2)$ concerns cuts of the walk DAG on $[0,T)$, where the chains are open. A label of the walk DAG costs at most
a constant more on the cover, so the conjectures already speak about the walk DAG; this is the elementary transfer in the
spirit of the fourth paper's remark on the cover and the walk DAG~\cite[Remark~5.23]{FourthPaper2026}, written for
served levels.

\begin{lemma}[From the walk DAG to the cover]\label{lem:wrap}
\notformalized
Let $s_0\ge m+2$, let $\xi$ be odd, let $N\ge1$ and $T=NP$, and let $\Psi\colon[0,T)\to\{0,\dots,q-1\}$ be a label of the
walk DAG of the pair $x_1=U(s_0,m,1)$, $x_2=x_1(\cdot+\xi)$ on $[0,T)$. Then $\Psi$, read as a label $G$ of the $N$-fold
cover, satisfies $\cost(G)\le\cost(\Psi)+2(2\Lambda+1)+1$.
\end{lemma}

\begin{informalproof}
Let $(p,u)$ be a hop of the cover, of head~$h$, lifted to~$\Z$ with $u\in[0,T)$, so that $p$ is the previous visit of
head~$h$ to the cell $x_h(u)$. If $p\ge0$, then $p$ is also the previous visit within $[0,T)$; since $U$ has no pauses,
$u\ge p+2$, so $(p,u)$ is a long hop of the walk DAG, and it is violated by~$G$ exactly when it is violated by~$\Psi$. If
$p<0$, then $u$ is the first visit of head~$h$ within $[0,T)$ to its cell; each head stays in the $2\Lambda+1$ cells of
$[-\Lambda,\Lambda]$, so there are at most $2(2\Lambda+1)$ such hops of the two heads. No pair is a hop of both heads, so
every long hop of the walk DAG is counted once on the cover: all turn times and turn positions are multiples of~$R_m$, so
a hop of~$x_1$ has $p+u\equiv0$ and a hop of~$x_2$ has $p+u\equiv-2\xi\pmod{R_m}$~\cite[Remark~5.23]{FourthPaper2026}, and
$2\xi\not\equiv0\pmod{R_m}$ because $\xi$ is odd and $R_m=2^{s_0-m}\ge4$. Finally, the descents of~$G$ at the times
$t\in[1,T)$ are those of~$\Psi$, and the cyclic comparison at $t=0$ adds at most one.
\end{informalproof}

\begin{corollary}[The conjectures on the walk DAG]\label{cor:wrap}
\notformalized
Suppose that Conjecture~A$''$ holds with constants $\theta_0$, $\eta$, $C$ and~$C'$. Let $0<\eps<\min(1/4,(1-\eta)\theta_0)$,
$\delta=((1-\eta)-\eps/\theta_0)/2$ and $N_0=\lceil1/\min(1/8-\eps/2,\delta\theta_0/2)\rceil$. For every~$m$, with $s_0=m+2$
and the shift~$\xi$ of the conjecture, and every $N\ge N_0$, every cut of at most $\eps NP$ edges of the walk DAG of the
pair on $[0,NP)$ has depth at least $2^{(\delta(m+1)-C')/C}-1$. So the deep-nesting pairs are depth-robust as $m\to\infty$
and $\TCP(2)$ fails. The same holds under Conjecture~A$'$, with $\eta=0$.
\end{corollary}

\begin{informalproof}
Let $F$ be a cut with at most $\eps NP$ edges and depth~$D$. By the converse of \cref{lem:labels}~\cite[Lemma~5.2]{FourthPaper2026}
there is a label~$\Psi$ with $q=D+1$ values and $\cost(\Psi)\le\eps NP$. By \cref{lem:wrap} the label~$G=\Psi$ of the
$N$-fold cover has $\cost(G)\le\eps NP+4\Lambda+3\le\eps_cn_{\mathrm h}$ with $\eps_c=\eps/2+1/(N(m+1))$, because
$n_{\mathrm h}=2NP$, $P=4(m+1)\Lambda$ and $4\Lambda+3\le8\Lambda$. Since $N\ge N_0$, $\eps_c\le1/8$ and
$2\eps_c/\theta_0\le\eps/\theta_0+\delta<1-\eta$, so $\eps_c<\theta_0/2$. By the per-class split
(\cref{sec:prelim-dnp}), $G$ doubly $\theta_0$-serves more than $(m+1)(1-2\eps_c/\theta_0)\ge(m+1)(1-\eps/\theta_0-\delta)$
levels, and Conjecture~A$''$ bounds their number by $C\log_2q+C'+\eta(m+1)$. As $1-\eta-\eps/\theta_0=2\delta$, this gives
$C\log_2q>\delta(m+1)-C'$, which is the claim.
\end{informalproof}

The proof uses Conjecture~A$''$ only through a bound on the number of doubly $\theta_0$-served levels by a function
of~$q$ plus $\eta(m+1)$. So a weaker statement already refutes $\TCP(2)$: that for every~$q$ there are $F(q)$
and~$m_0(q)$ such that for every $m\ge m_0(q)$, with $s_0=m+2$, some odd shift makes every label with $q$ values and cost at
most $n_{\mathrm h}/8$ on every cover doubly $\theta_0$-serve at most $F(q)+\eta(m+1)$ levels. Indeed, if $\TCP(2)$ held with
depth~$D$ at some $\eps<\min(1/4,(1-\eta)\theta_0)$, take $q=D+1$ and $m\ge\max(m_0(q),F(q)/\delta)$; the proof above gives
$F(q)+\eta(m+1)>(m+1)(\eta+\delta)$, a contradiction. The function~$F$ need not be monotone, and the shift may depend on~$q$.


\section{Consequences for transposers}\label{sec:transposers}

Theorem~D of the fourth paper (\cref{thm:D4}) excludes oblivious linear-time transposers under two hypotheses, depth
reduction and the crux. For machines whose work heads enter every cell boundedly often, Theorem~F supplies the first: the
crux alone excludes them (Corollary~F). \emph{Directed Recovery Structures in Tape Computations} shows that the cone
relations of every linear-time schedule are paired-distance separated after paid trimming, and asks whether all codes on
such relations have rate at most $\frac12+o(1)$~\cite[Question~9.1]{DirectedRecovery2026}; an affirmative answer excludes
bounded-entry transposers, oblivious (Corollary~F$'$) or of low adaptivity (Corollary~F$''$). Theorem~M gives the analogous
statement for machines with one unrestricted work head (\cref{cor:M-machine}).

\subsection{Corollary F}\label{sec:trans-F}

\begin{corollary}[Corollary F: the crux excludes bounded-entry oblivious transposers]\label{cor:F}
\notformalized
Let $k\ge1$, $\beta\ge2$ and $V\ge1$ be integers, $C\ge1$ and $0<\eps\le1$, and put
\[
  D_1=\DF\bigl(k,2V-1,\eps/(4\beta C)\bigr),\qquad P=\lceil4\beta C/\eps\rceil,\qquad D=D_1P .
\]
Suppose that $\Crux(\mathcal S_{k,C,\beta,D},\rho)$ holds for some $\rho>\eps$. Then for every machine~$M$ with $k$ work
tapes and value width at most~$\beta$ there is~$d_0$ such that for no $d\ge d_0$ do the following hold together at
$L=d^2b(d)$: $M$ is oblivious at length~$L$ with halting time $T\le C(L+1)$; every work-tape cell is entered at most $V$
times on $[0,T)$; and $M$ transposes all $d\times d$ matrices with $b(d)$-bit entries.
\end{corollary}

\begin{informalproof}
Let $d_0$ be given by the crux for $\eta=(\rho-\eps)/2$, and suppose the three conditions hold at some $d\ge d_0$. By
\cref{cor:owner-load} the long hops of the work-head walks have owner-load at most $2V-1$, so \cref{thm:F}, at
$\eps'=\eps/(4\beta C)\in(0,1/8]$, gives a cut of the walk DAG with at most $\eps'T$ edges and depth less than~$D_1$. By
\cref{lem:cut-genie} the dependency DAG has a genie~$\Gamma$ of cost at most $\eps L$ and full depth less than~$D$. As in the
proof of Theorem~D, $M$ appends exactly $L$ output bits, $I_\Gamma\in\mathcal S_{k,C,\beta,D}$, and the revealed word
$G_\Gamma$ is an index code for~$I_\Gamma$ with at most $\eps L$ bits, by the cone lemma~\cite[Corollary~4.5]{FourthPaper2026}
(a written proof there), while the
crux requires at least $(\rho-\eta)L>\eps L$ bits. The input head needs no hypothesis: it adds no edge to the dependency DAG.
\end{informalproof}

So the crux alone, at the depth~$D$ and a rate above~$\eps$, excludes oblivious linear-time transposers whose work heads
enter every cell $O(1)$ times, whatever the coupling, drift and scales of the heads and whatever the input head does.
Here $\log_2D\approx512k^3(2V-1)^3\beta C/\eps$, so the crux is needed at a rate exceeding about
$512k^3(2V-1)^3\beta C/\log_2D$, the analogue of the rate $456k^3(1+A)\beta C/\log_2D$ needed for home-based
heads~\cite[Remark~4.15]{FourthPaper2026}.

\begin{remark}[Which machines are new]\label{rem:calibration}
Bounded entry allows unboundedly many reversals. A machine whose work heads make at most $r$ reversals each enters every
cell at most $r+1$ times, since each head's motion splits into $r+1$ weakly monotone pieces, each of which visits a cell in
one run; a bound on the number of visits counted in time steps does not cover such machines, because a head may wait on a
cell. Such a head has at most $r+1$ legs, each contributing at most~$T$ to its leg area about its start cell, so it satisfies
the area hypothesis of the fourth paper's Corollary~B~\cite[Corollary~4.23]{FourthPaper2026} with $A=r+1$, and that corollary, like Corollary~F, puts no
condition on the input head. Relative to Corollary~B, Corollary~F is therefore new exactly for bounded-entry heads whose leg
area is superlinear about every home; such heads reverse unboundedly often. Conversely, home-based heads may enter cells
near home unboundedly often, so the two corollaries are incomparable. Relative to the unconditional bounds, Corollary~F
also allows the input head to move freely: machines whose heads, the input head included, make $o(\log d)$ reversals in
all cannot transpose at all, oblivious or not, by the first paper's pass bound, Lean-verified
there~\cite[Theorems~5.9 and~5.12]{FirstPaper2026}, and conservative reversal-bounded machines are excluded by Kirchherr's
bound~\cite{Kirchherr1988}.
\end{remark}

\subsection{Corollaries F\texorpdfstring{$'$}{'} and F\texorpdfstring{$''$}{''}}\label{sec:trans-Fprime}

Question~9.1 of \emph{Directed Recovery Structures} asks: for a fixed degree bound~$\Delta$ and recovery relations $N_j$ on
finite sets $V_j$ with $\card{V_j}\to\infty$, whose incidence graphs have maximum degree at most~$\Delta$ and satisfy
the paired-distance condition, must every sequence of storage codes have rate at most $\frac12+o(1)$, for every
alphabet?~\cite[Question~9.1]{DirectedRecovery2026} The incidence graph has vertices $L_p$ and $R_p$ for every
coordinate~$p$ and an edge $L_pR_q$ whenever $p\in N(q)$; the paired-distance condition says that, for every fixed~$r$, only
$o(\card{V_j})$ coordinates have $L_p$ and $R_p$ at distance at most~$r$~\cite[Definition~2.3]{DirectedRecovery2026}. The
same paper proves that the cone relations $N_\Gamma(q)=I_\Gamma(\tpos(q))$ of every schedule with $T\le c(L+1)$, $k$ work
tapes and a genie of fixed full depth satisfy this condition after the paid deletion of a small fixed fraction of
coordinates of high reuse~\cite[Theorem~3.3 and Corollary~3.8]{DirectedRecovery2026}, both partly formalized there, with
written proofs of the native realization and transport estimates and of the cone-size premise; this holds also for a single
frozen trace of an adaptive machine, as argued in prose in~\cite[Remark~3.9]{DirectedRecovery2026}. It proves, with a written
proof that is not formalized, that an affirmative answer to Question~9.1 for every~$\Delta$ gives rate at most
$\frac12+o(1)$ for all codes on these cone relations, uniformly over the schedules and
genies~\cite[Proposition~9.2]{DirectedRecovery2026}; the proof of that proposition uses the schedule only through the cone
bound, the paid trimming and the paired-distance theorem.

\begin{corollary}[Corollary F$'$: Question 9.1 excludes bounded-entry oblivious transposers]\label{cor:Fprime}
\notformalized
Suppose that Question~9.1 of~\cite{DirectedRecovery2026} has an affirmative answer for every fixed~$\Delta$; the binary
alphabet suffices. Then for all integers $k,V\ge1$ and every $c>0$, no machine with $k$ work tapes is, for infinitely
many~$d$, oblivious at length $L=d^2b(d)$, a transposer of all $d\times d$ matrices within $c(L+1)$ steps, and such that
every work-tape cell is entered at most $V$ times.
\end{corollary}

\begin{informalproof}
Let $M$ be such a machine, $\beta=\beta(M)$, $C=\max(c,1)$ and $\eps=1/4$. At each such~$d$, as in the proof of
Corollary~F, \cref{thm:F} and \cref{lem:cut-genie} give a genie~$\Gamma$ of cost at most $L/4$ and full depth at most a
constant~$D$ depending only on $k$, $V$, $\beta$ and~$C$. The revealed word is an index code for~$I_\Gamma$, so its fibres
are binary storage codes for the cone relation $N_\Gamma$~\cite[Remark~4.9(b)]{FourthPaper2026}, and some fibre has at least
$2^{3L/4}$ words. By~\cite[Proposition~9.2]{DirectedRecovery2026}, applied with $k$ work tapes, linear-time constant~$c$ and
cone depth~$D$, every binary storage code on these relations has at most $2^{(1/2+o(1))L}$ words. This is a contradiction
for large~$d$.
\end{informalproof}

For home-based schedules the same conclusion is the remark of \emph{Directed Recovery Structures} that every class of
schedules with arbitrarily cheap genies of fixed depth would be excluded by the all-code half bound
\cite[Section~9.1]{DirectedRecovery2026}, combined with the fourth paper's Corollary~B; the bounded-entry class is supplied by
Theorem~F. The same argument extends to adaptive machines whose traces have entropy at most $\eta L$ for a fixed
$\eta<\frac12$, provided that every trace of probability at least $2^{-\eta L}$ has bounded entry.

\begin{corollary}[Corollary F$''$: low adaptivity]\label{cor:Fsecond}
\notformalized
Let $k,V\ge1$ be integers, $c>0$ and $0\le\eta<1/2$. Suppose that Question~9.1 of~\cite{DirectedRecovery2026} has an
affirmative answer for every fixed~$\Delta$. Then for every machine~$M$ with $k$ work tapes there is~$L_0$ such that at no
length $L=d^2b(d)\ge L_0$ does $M$ have the following three properties:
\begin{enumerate}[label=(\roman*)]
\item on every input~$X$ it halts within $c(L+1)$ steps with output $X^\T$;
\item every trace of probability at least $2^{-\eta L}$, for a uniformly random input, enters every work-tape cell at most
$V$ times (for instance, every trace does);
\item the trace $\tau(X)$ has Shannon entropy $H(\tau(X))\le\eta L$ for a uniformly random input (for instance, at most
$2^{\eta L}$ distinct traces occur).
\end{enumerate}
For $\eta=0$ this is Corollary~F$'$. The proof uses only one trace: (ii) and~(iii) may be replaced by the single condition
that some trace of probability at least $2^{-\eta L}$ enters every work-tape cell at most $V$ times.
\end{corollary}

\begin{informalproof}
Fix $\eps\in(0,1]$ with $\eta+\eps<\frac12$, and let $D(\eps)$ be the full depth that \cref{thm:F} and
\cref{lem:cut-genie} give, as in the proof of Corollary~F$'$, for $k$, $V$, $\beta(M)$, $c$ and~$\eps$. Suppose that $M$ has the
three properties at a length~$L$. Since $H(\tau(X))\ge-\log_2\max_\tau\Pr[\tau(X)=\tau]$, some trace~$\tau_0$ has
probability at least $2^{-\eta L}$; let $W$ be the set of inputs that follow it, so $\card W\ge2^{(1-\eta)L}$. All runs on~$W$
share the dependency DAG $H_{\tau_0}$ of the trace (\cref{sec:prelim-tcp}), and the local recursion and the cone lemma of the
fourth paper~\cite[Lemma~4.2, Lemma~4.4 and Corollary~4.5]{FourthPaper2026}, which have written proofs there, hold for all
pairs of inputs in~$W$, because their proofs use only the common schedule. By~(ii), $\tau_0$ has bounded entry, so
\cref{thm:F} and \cref{lem:cut-genie}, applied to $H_{\tau_0}$, give a genie~$\Gamma$ of cost at most $\eps L$ and full depth
at most~$D(\eps)$. For a value~$g$ of the revealed word,
the set $W_g$ of inputs in~$W$ with revealed word~$g$ is a storage code for the cone relation of $\tau_0$: two of its inputs that
agree on $N_\Gamma(q)$ produce the same output bit at position $\tpos(q)$, which by correctness is their bit~$q$. Some~$g$ has
$\card{W_g}\ge2^{(1-\eta-\eps)L}$. The paired-distance theorem and its corollaries hold for one frozen trace
\cite[Remark~3.9]{DirectedRecovery2026}, and the proof of~\cite[Proposition~9.2]{DirectedRecovery2026} uses the schedule only
through them, so it bounds $\log_2\card{W_g}$ by $(\frac12+o(1))L$, uniformly in the trace, with $o(1)$ depending only on
$k$, $c$ and $D(\eps)$. Since $\eta+\eps<\frac12$, this is a contradiction for all $L\ge L_0$, for a suitable $L_0$.
\end{informalproof}

The trace's information is paid for once, in the choice of~$\tau_0$. This is the kind of paid treatment of adaptivity that
the fourth paper anticipated, with a coding bound for codes that are correct only on a large set of
inputs~\cite[Remark~4.24]{FourthPaper2026}, and that \emph{Directed Recovery Structures} leaves open in
general~\cite[Section~9.3]{DirectedRecovery2026}; here it is available only for traces of entropy below~$L/2$.

\paragraph{The entropy hypothesis is not robust.} Hypothesis~(iii) can be broken by changes that do not affect the
computation. Given a transposer~$M$, let $M'$ first sweep its input head across the input, adding a step out and back
at every cell of the first half of the input that holds a~$1$, and then run~$M$. Then $M'$ is a transposer, it takes
$O(L)$ more steps, and its work-tape entries are those of~$M$, so (i) and~(ii) are unaffected; but its trace determines
the first half of the input, so $H(\tau(X))\ge L/2$ and (iii) fails for every $\eta<1/2$. A form of Corollary~F$''$
that survives such padding would bound the entropy of a coarse skeleton of the trace rather than that of the trace.

\begin{remark}[A weak form of Question 9.1 suffices]\label{rem:weak-crux}
Corollary~F$'$ does not need rate $\frac12$. The paired-distance condition is a property of sequences of relations.
Suppose that along every sequence of relations of incidence degree at most~$\Delta$ that satisfies it, all codes have rate at
most $1-\rho(\Delta)+o(1)$. The argument of Corollary~F$'$, with the paid trimming
of~\cite[Corollary~3.8]{DirectedRecovery2026} at the threshold $\lceil B/\delta\rceil$, where $B=\sum_{i\le D(\eps)}(k+1)^i$
bounds the cone sizes, then works as soon as $(1-\delta)\rho(\Delta)>\eps$ for some $\delta\in(0,1)$ and $\eps\in(0,1]$
with $\Delta=\lceil B/\delta\rceil$. Since $D(\eps)$ is exponential in $1/\eps$, $\log\log\Delta$ is of order $1/\eps$, and a
deficit with $\rho(\Delta)\log\log\Delta\to\infty$ suffices. For Corollary~F$''$ the deficit must exceed $\eta$ as well.
\emph{Directed Recovery Structures} makes the same qualitative point: a deficit works if the cost of the genie can be chosen
below it before the depth is fixed~\cite[Section~9.1]{DirectedRecovery2026}.
\end{remark}

\subsection{One unrestricted head}\label{sec:trans-M}

\begin{corollary}[Machine form of Theorem M]\label{cor:M-machine}
\notformalized
Let $k\ge2$, $\beta\ge2$, $\ell\ge1$, $q\ge1$ and $h\ge0$ be integers, $C\ge1$ and $0<\eps\le1$, and put
$D=\DM(k-1,\ell,q,h,\eps/(4\beta C))\cdot\lceil4\beta C/\eps\rceil$. Suppose that $\Crux(\mathcal S_{k,C,\beta,D},\rho)$ holds
for some $\rho>\eps$. Then for every machine~$M$ with $k$ work tapes and value width at most~$\beta$ there is~$d_0$ such that
for no $d\ge d_0$ is $M$ oblivious at length $L=d^2b(d)$ with halting time at most $C(L+1)$, a transposer of all
$d\times d$ matrices, and such that all its work heads but one have long hops of owner-load at most~$\ell$ whose family bands
are covered by $q$ windows of width~$h$. The same holds with Question~9.1, in the form of Corollary~F$'$.
\end{corollary}

\begin{informalproof}
The proof of Corollary~F uses Theorem~F only through a cut of the walk DAG with at most $\eps T/(4\beta C)$ edges and depth less
than~$D_1$. \Cref{thm:M}, applied with the unrestricted work head as~$x_0$ and the long hops of the other $k-1$ work heads
as~$H'$, at $\eps/(4\beta C)$, supplies such a cut with $D_1=\DM(k-1,\ell,q,h,\eps/(4\beta C))$.
\end{informalproof}

Corollary~F also holds with an owner-load hypothesis on all work heads instead of bounded entry, with
$D_1=\DF(k,\ell,\eps/(4\beta C))$, since its proof uses only the owner-load. Together the two corollaries say the following.
Let $0<\eps\le1$ and $\rho>\eps$, and let $M$ be oblivious with $k\ge2$ work tapes and value width at most~$\beta$,
transposing within $C(L+1)$ steps. For fixed $\ell$, $q$ and~$h$, if $\Crux(\mathcal S_{k,C,\beta,D},\rho)$ holds at the depth
$D=\DF(k,\ell,\eps/(4\beta C))\cdot\lceil4\beta C/\eps\rceil$ and at the depth~$D$ of \cref{cor:M-machine}, then at all
large lengths $M$ has at least two work heads of owner-load greater than~$\ell$, or exactly one, while the family bands of
the other work heads are not covered by $q$ windows of width~$h$. Which of the two a transposer would use is open. The first paper's pass bound forces $\Omega(\log d)$
reversals~\cite[Theorem~5.12]{FirstPaper2026}, but one head of owner-load one and a single scale supplies them: a zigzag of
amplitude about $L/\log d$ has one family band, up to the families truncated at the ends, and makes $\log d$ reversals of
amplitude larger than any smaller scale. So the pass bound does not force the case of many scales.

\section{Linear machines: the linear crux as rigidity}\label{sec:linear}

For machines whose data path is linear over a field, Theorem~D has a linear analogue, Theorem~$\Dlin$ (\cref{thm:Dlin}):
\emph{linear depth-TCP} together with a \emph{linear crux} excludes oblivious linear-time transposers. Linear depth-TCP weakens
the structural hypothesis, replacing depth by \emph{linear} depth, in which long paths may cancel; the linear crux is required
for a larger class of systems, the ball systems of the schedules, and it can be stated exactly as relative rigidity of the
transposition matrix with respect to the dependency pattern of the schedule. Relative rigidity alone excludes no machine.

\subsection{Linear data paths and path sums}\label{sec:lin-setting}

Let $F$ be a field; input bits are read as $0,1\in F$. Let $M$ be oblivious at length~$L$, with dependency DAG $H_M$ and step
values $v_t(w)$ as in \cref{sec:prelim-tcp}, and write $\mathrm{out}(t)$ for the set of defined out-neighbours $t-1$,
$\prev_1(t),\dots,\prev_k(t)$ of~$t$ in~$H_M$.

\begin{definition}[Linear data path]\label{def:linear-path}
\partformal{Rigidity.lean}{LinNet, LinNet.val}{the machine, its step values and their encodings; in Lean the data path is
given as a network, and the output steps and the functionals $\varphi_\ell$ are arguments of \lean{LinNet.rigidity}.}
$M$ has an \emph{$F$-linear data path} at length~$L$ if there are maps $\mathrm{enc}_t$ from the possible values of step~$t$ to
$F^m$, vectors $c_t,a_t\in F^m$ and matrices $A_{t,s}\in F^{m\times m}$, none depending on the input, such that for every
$w\in\{0,1\}^L$
\[
  \mathrm{enc}_t(v_t(w))=c_t+\sum_{s\in\mathrm{out}(t)}A_{t,s}\,\mathrm{enc}_s(v_s(w))+a_t\,w_{x_0(t)},
\]
where the last term is present only if the input head reads a position $x_0(t)\in[0,L)$, and such that output bit~$\ell$,
appended at step~$t_\ell$, equals $\varphi_\ell(\mathrm{enc}_{t_\ell}(v_{t_\ell}(w)))$ for a fixed affine functional
$\varphi_\ell$ with linear part $\varphi_\ell^{\mathrm{lin}}$. Such a \emph{representation} need not be unique.
\end{definition}

Write $v_t$ for $\mathrm{enc}_t(v_t)$. For a genie~$\Gamma$, expanding the recursion along paths and stopping at steps
of~$\Gamma$ gives
\[
  v_t=c'_t+\sum_{s\in\Gamma}T_{t,s}v_s+\sum_jC_{t,j}w_j .
\]
Here, for a path $\pi=(t=t_0\to\dots\to t_n)$ of $H_M$ with weight $W(\pi)=A_{t_0,t_1}\cdots A_{t_{n-1},t_n}$, $T_{t,s}$ is
the sum of $W(\pi)$ over the paths from~$t$ to $s\in\Gamma$ with $n\ge1$ whose interior avoids~$\Gamma$; $C_{t,j}$ is the sum
of $W(\pi)a_{t_n}$ over the $\Gamma$-free paths whose last step reads position~$j$; and $c'_t$ is the sum of $W(\pi)c_{t_n}$
over all $\Gamma$-free paths from~$t$. The sums are finite because $H_M$ is acyclic. In matrix form they are the blocks of
$(I-A_\Gamma)^{-1}(A-A_\Gamma)$ and $(I-A_\Gamma)^{-1}B$, where $A_\Gamma$ is the block matrix~$A$ with the columns of~$\Gamma$
set to zero and $B$ is the input matrix: the transfer matrices of the linear network cut at~$\Gamma$, in the sense of Koetter and
Médard~\cite{KoetterMedard2003}.

\begin{definition}[Linear depth]\label{def:linear-depth}
\partformal{Rigidity.lean}{LinNet.GFree, LinNet.C, LinNet.reach}{the condition of linear depth itself, which in Lean is a
hypothesis of \lean{LinNet.rigidity}.}
Let $I^D_\Gamma(\ell)$ be the set of input positions read at steps reachable from~$t_\ell$ by $\Gamma$-free paths of length
at most~$D$, and $B^D(\ell)\supseteq I^D_\Gamma(\ell)$ the set of positions read at steps reachable from~$t_\ell$ by directed paths of
length at most~$D$ in~$H_M$, so that $B^D(\ell)=I^D_\emptyset(\ell)$. Output step~$t_\ell$ has \emph{linear depth at most~$D$ for~$\Gamma$} if
$\varphi_\ell^{\mathrm{lin}}C_{t_\ell,j}=0$ for every $j\notin I^D_\Gamma(\ell)$.
\end{definition}

For such~$j$ every $\Gamma$-free path to a read of~$j$ is longer than~$D$, so the condition says that the long contributions
to inputs outside $I^D_\Gamma(\ell)$ cancel; output bit~$\ell$ is then an affine function of the genie values and of the
bits~$w_j$, $j\in I^D_\Gamma(\ell)$. Structural depth at most~$D$ implies linear depth at most~$D$ at every output step.
Linear depth is not monotone in~$\Gamma$: adding a step to~$\Gamma$ can break a cancellation between a path through it and a
path that avoids it.

\subsection{Rigidity relative to the schedule}\label{sec:lin-rigid}

Let $P_{\tpos}$ be the $L\times L$ matrix whose row~$\ell$ is the unit vector $e_{\tpos(\ell)}$. The \emph{minrank} of a
system $(I(\ell))_\ell$ over~$F$ is the least rank of a matrix~$R$ such that row~$\ell$ of $P_{\tpos}-R$ is supported in
$I(\ell)$ for every~$\ell$; it is the least length of an $F$-linear index code for the system, decodable on all of~$F^L$ with
linear decoders~\cite{Haemers1979,BBJK2011}.

\begin{lemma}[Per-machine rigidity]\label{lem:rigidity}
\partformal{Rigidity.lean}{LinNet, LinNet.val, LinNet.C, LinNet.reach, LinNet.reach_subset, LinNet.Pmat, LinNet.minrank, LinNet.rigidity, LinNet.minrank_le}{the
translation from the machine to a linear network: the machine model, the encodings $\mathrm{enc}_t$ and the step
values~$v_t$.}
Suppose that $M$ has an $F$-linear data path at length~$L$ (with some representation), transposes all $d\times d$ matrices,
and has a genie~$\Gamma$ under which every output step has linear depth at most~$D$. Then
\[
  P_{\tpos}=R+S,\qquad\rank R\le\beta\card\Gamma,\qquad\supp(\text{row }\ell\text{ of }S)\subseteq I^D_\Gamma(\ell)
  \subseteq B^D(\ell);
\]
equivalently, the system $(I^D_\Gamma(\ell))_\ell$ has minrank at most $\beta\card\Gamma$ over~$F$.
\end{lemma}

\begin{outline}
The Lean statement concerns a linear network, \lean{LinNet}, with out-neighbours at earlier steps, the vectors $c_t$ and
$a_t$, the matrices $A_{t,s}$, and the input position read by each step, if any. Its formal values $V_t(w)$, $w\in F^L$, are
defined by the recursion of the data path (\lean{LinNet.val}), and the hypotheses are stated for them: on binary inputs the
affine output functionals give the transposed bits, each step of~$\Gamma$ takes at most $2^\beta$ values, and every output
step has linear depth at most~$D$, with the path sums \lean{LinNet.C} over the $\Gamma$-free paths (\lean{LinNet.GFree}) and
the sets $I^D_\Gamma(\ell)$ given by \lean{LinNet.reach}; for $\Gamma=\emptyset$ these are the ball sets, and
$I^D_\Gamma(\ell)\subseteq B^D(\ell)$ (\lean{LinNet.reach_subset}).

The values split into a constant and a linear part, $V_t(w)=V_t(0)+V^{\mathrm{lin}}_t(w)$ (\lean{LinNet.val_eq_add}), and the
linear part has the path-sum expansion
$V^{\mathrm{lin}}_t(w)=\sum_{s\in\Gamma}T_{t,s}V^{\mathrm{lin}}_s(w)+\sum_jC_{t,j}w_j$ identically in $w\in F^L$, being the
unrolled recursion (\lean{LinNet.C_eq}, \lean{LinNet.vlin_expand}); in Lean $T_{t,s}$ is defined by the recursion that it
satisfies (\lean{LinNet.Tm}). Let $E(w)=(V^{\mathrm{lin}}_s(w))_{s\in\Gamma}$, the linear parts of the revealed values, let
$R$ be the matrix of the linear map $w\mapsto\bigl(\sum_{s\in\Gamma}\varphi_\ell^{\mathrm{lin}}T_{t_\ell,s}E_s(w)\bigr)_\ell$,
and let $S_{\ell j}=\varphi_\ell^{\mathrm{lin}}C_{t_\ell,j}$. Transposition at the inputs $0$ and~$e_j$ gives
$\varphi_\ell^{\mathrm{lin}}V^{\mathrm{lin}}_{t_\ell}(e_j)=(e_j)_{\tpos(\ell)}$, and the expansion turns this into
$(P_{\tpos})_{\ell j}=R_{\ell j}+S_{\ell j}$. By linear depth, $S_{\ell j}=0$ for $j\notin I^D_\Gamma(\ell)$. On $\{0,1\}^L$,
$E$ takes at most $2^{\beta\card\Gamma}$ values, since each $V_s$, $s\in\Gamma$, takes at most $2^\beta$ values there. A
linear map of rank~$\rho$ takes at least $2^\rho$ values on $\{0,1\}^L$: it is injective on the binary vectors supported on
$\rho$ coordinates whose columns are independent, because differences of such vectors have coefficients in $\{-1,0,1\}$
(\lean{LinNet.two_pow_finrank_le}). So $\rank R\le\rank E\le\beta\card\Gamma$ (\lean{LinNet.rigidity}); in other words, $E$
is an $F$-linear index code of rank at most $\beta\card\Gamma$ for $(I^D_\Gamma(\ell))_\ell$, decodable with linear decoders.
The minrank form is \lean{LinNet.minrank_le}, with the minrank as the least such rank (\lean{LinNet.minrank}).
\end{outline}

\begin{proof}[Proof of the translation (not formalized)]
The dependency DAG of~$M$, with out-neighbours $\mathrm{out}(t)$, the data $c_t,A_{t,s},a_t$ of the representation and the
input positions read define a linear network. By induction on~$t$, its formal values satisfy $V_t(w)=\mathrm{enc}_t(v_t(w))$
for $w\in\{0,1\}^L$, since both satisfy the recursion. As $M$ transposes, it appends exactly $L$ bits, and output bit~$\ell$
equals $\varphi_\ell(V_{t_\ell}(w))$ on binary inputs; each $v_s$ takes at most $2^\beta$ values; and linear depth is a
statement about the path sums of the representation.
\end{proof}

\begin{definition}[Linear depth-TCP and the linear crux]\label{def:lin-hyp}
\emph{Linear depth-TCP}$(k,C,\beta,\eps,D;F)$ says: for every machine~$M$ with $k$ work tapes and value width at most~$\beta$
there is~$d_0$ such that at every $d\ge d_0$ at which $M$ is oblivious at length~$L$, has an $F$-linear data path and halting
time at most $C(L+1)$, there are a representation of its data path and a genie~$\Gamma$ with $\beta\card\Gamma\le\eps L$ under
which every output step has linear depth at most~$D$. Let $\Sball_{k,C,\beta,D}$ be the class of the ball systems
$(B^D(\ell))_{\ell\in[0,L)}$, at the lengths $L=d^2b(d)$, of the oblivious machines with $k$ work tapes, value width at
most~$\beta$ and halting time at most $C(L+1)$ that append exactly $L$ bits, as for $\mathcal S_{k,C,\beta,D}$; no linear data
path is required. $\LinCrux_F(\Sball_{k,C,\beta,D},\rho)$ says: for every $\eta>0$ there is~$d_0$ such
that every system of $\Sball_{k,C,\beta,D}$ at a size $d\ge d_0$ has minrank at least $(\rho-\eta)L$ over~$F$.
\end{definition}

\begin{theorem}[Theorem $\Dlin$]\label{thm:Dlin}
\partformal{Rigidity.lean}{LinNet.dlin_step}{the quantification over machines and sizes, the two hypotheses as statements
about machines, and the translation of a machine into a linear network (\cref{lem:rigidity}); the Lean theorem is the
argument at one size, for a network.}
Let $k\ge1$, $D\ge0$ and $\beta\ge2$ be integers, $C>0$, $0<\eps<\rho$ and $F$ a field. If linear depth-TCP$(k,C,\beta,\eps,D;F)$
and $\LinCrux_F(\Sball_{k,C,\beta,D},\rho)$ hold, then for every machine~$M$ with $k$ work tapes and value width at
most~$\beta$ there is~$d_0$ such that for no $d\ge d_0$ is $M$ oblivious at length~$L$ with an $F$-linear data path, halting
time at most $C(L+1)$, and a transposer of all $d\times d$ matrices.
\end{theorem}

\begin{outline}[Outline of the formal proof at one size]
For a transposing network with a genie~$\Gamma$ of cost $\beta\card\Gamma\le\eps L$ under which every output step has
linear depth at most~$D$, \cref{lem:rigidity} bounds the minrank of the ball system by $\beta\card\Gamma\le\eps L$; with
$\eta=(\rho-\eps)/2$ and $L\ge1$ this contradicts a minrank of at least $(\rho-\eta)L$ (\lean{LinNet.dlin_step}).
\end{outline}

\begin{proof}[Proof of the reduction to one size (not formalized)]
Let $\eta=(\rho-\eps)/2$ and let $d_0$ be the larger of the thresholds of the two hypotheses. If $M$ were such a transposer at
some $d\ge d_0$, linear depth-TCP would give a representation and a genie~$\Gamma$ with $\beta\card\Gamma\le\eps L$ under which
every output step has linear depth at most~$D$, and $\LinCrux_F$ would give the ball system
$(B^D(\ell))_\ell\in\Sball_{k,C,\beta,D}$ a minrank of at least $(\rho-\eta)L$; the argument at one size, applied to the
network of the representation, gives a contradiction.
\end{proof}

This is the proof of Theorem~D with the cone lemma replaced by the path-sum expansion, which is exact for linear data paths.

\paragraph{Why it matters.} Read as a matrix statement, $\LinCrux_F(\Sball,\rho)$ says that $P_{\tpos}$ is $\rho L$-rigid
relative to the $D$-ball pattern of every linear-time schedule: it is not the sum of a matrix of rank below $\rho L$ and a
matrix supported on that pattern. So, for linear machines, linear depth-TCP together with a
linear crux excludes oblivious linear-time transposers, and the linear crux can be stated exactly as relative rigidity of the
transposition matrix; \cref{lem:rigidity} applies only to a machine that has a cheap genie under which every output has
linear depth at most~$D$. This is a relative form of Valiant's route from depth reduction and rigidity to lower bounds for
linear circuits~\cite{Valiant1977}. The route from small depth to index coding also goes back to Riis's guessing
games~\cite{Riis2007}, and it was used in the language of min-rank by Jukna and Schnitger~\cite{JuknaSchnitger2011}. For
static data structures, lower bounds imply rigidity~\cite{DGW2019}, and for systematic linear data structures the two are
equivalent~\cite{NRR2020}. Without the pattern the statement is false for permutation matrices: splitting the receivers into
groups of $s+1$, each knowing the bits demanded by the others in its group, gives rank $\lceil L/(s+1)\rceil$ with $s$
off-diagonal entries per row~\cite[Section~4.3]{FourthPaper2026}. Any proof must use the schedule's structure. On the construction
side, a linear oblivious transposer at size~$d$ either has no genie of cost $\eps L$ that makes every output linearly
$D$-shallow, or its formal revealed map is a linear index code of rank at most $\eps L$ for its own ball system.

\begin{remark}[Scope]\label{rem:Dlin-scope}
(a) The genie plays no role in the crux class: since $I^D_\Gamma(\ell)\subseteq I^D_\emptyset(\ell)=B^D(\ell)$ and the empty
genie is allowed, the linear crux for the systems $(I^D_\Gamma(\ell))_\ell$ over all genies is equivalent to the linear crux
for the ball systems; and $\Sball$ dominates the fourth paper's class $\mathcal S_{k,C,\beta,D}$ pointwise. Neither
Theorem~$\Dlin$ nor the linear form of Theorem~D implies the other: $\Dlin$ trades a weaker structural hypothesis for a larger
crux class. For home-based linear machines $\Dlin$ gives nothing new, since Corollary~B supplies structural genies.
(b) The trade-off between cost and depth survives: the machine that appends, at every step, the parity of all bits read so
far forces at least $\lfloor(L-1)/(D+1)\rfloor$ revealed steps for linear depth~$D$, as in the fourth paper's forced trade-off
\cite[Proposition~4.13]{FourthPaper2026}.
(c) Linear depth is robust to dummy traffic on separate coordinates: copies of an input bit carried along deep paths of a
dummy schedule on extra tapes and combined so that they cancel contribute $W(\pi_1)+W(\pi_2)=0$ to the path sums, whatever their
lengths, provided the genie reveals single coordinates (fields, as in~\cite[Remark~4.6]{FourthPaper2026}) and does not cut one
of the two paths; Lemma~\ref{lem:rigidity} and Theorem~$\Dlin$ hold verbatim for such coordinate genies, with cost equal to the
total bit width of the revealed coordinates. With whole-step genies the robustness fails, and it fails for dummy traffic that is
entangled with revealed coordinates, for instance a value masked by a pad and unmasked through a deep partner copy.
(d) Linear depth-TCP quantifies over all linear machines, and the theorem uses it only for the transposer itself. It is not a
statement about head schedules alone: it depends on the data path and on its representation.
\end{remark}

\section{The crux in special linear settings}\label{sec:crux}

The crux asks for a deficit below rate~$1$, and its undirected core asks whether storage codes on locally bipartite side
information of bounded degree can beat rate~$\frac12$~\cite[Question~6.1]{FourthPaper2026}. For all codes this is open beyond
the classes of the companion paper, \emph{Nested Elimination and Half-Rate Bounds for Storage Codes}~\cite{SeventhPaper2026},
which proves the half bound for every storage code on triangle-free graphs of maximum degree three, among other classes, with
the guessing number equal to the fractional matching number there. This section proves two linear cases in which a natural
source of high-rate codes fails: codes lifted along towers of covers in characteristic~$2$ (\cref{sec:crux-towers}), and scalar
codes at degree four whose exceptional vertices are few or far apart (\cref{sec:crux-deg4}). By \cref{sec:linear}, linear codes
are the relevant ones for linear machines.

\subsection{Lifted codes along towers of covers}\label{sec:crux-towers}

Let $H$ be a finite connected graph with vertex set~$V$, $\card V=N$, and let $M\in\F_2^{V\times V}$ have unit diagonal and
off-diagonal support in the edges of~$H$; the entries $M_{uv}$ and $M_{vu}$ may differ. Its kernel is a linear storage code
on~$H$~\cite[Lemma~6.14]{FourthPaper2026}. Let $\mathcal T$ be the universal cover of~$H$, a tree, and $\pi=\pi_1(H)$, a free
group acting freely on~$\mathcal T$. A \emph{finite cover} is $\mathcal T/\Lambda$ for a subgroup $\Lambda\le\pi$ of finite
index; more generally, a finite $\pi$-set~$X$ defines a cover with vertex set $X\times V$. Choosing a spanning tree and a base
vertex turns~$M$ into a matrix $\tilde M\in\mathrm{Mat}_V(\F_2[\pi])$, whose entry $(u,v)$ is the sum of $M_{uv}g_e$ over the
edges~$e$ from $u$ to~$v$, where $g_e\in\pi$ is the closed walk along the tree to~$u$, across~$e$ and back along the tree, and
$\tilde M_{vv}=1$. The \emph{lift} $M_X$ is the image of $\tilde M$ under the permutation representation of~$\pi$ on
$\F_2[X]$, applied entrywise: a matrix with unit diagonal and support in the edges of the cover, whose kernel is a linear
storage code there, of rate $\dim\ker M_X/(N\card X)$. The \emph{inner rank} $\irk\tilde M$ is the least~$\rho$ such that
$\tilde M=PQ$ with $P\in\mathrm{Mat}_{N\times\rho}(\F_2[\pi])$ and $Q\in\mathrm{Mat}_{\rho\times N}(\F_2[\pi])$. The group algebra
of a free group is a free ideal ring, and its inner rank equals the rank over its universal field of fractions, Cohn's free
field~\cite{Cohn2006}. The \emph{universal rate} is
\[
  r_U(M)=\bigl(N-\irk\tilde M\bigr)/N .
\]
Let $\chi\colon\pi\to\Z/2$ be the parity of the length of closed walks, $\pi^+=\ker\chi$, $\hat\pi$ the pro-$2$ completion of
$\pi$ and $\hat\pi^+$ the closure of~$\pi^+$ in it. A cover $\mathcal T/\Lambda$ is bipartite if and only if $\Lambda\le\pi^+$.

\begin{theorem}[The equivariant crux in characteristic two]\label{thm:eqcrux}
\partformal{TwistPairing.lean}{lift2, twistPairing, finrank_ker_lift2, excess_lift2, rate_lift2_le}{parts (a), (b), (d),
(e) and~(f), and in~(c) the passage from the tower to its levels: the Lean theorems are the step from one level to the next,
for a double cover given by its twisted pairs.}
Let $H$ be non-bipartite.
\begin{enumerate}[label=(\alph*)]
\item Every finite cover~$X$ satisfies $\operatorname{rate}(M_X)\ge r_U(M)$.
\item $r_U(M)\le\frac12$.
\item Along every tower $\pi=\Lambda_0\ge\Lambda_1\ge\cdots$ in which each $\Lambda_{n+1}$ has index~$2$ in~$\Lambda_n$, the rate
is non-increasing. More precisely, with $N_n=2^nN$ and $M_n$ the lift at level~$n$, the excess $e_n=\dim\ker M_n-N_n/2$
satisfies $e_{n+1}=2e_n-r_n$, where $r_n\ge0$ is the rank of the \emph{twist pairing}
$\ker M_n^\T\times\ker M_n\to\F_2$, $(y,x)\mapsto\sum(M_{uv}y_ux_v+M_{vu}y_vx_u)$, the sum over the edges $uv$ of level~$n$
that the double cover twists.
\item Let $C=\bigcap_n\overline{\Lambda_n}\le\hat\pi$. Some level of the tower is bipartite if and only if $C\le\hat\pi^+$, and
from that level on the rate is at most~$\frac12$.
\item Let $\iota\colon\pi\to\mathbf F$ be an injective homomorphism into a free-by-$\Z_2$ pro-$2$ group~$\mathbf F$, where
$\Z_2$ denotes the $2$-adic integers (for instance a finitely generated free pro-$2$ group), and let
$\mathbf F=\mathbf N_1\ge\mathbf N_2\ge\cdots$ be open normal subgroups with trivial intersection. Then
$\operatorname{rate}(M_{\iota^{-1}(\mathbf N_j)})\to r_U(M)$.
\item For the Clebsch graph and $M=I+A$, $r_U(M)=0$.
\end{enumerate}
\end{theorem}

\begin{outline}[Outline of the formal proof of the step in~(c)]
In Lean the double cover of level~$n$ has vertex set $V_n\times\{0,1\}$ and is given by its twisted pairs; the lift
(\lean{lift2}) connects equal sheets by the untwisted entries of~$M_n$ and opposite sheets by the twisted ones, whose matrix
is $M_n^{\mathrm{tw}}$ (\lean{twistPart}). The field is any field of characteristic~$2$. In the coordinates $a=x_0+x_1$,
$b=x_1$ (\lean{sheetEquiv}) the lift becomes $(a,b)\mapsto(M_na,M_nb+M_n^{\mathrm{tw}}a)$ (\lean{blockMap_sheetEquiv}).
(Equivalently, in $\F_2[\Z/2]\cong\F_2[t]/(t^2)$ with $g=1+t$, the lift is $\varphi=M_n+tM_n^{\mathrm{tw}}$ in the basis
$1,t$; the Lean proof works with the block map.) Its kernel consists of the
pairs $(a,b)$ with $M_na=0$ and $M_nb=M_n^{\mathrm{tw}}a$. Since $\operatorname{im}M_n=(\ker M_n^\T)^\perp$
(\lean{mem_range_iff_dotProduct_ker_transpose}), the admissible~$a$ form the kernel of the twist pairing
$x\mapsto(y\mapsto y^\T M_n^{\mathrm{tw}}x)$ on $\ker M_n$ (\lean{twistPairing}), of dimension $\dim\ker M_n-r_n$, and each
admits a coset of $\ker M_n$ of choices of~$b$ (\lean{finrank_ker_blockMap}). So $\dim\ker M_{n+1}=2\dim\ker M_n-r_n$
(\lean{finrank_ker_lift2}), which is the recursion $e_{n+1}=2e_n-r_n$ (\lean{excess_lift2}), and the rate does not increase
(\lean{rate_lift2_le}). In Lean the twisted pairs are an arbitrary relation; for a double cover they are the two
orientations of the twisted edges, and then $y^\T M_n^{\mathrm{tw}}x$ is the sum in~(c), which belongs to the passage from the
tower to its levels.
\end{outline}

\begin{proof}[Proof of (a), (b), (d), (e), (f) and of the tower in~(c) (not formalized)]
(a) Write $\tilde M=PQ$ with inner dimension $\rho=\irk\tilde M$. The permutation representation is a ring homomorphism, so
$M_X=P_XQ_X$ has rank at most $\card X\rho$ and nullity at least $\card X(N-\rho)$.

(b) Apply~(a) to the bipartite double cover $H\times K_2=\mathcal T/\pi^+$, on which every edge is twisted. With $M=I+B$, its
lift is $\begin{psmallmatrix}I&B\\B&I\end{psmallmatrix}$, and $(x,y)$ lies in its kernel exactly when $x=By$ and $(I+B^2)y=0$;
so its nullity is the nullity of $M^2=I+B^2$, at most~$N$. Hence $2Nr_U(M)\le N$. The bound is also contained in
\emph{Covers and Local Generation}~\cite[Proposition~3.2]{Covers2026}, a lower bound $\lceil N/2\rceil$ for the inner rank
over every field, whose separation hypothesis holds for graph templates and whose proof uses the specialization argument
of~(a).

(c) A double cover of level~$n$ is given by voltages in $\Z/2=\{1,g\}$ on the edges: the subgroup $\Lambda_{n+1}$ of
index~$2$ in~$\Lambda_n$ defines the double cover $\mathcal T/\Lambda_{n+1}\to\mathcal T/\Lambda_n$, and the lift $M_{n+1}$ is
the lift of~$M_n$ along it, with the twisted edges given by the voltages. The step above then gives the recursion at every
level. For symmetric checks, the
adjacency checks $M=I+A$, the recursion is the lift identity of \emph{Covers and Local Generation}~\cite[Lemma~9.1]{Covers2026},
for arbitrary binary lifts of finite connected graphs; the monotonicity of the normalized nullity under refinements with
$2$-group kernels in characteristic~$2$ is its Lemma~4.1, and it is also the index-$2$ case of the argument of Bergeron,
Linnell, Lück and Sauer~\cite[Lemma~4.1 and Theorem~1.6]{BLLS2014} (see also Kionke~\cite[Lemma~2.6]{Kionke2017}). Part~(c)
gives the exact identity, with the twist pairing, for every unit-diagonal, possibly asymmetric, binary check, along arbitrary
towers of index-$2$ subgroups.

(d) If every level were non-bipartite, the sets $\overline{\Lambda_n}\cap(\hat\pi\setminus\hat\pi^+)$ would form a decreasing
sequence of nonempty compact sets, so $C\not\le\hat\pi^+$. Conversely, a bipartite level gives
$C\le\overline{\Lambda_n}\le\hat\pi^+$. Covers of bipartite graphs are bipartite, and every storage code on a bipartite graph has
rate at most~$\frac12$~\cite[Corollary~6.9]{FourthPaper2026}.

(e) A free pro-$2$ group of positive rank maps onto~$\Z_2$ with a free pro-$2$ kernel, so it is free-by-$\Z_2$. By
Jaikin-Zapirain and Souza~\cite[Theorem~8.1]{JaikinSouza2026}, applied to the finitely generated subgroup $\iota(\pi)\cong\pi$,
the normalized dimensions of the kernels converge to the dimension over the universal division ring of~$\F_2[\iota(\pi)]$; in the
case of a finitely generated free pro-$2$ group, which is all that \cref{prop:eqtower} uses, this is
Jaikin-Zapirain's~\cite[Proposition~3.5]{JaikinZapirain2024}. Since $\iota$ is injective, this division ring is Cohn's free field
of $\F_2[\pi]$, and the limit is $r_U(M)$.

(f) The Clebsch graph contains the $4$-cube, whose Gray code is a Hamiltonian cycle, so it has a spanning $2$-factor, and
$I+A$ has unit diagonal and nonzero scalars on all edges. For such checks, \emph{Covers and Local Generation} shows that the
abelianization of~$\tilde M$, over the Laurent polynomial ring in the chord generators, has nonzero
determinant~\cite[Lemma~6.1]{Covers2026}, with a written proof there. A factorization $\tilde M=PQ$ through fewer than $N$
columns would specialize to a factorization of this invertible matrix, so $\irk\tilde M=N=16$, as in the proof
of~\cite[Theorem~6.2]{Covers2026}, and $r_U(M)=0$. As an independent computational check, the script
\leanfile{tools/clebsch_lift.py} constructs an explicit $16$-sheeted lift~$X_0$ of the Clebsch graph, by permutation voltages
on the $25$ edges outside a breadth-first spanning tree, and computes, by exact elimination over~$\F_2$, that the lift of
$I+A$ has nullity~$10$; by~(a), $10\ge16(16-\irk\tilde M)$, so again $\irk\tilde M=16$.
\end{proof}

\begin{proposition}[Towers that never become bipartite]\label{prop:eqtower}
\notformalized
There are towers of normal subgroups of $\pi_1(\mathrm{Clebsch})$ with trivial intersection whose covers are non-bipartite at
every level and have girth tending to infinity, and along which the rate of the lifted Clebsch code tends to~$0$.
\end{proposition}

\begin{informalproof}
The Clebsch graph has $16$ vertices and $40$ edges, so $\pi_1(\mathrm{Clebsch})$ is free of rank~$25$. The kernel
$K=\ker(F_9\to\Z/3)$ of a surjection from the free group of rank~$9$ is free of rank $1+3\cdot8=25$ by Schreier's formula. It is
dense in the pro-$2$ completion $\hat F_9$, since for every normal subgroup $N$ of $2$-power index the index of $KN$ divides both
$3$ and a power of~$2$; and $K$ embeds in $\hat F_9$, since free groups are residually~$2$. A character of~$K$ is continuous for
the induced topology exactly when it is the restriction of a character of~$F_9$, and restriction is injective on characters
because $F_9/K$ has order~$3$; so exactly $2^9$ of the $2^{25}$ characters of~$K$ are continuous. Since $\mathrm{Aut}(F_{25})$
maps onto $\mathrm{GL}_{25}(\F_2)$, it acts transitively on nonzero characters, so there is an isomorphism
$\varphi\colon\pi_1(\mathrm{Clebsch})\to K$ for which $\chi\circ\varphi^{-1}$ is not continuous. For a descending chain of open
normal subgroups $U_j\le\hat F_9$ with trivial intersection, put $\Lambda_j=\varphi^{-1}(K\cap U_j)$. These are normal with
trivial intersection; $\Lambda_j\not\le\pi^+$ for every~$j$, since otherwise $\chi\circ\varphi^{-1}$ would vanish on $K\cap U_j$
and be continuous; and the girth tends to infinity, because only finitely many conjugacy classes have bounded translation
length and each nontrivial element eventually leaves the normal subgroups~$\Lambda_j$. By \cref{thm:eqcrux}(e), with $\iota$ the
composite of $\varphi$ and the inclusion $K\subseteq\hat F_9$, and by \cref{thm:eqcrux}(f), the rate tends to $r_U=0$; the limit
also follows from~\cite[Theorem~6.2]{Covers2026}.
\end{informalproof}

\paragraph{Why it matters.} A construction cannot obtain its coding gain from a fixed small code lifted along towers that
become bipartite, nor, for binary linear codes, along towers pulled back injectively from free pro-$2$ groups: bipartite levels
have rate at most~$\frac12$ for every code (\cref{thm:eqcrux}(d)), lifting a binary check along double covers never raises the
rate (\cref{thm:eqcrux}(c)), and along towers of type~(e) the rate tends to the universal rate, at most~$\frac12$, and to~$0$
for the Clebsch code, whose rate is $\frac58$ on its own graph~\cite[Corollary~6.11]{FourthPaper2026}. The theorem leaves
open: normal
towers with trivial intersection and $C\not\le\hat\pi^+$ that are not of type~(e), such as $2$-adic analytic towers, where a
positive limit for the Clebsch code would contradict the mod-$2$ case of the Lück approximation conjecture in positive
characteristic, which Jaikin-Zapirain and Souza discuss for chains with trivial intersection~\cite[Section~8]{JaikinSouza2026};
non-normal towers; towers with short even cycles at positive density; and codes that are not lifts of a code on the base.
Trivial intersection matters: for the Clebsch adjacency check, the cyclic covers of degree $2^j$ with the voltage on a single
edge form a normal tower of binary lifts whose intersection is not trivial; odd cycles that avoid that edge lift to odd
cycles, so no level is bipartite, and the dimension $N/2+2$ at every level~\cite[Theorem~9.5]{Covers2026} makes the rate tend
to~$\frac12$, not to $r_U=0$.

\paragraph{Relation to earlier work.} \emph{Covers and Local Generation in Linear Storage Codes}~\cite{Covers2026} contains
several parts of the theorem. Its lift identity~\cite[Lemma~9.1]{Covers2026} is the recursion of~(c) for symmetric binary
checks, and its Lemma~4.1 gives the monotonicity in~(c) for towers of normal subgroups. Its Proposition~3.2 contains~(b). Its Lemma~6.1 and Theorem~6.2, full inner rank and rate tending
to~$0$ along every chain from a faithful embedding into a free-by-$\Z_p$ pro-$p$ group, for every check with nonzero edge
entries on a graph with a spanning $2$-factor, give~(f) without computation and the limit in \cref{prop:eqtower}. It also
proves, under pro-$p$ and separation hypotheses, an asymptotic half-rate bound for homogeneous checks on a fixed free-group
voltage template in every characteristic~$p$, uniformly over coefficient fields and values~\cite[Theorem~4.3]{Covers2026},
through Jaikin-Zapirain and Souza's Theorem~8.1, and the exact finite bound $N/2+2$ for the Clebsch adjacency code along
every sequence of connected binary lifts~\cite[Theorem~9.2]{Covers2026}. New here are the twist-pairing recursion of~(c) for
asymmetric binary checks along arbitrary towers of index-$2$ subgroups, part~(d), and the towers of \cref{prop:eqtower}, which
never become bipartite, have girth tending to infinity and are pulled back injectively from~$\hat F_9$.

\subsection{Scalar codes at degree four}\label{sec:crux-deg4}

Let $F$ be a field and $M\in F^{V\times V}$ a matrix with unit diagonal, $N=\card V$. The \emph{recovery set} of~$v$ is
$S_v=\{u\ne v:M_{vu}\ne0\}$, and we say that $u$ \emph{points to}~$v$ if $v\in S_u$. The matrix has \emph{no recovery triangle}
if $q\notin S_p$ whenever $p,q\in S_r$ and $p\ne q$; this holds whenever the off-diagonal support of~$M$ lies in the edges of a
triangle-free graph. Following the fourth paper's Theorem~S~\cite[Theorem~6.13]{FourthPaper2026}, consider at a vertex~$v$ the
conditions
\begin{itemize}[leftmargin=3em]
\item[$(\star\star)$] $\card{S_v}\le2$, or $\card{S_v}=3$ and $r\in S_v$ whenever $v\in S_r$;
\item[$(\star4)$] $\card{S_v}\le3$, or $\card{S_v}=4$ and $r\in S_v$ whenever $v\in S_r$.
\end{itemize}
If the support of~$M$ lies in a graph of maximum degree at most~$4$, every vertex satisfies $(\star4)$, since
$\card{S_v}=4$ forces $S_v$ to be the whole neighbourhood. A vertex violating $(\star\star)$ is \emph{heavy}, and $H$ is the set
of heavy vertices. The \emph{separation} conditions are: (S1) no vertex has two heavy vertices in its recovery set; (S2) no
heavy vertex has a heavy vertex in its recovery set. If the support lies in a graph~$G$ and the heavy vertices are pairwise at
distance at least~$3$ in~$G$, both hold.

\begin{theorem}[Theorem SH: separated heavy vertices]\label{thm:SH}
\formal{TheoremSH.lean}{theoremSH, theoremSH_graph}
Let $M$ have unit diagonal and no recovery triangle, and let every vertex satisfy $(\star4)$. If $\card H\le2$, or if the heavy
vertices satisfy (S1) and~(S2), then $\dim\ker M\le\lfloor N/2\rfloor$, over every field. In graph form: if $G$ is triangle-free
of maximum degree at most~$4$, the off-diagonal support of~$M$ lies in the edges of~$G$, and the vertices violating
$(\star\star)$ number at most two or lie pairwise at distance at least~$3$ in~$G$, then $\rank M\ge\lceil N/2\rceil$.
\end{theorem}

Theorem~S of the fourth paper is the case $H=\emptyset$. Since $\dim\ker M^\T=\dim\ker M$ and the absence of recovery
triangles is invariant under transposition (\lean{finrank_ker_transpose}, \lean{recTriangleFree_transpose}), the same holds
when $M^\T$ satisfies $(\star4)$ and its heavy vertices, defined through the sets $\{u\ne v:M_{uv}\ne0\}$, satisfy the
hypothesis instead (\lean{theoremSH_transpose}).

\begin{corollary}[Barrier and matching forms]\label{cor:SH}
\formal{TheoremSH.lean}{finrank_ker_le_barrier_SH, rho3, finrank_ker_le_matching_rho3, finrank_ker_le_matching_rho3_transpose}
Let $M$ have unit diagonal and no recovery triangle, let $U\subseteq V$, and let $V\setminus U$ be partitioned into blocks with
$M_{uv}=0$ for $u,v$ in different blocks. If on every block~$B$ every vertex satisfies $(\star4)$ on~$B$ and the vertices
violating $(\star\star)$ on~$B$ are at most two or satisfy (S1) and~(S2) on~$B$, then
$\dim\ker M\le\card U+\sum_B\lfloor\card B/2\rfloor$. If the support of~$M$ lies in a triangle-free graph~$G$ of maximum degree at
most~$4$, then
\[
  \dim\ker M\;\le\;\nu(G)+\card H-\max\bigl(\rho_3(H),\min(\card H,2)\bigr),
\]
where $\nu(G)$ is the matching number and $\rho_3(H)$ the largest size of a subset of~$H$ at pairwise distance at least~$3$
in~$G$; the same holds with the heavy set of~$M^\T$ in place of~$H$.
\end{corollary}

\begin{outline}
In Lean the unit diagonal and the absence of recovery triangles form the hypothesis \lean{MotionDepthCodes.ThmSHyp}, and
the other hypotheses of \cref{thm:SH}, relative to a vertex set~$W$, form \lean{SHOn}: $(\star4)$ on~$W$, and at most two heavy
vertices or (S1) and (S2) on~$W$, where the heavy vertices are those that violate $(\star\star)$ on~$W$ (\lean{Star4On},
\lean{heavySet}, \lean{StarOn}, \lean{S1On}, \lean{S2On}). The half bound holds for the code of~$M$ on every such~$W$
(\lean{finrank_subCode_le_half_SH}). The kernel of $M$ meets $\{x_U=0\}$ in a
subspace of codimension at most $\card U$, which lies in the code on $V\setminus U$; this code splits over the blocks
(\lean{MotionDepthCodes.finrank_subCode_biUnion_le}, for the graph of the nonzero entries), and the half bound applies to each
block (\lean{finrank_ker_le_barrier_SH}). For the matching form (\lean{finrank_ker_add_card_le_matching}), let $D$ be a set of
heavy vertices with at most two elements or at pairwise distance at least~$3$, that is, not adjacent and without a common
neighbour (\lean{Far3}), and put $U=H\setminus D$. The Tutte--Berge formula~\cite{Berge1958}, in the barrier form of
lean-machine (\lean{SimpleGraph.exists_card_add_sum_card_div_two_eq_matchingNumber}), gives a set~$U'$ with
$\card{U'}+\sum_B\lfloor\card B/2\rfloor=\nu(G)$ over the components~$B$ of $G-U'$. Apply the barrier form with the set
$U\cup U'$ and the blocks $B\setminus U$. They are separated, and on each of them the heavy vertices lie in~$D$, since a vertex
heavy on a subset is heavy (\lean{heavySet_mono}); so the hypotheses hold on each block, (S1) and (S2) following from the
distances. Hence $\dim\ker M\le\card U+\card{U'}+\sum_B\lfloor\card{B\setminus U}/2\rfloor\le\card H-\card D+\nu(G)$. A largest
set at pairwise distance at least~$3$ (\lean{rho3}, \lean{exists_rho3}) and a set of $\min(\card H,2)$ heavy vertices are such
sets~$D$, which gives the bound as stated (\lean{finrank_ker_le_matching_rho3}). For $M^\T$, whose support lies in~$G$ as
well, apply the bound to~$M^\T$ (\lean{finrank_ker_le_matching_rho3_transpose}).
\end{outline}

\paragraph{Why it matters.} The theorem extends the fourth paper's Theorem~S, the case $H=\emptyset$, to matrices with few
or scattered heavy vertices at degree four. A scalar linear code $\ker M$ beating $\lfloor N/2\rfloor$ by $\eps N$ on a
triangle-free graph of maximum degree four needs $\card H-\rho_3(H)\ge\eps N$, for~$M$ and for~$M^\T$ alike: a linear number
of heavy vertices that cannot be scattered, that is, dense clusters of vertices whose rows recover them from four neighbours
or from three non-reciprocally. For all codes, guessing numbers of triangle-free graphs need not equal the lower bound
given by fractional clique covers, as Cameron, Dang and Riis showed with the Higman--Sims graph~\cite{CameronDangRiis2016}. This is the regime of the fourth paper's tree code, in which every
interior vertex is heavy and no locally confined argument can work~\cite[Proposition~6.28]{FourthPaper2026}. For linear machines,
\cref{thm:Dlin} needs the linear crux, and the theorem says where a linear gain at degree four would have to live.

\paragraph{Proof of \cref{thm:SH}.} Write $T_v=\{v\}\cup S_v$ for the support of the row $z_v$ of~$M$. The proof works with the
codes of the fourth paper's formalization on vertex sets (\lean{MotionDepthCodes.subCode}): the code on~$W$ consists of the
vectors supported on~$W$ that satisfy the rows of~$W$, and it is the kernel of $M[W,W]$. For a vertex set~$W$, ``on~$W$'' means
that $S_v$ is replaced by $S_v\cap W$ (\lean{rs}) and only rows $r\in W$ are considered. The unit diagonal and the absence of recovery
triangles are hypotheses on all of~$M$ (\lean{MotionDepthCodes.ThmSHyp}). The class~$\mathcal H$ of vertex sets~$W$ on which
$(\star4)$ and the condition on the heavy vertices hold (\lean{SHOn}) is \emph{hereditary} (\lean{SHOn.mono}): $(\star4)$
passes to subsets, since a recovery set either keeps its size or loses elements, and the reciprocity condition only loses
in-neighbours; a vertex heavy on $W'\subseteq W$ is heavy on~$W$, for the same reason
(\lean{heavySet_mono}); so $\card H\le2$, (S1) and (S2) on~$W$ imply them on~$W'$. We prove the half bound for every
$W\in\mathcal H$ by strong induction on~$W$ (\lean{finrank_subCode_le_half_SH}). Let $W\in\mathcal H$ violate it while all its
proper subsets satisfy it (\lean{MinCE}); below, $V$ stands for~$W$, $N=\card V$, $\mathcal C$ is its code and $H$ its set of
heavy vertices. The two cases, $\card H\le2$ and the separated case, are treated together.

For $X\subseteq V$ let $c(X)$ be the dimension of the codewords vanishing on~$X$ (\lean{cdim}). For $u\in X$, the
\emph{pending set} of~$u$ at~$X$ is $P_X(u)=S_u\setminus X$ (\lean{pend}), and $u$ is \emph{closed} at~$X$ if
$P_X(u)=\emptyset$. A \emph{move} at~$X$ is a
pair $(r,P_X(r))$ with $r\in X$ and $P_X(r)\ne\emptyset$; it leads to $X\cup P_X(r)$, and $P_X(r)$ is its \emph{block}.

\begin{lemma}[Axioms of a minimal counterexample]\label{lem:SH-axioms}
\formal{SHBase.lean}{MinCE.odd, MinCE.dim, MinCE.ub, MinCE.two_le_rs, MinCE.out_card}
\begin{enumerate}[label=(\alph*)]
\item $N=2m+1$ and $\dim\mathcal C=m+1$.
\item For every nonempty $X\subseteq V$, $c(X)\le\lfloor(N-\card X)/2\rfloor$.
\item $\card{S_v}\ge2$ for every~$v$.
\item For every nonempty $X\subsetneq V$, the set $\mathrm{Out}(X)=\{v\in X:S_v\not\subseteq X\}$ has at least two elements if
$\card X$ is even and at least one if $\card X$ is odd.
\end{enumerate}
\end{lemma}

By~(a), $\rank M=m$. By duality, (b) says that $\dim(\mathcal C^\perp\cap F^X)\le\lceil\card X/2\rceil-1$, where
$\mathcal C^\perp$ is the row space; in particular the rows supported in~$X$ span at most $\lceil\card X/2\rceil-1$
dimensions.

\begin{outline}
(a) Minimality gives $\dim\ker M[Y,Y]\le\lfloor\card Y/2\rfloor$ for $Y\subsetneq V$. The codewords vanishing at a vertex~$v$
lie in the code on $V-v$, so $\dim\mathcal C\le\lfloor(N-1)/2\rfloor+1$, which exceeds $\lfloor N/2\rfloor$ only if $N=2m+1$ and
$\dim\mathcal C=m+1$ (\lean{MotionDepthCodes.start_facts}); then $c(\{v\})=m$ for every~$v$ (\lean{MinCE.single}). (b) The
codewords vanishing on~$X$ lie in the code on the proper subset $V\setminus X$
(\lean{MotionDepthCodes.subCode_inf_zeroOn_le}). (c) If $S_v=\emptyset$, the row of~$v$ forces $x_v=0$, so $c(\{v\})=m+1$. If
$S_v=\{u\}$, the codewords vanishing at~$v$ vanish at~$u$, so $c(\{v,u\})=c(\{v\})=m>\lfloor(N-2)/2\rfloor$
(\lean{cdim_move}, \lean{MinCE.single_odd_false}). Both contradict~(b). (d) Let $I=X\setminus\mathrm{Out}(X)$ and let $Q$ be
the space of the codewords vanishing on $\mathrm{Out}(X)$, of dimension at least $m+1-\card{\mathrm{Out}(X)}$. The rows of the
vertices of~$I$ are supported in~$X$, so restriction to~$I$ maps~$Q$ into the code on~$I$, of dimension at most
$\lfloor\card I/2\rfloor$, and its kernel is the space of the codewords vanishing on~$X$, of dimension at most
$\lfloor(N-\card X)/2\rfloor$ (\lean{exists_restrict}). So
$m+1-\card{\mathrm{Out}(X)}\le\lfloor\card I/2\rfloor+\lfloor(2m+1-\card X)/2\rfloor$, which gives the claim.
\end{outline}

\emph{Good states and runs.} A nonempty set~$X$ is a \emph{good state} if equality holds in~(b) (\lean{Good}). The row of
the mover makes one pending coordinate a combination of the others, so a move lowers~$c$ by at most $\card{P_X(r)}-1$
(\lean{cdim_move}). Together with \cref{lem:SH-axioms}(b) this gives the effect of a move from a good state~$X$: a block of
size~$1$ is a contradiction if $\card X$ is odd and leads to a good state if $\card X$ is even; a block of size~$2$ leads to a
good state; a block of size~$3$ leads to a good state if $\card X$ is odd; larger blocks, and blocks of size~$3$ at even
states, may leave the good states. The moves of the three kinds that lead to good states, a block of size~$2$, of size~$1$ at
an even state, or of size~$3$ at an odd state, are \emph{neutral}; they lower~$c$ by exactly $\card P-1$
(\lean{MinCE.good_move}, \lean{MinCE.single_odd_false}). A \emph{run} from~$v_0$ is a sequence of neutral moves starting at
the good state $\{v_0\}$; its first move is $(v_0,S_{v_0})$. A \emph{state} of a run is a set that it reaches after at least
one move. The states are good, contain $v_0$ and~$S_{v_0}$, and every other vertex~$u$ of a state has a \emph{parent}, a
vertex of the state that points to~$u$ (\lean{St}, \lean{St.first}, \lean{St.move}). A heavy vertex~$g$ is \emph{non-closed} at~$X$ if $g\notin X$ or
$P_X(g)\ne\emptyset$; the set $\mathrm{NC}(X)$ of these vertices only shrinks as $X$ grows (\lean{NC}, \lean{NC_anti}).

\emph{Test vectors.} At a neutral move $(r,P)$ from~$X$, the codewords vanishing on~$X$ restrict onto the whole hyperplane of
the row of~$r$ on~$P$: for distinct $u,u'\in P$, some codeword vanishes on~$X$ and on $P\setminus\{u,u'\}$ but not at~$u$, nor,
by the row of~$r$, at~$u'$ (\lean{test_vec}). Indeed, imposing the zeros on $P\setminus\{u,u'\}$ lowers~$c$ by at most
$\card P-2$, while a codeword that also vanishes at~$u$ vanishes at~$u'$ and on all of $X\cup P$, where $c$ is lower by
$\card P-1$. Test vectors and rows replace the arguments with the span of the rows of a run.

\begin{lemma}[Pending sizes, immediate closure, a Hall count]\label{lem:SH-tools}
\formal{SHBase.lean}{St.pend_start, St.pend_light, St.pend_le_three, imm_closure, MinCE.hall_count}
\begin{enumerate}[label=(\alph*)]
\item At every state~$X$ of a run from~$v_0$, $v_0$ is closed, a light vertex $u\in X$ other than~$v_0$ has
$\card{P_X(u)}\le2$, and a heavy vertex $u\in X$ other than~$v_0$ has $\card{P_X(u)}\le3$.
\item (\emph{Immediate closure}.) After a neutral move with a block~$P$ of size at least~$2$, every member~$w$ of~$P$ has
$P_{X\cup P}(w)\ne\emptyset$. Consequently a run cannot end at~$V$ with a last block of size at least~$2$.
\item If $I\subseteq V$ is nonempty and no vertex of~$I$ points to another, then $\card{I\cup\bigcup_{v\in I}S_v}\ge2\card I+1$.
\end{enumerate}
\end{lemma}

\begin{outline}
(a) If $\card{S_u}=3$ and $u$ is light, $(\star\star)$ puts its parent, which points to~$u$, into~$S_u$; if $\card{S_u}=4$,
$(\star4)$ does the same; the parent lies in~$X$. (b) Let $u'\in P$ with $u'\ne w$. If $u'\in S_w$, then $w,u'\in S_r$ form a
recovery triangle. Otherwise a test vector for $(w,u')$ vanishes on~$X$ and on $P\setminus\{w,u'\}$ but not at~$w$; if
$S_w\subseteq X\cup P$, it vanishes on all of~$S_w$, and the row of~$w$ forces it to vanish at~$w$. At~$V$ every vertex is
closed (\lean{rs_subset}). (c) The rows $z_v$, $v\in I$, are supported in $U=I\cup\bigcup_{v\in I}S_v$, and they vanish on $I\setminus\{v\}$ since no
vertex of~$I$ points to another. So the codewords vanishing on $U\setminus I$ vanish on~$U$, and by \cref{lem:SH-axioms}(a) and~(b)
$m+1\le\card{U\setminus I}+c(U)\le\card U-\card I+\lfloor(2m+1-\card U)/2\rfloor$, which gives the claim.
\end{outline}

\begin{lemma}[The odd-state lemma]\label{lem:SH-OS}
\formal{SHPhase.lean}{OSAt, odd_state}
Assume $\card H\le2$, or (S1) and~(S2). For every $t\ge0$, no state $X\ne V$ of a run, with $\card X$ odd and
$\card{\mathrm{NC}(X)}\le t$, exists.
\end{lemma}

\begin{outline}
In Lean the states are the sets with the properties of \lean{St}, which include the states of runs. By induction on~$t$; the
cases that invoke the hypothesis for $t-1$ do not arise when $t=0$. Extend the run by neutral moves only, so that all states are
good.

\emph{Phase O.} At the current odd state $X'\ne V$: (O1) a vertex of~$X'$ with exactly one pending vertex gives a
contradiction; (O2) otherwise, if a vertex~$u$ of~$X'$ has $\card{P(u)}=2$, make that move, which keeps the state odd. If it
yields~$V$, its last block is a pair, contradicting immediate closure. Otherwise, if $u$ is light, continue Phase~O at the new
state; if $u$ is heavy, it is closed at the new state, which therefore has at most $t-1$ non-closed heavy vertices, and the
induction hypothesis gives a contradiction. (O3) Otherwise every vertex of $\mathrm{Out}(X')$ is heavy with exactly three
pending vertices, by \cref{lem:SH-tools}(a) and because $v_0$ is closed. By
\cref{lem:SH-axioms}(d), $\mathrm{Out}(X')\ne\emptyset$, so $t\ge1$; go to Phase~E with $Z=X'$ and some $h\in\mathrm{Out}(Z)$.
Every move enlarges the state, so Phase~O ends.

\emph{Phase E} (\lean{F3}). Let $P=P_Z(h)=\{t_1,t_2,t_3\}$. The move $(h,P)$ is neutral at the odd good state~$Z$; $X_0=Z\cup P$ is an
even state, $h$ is closed, and $\card{\mathrm{NC}(X_0)}\le t-1$. No $t_i$ points to another, since $M$ has no recovery triangle.
If (S1) and (S2) hold, all $t_i$ are light by~(S2); if $\card H\le2$, then $\card{\mathrm{NC}(X_0)}\le1$. Below, a \emph{state} is a state reached
from~$X_0$ by neutral moves, and $Y=V\setminus X$; $\card Y$ is odd at even states.

\emph{(F1) No early singles} (\lean{F1}). At an even state with $\card Y\ge3$, no vertex of the state has exactly one pending vertex:
that move is neutral and leads to an odd state $\ne V$ with at most $t-1$ non-closed heavy vertices, contradicting the induction
hypothesis.

\emph{(F2) Pair coincidence} (\lean{F2}). At an even state~$X$ with $\card Y\ge3$, any two pending sets of size exactly~$2$ of
vertices of~$X$ are equal. We argue by induction on~$\card Y$. Let $u_1,u_2\in X$ with $q_1=P_X(u_1)\ne q_2=P_X(u_2)$ of
size~$2$. If $\card Y=3$, then
$q_1\cup q_2=Y$ and $q_1\cap q_2=\{p\}$. A test vector for the move $(u_1,q_1)$ vanishes on~$X$ but not at~$p$. The other two
vertices of~$Y$ lie with~$p$ in $S_{u_1}$ or in $S_{u_2}$, so they are not in~$S_p$, and $S_p\subseteq X$; the row of~$p$ forces
the test vector to vanish at~$p$. If $\card Y\ge5$ and $\card{q_1\cap q_2}=1$, the move $(u_1,q_1)$ leaves $u_2$ with a single
pending vertex at an even state with $\card Y-2\ge3$, contradicting~(F1). So $q_1\cap q_2=\emptyset$. Let $X_2=X\cup q_2$, after
the move of~$u_2$, and $X_{12}=X\cup q_1\cup q_2$, after the moves of $u_1$ and then~$u_2$; both have last block~$q_2$. Let
$c\in q_2$. By immediate closure at~$X_2$, $P_{X_2}(c)\ne\emptyset$, and by~(F1) at~$X_2$ it does not have size~$1$. If it has
size~$2$, the inductive hypothesis at~$X_2$, applied to $c$ and~$u_1$, whose pending set there is still~$q_1$, gives
$P_{X_2}(c)=q_1$; then $P_{X_{12}}(c)=\emptyset$, contradicting immediate closure at~$X_{12}$. So $\card{P_{X_2}(c)}\ge3$ and $c$
is heavy, by \cref{lem:SH-tools}(a). Then both members of~$q_2$ are heavy, non-closed and in~$S_{u_2}$, contradicting~(S1),
and, if $\card H\le2$, the bound $\card{\mathrm{NC}(X)}\le1$ (\lean{F2Mode}).

\emph{(F3) The final count.} By immediate closure at~$X_0$, every~$t_i$ has $P_{X_0}(t_i)\ne\emptyset$, so a heavy member
of~$P$ is non-closed at~$X_0$; and since no $t_i$ points to another, $S_{t_i}\cap X_0\subseteq Z$ for every~$i$. If $V\setminus X_0=\{w\}$, every~$t_i$ points
to~$w$. A test vector for the move $(h,P)$ and the pair $(t_2,t_3)$ vanishes on~$Z$ and at~$t_1$ but not at~$t_2$; the row
of~$t_1$, with $S_{t_1}\subseteq Z\cup\{w\}$, forces it to vanish at~$w$, and then the row of~$t_2$ forces it to vanish at~$t_2$,
a contradiction. So $\card{V\setminus X_0}\ge3$.

If all~$t_i$ are light, then by~(F1) and \cref{lem:SH-tools}(a) $\card{P_{X_0}(t_i)}=2$, and by~(F2) all three pending pairs
equal one pair~$p$. A light $t_i$ with $\card{S_{t_i}}=3$ has $h\in S_{t_i}$ by $(\star\star)$, since $h$ points to~$t_i$, so
$S_{t_i}=\{h\}\cup p$; otherwise $S_{t_i}=p$. Then $\card{P\cup\bigcup_iS_{t_i}}\le6<7$, contradicting \cref{lem:SH-tools}(c)
with $I=P$.

Otherwise, by~(S2), $\card H\le2$ and $\card{\mathrm{NC}(X_0)}\le1$; let $t_3=g$ be heavy and $t_1,t_2$ light. As above,
$t_1$ and~$t_2$ have the same pending pair $p_1=\{a,b\}$. If $\card{V\setminus X_0}\ge5$, make the move $(t_1,p_1)$, reaching an
even state~$X_1$ with $\card{V\setminus X_1}\ge3$. The vertices $a,b$ are light, since $g$ is the only non-closed heavy vertex.
By immediate closure, (F1) at~$X_1$ and \cref{lem:SH-tools}(a) each of them has a pending pair, the same pair~$p_2$
by~(F2). Both $t_1$ and~$t_2$ point
to~$a$, so if $\card{S_a}=3$, $(\star\star)$ would give $S_a\supseteq\{t_1,t_2\}\cup p_2$, four elements; hence $S_a=p_2$, and
likewise $S_b=p_2$, and \cref{lem:SH-tools}(c) with $I=\{a,b\}$ is contradicted, as $\card{\{a,b\}\cup p_2}=4<5$. If
$\card{V\setminus X_0}=3$, then $c(X_0)=1$, and the state $X_1=X_0\cup p_1$ after the move of~$t_1$ has $c(X_1)=0$. A codeword
that vanishes on~$X_0$ and at~$a$ vanishes at~$b$, by the row of~$t_1$, so some codeword~$y$ vanishes on~$X_0$ with $y_a\ne0$. The
rows of $t_1$ and~$t_2$ give $M_{t_1a}y_a+M_{t_1b}y_b=0=M_{t_2a}y_a+M_{t_2b}y_b$. A test vector~$x$ for the move $(h,P)$ and the
pair $(t_2,g)$ vanishes on~$Z$ and at~$t_1$ but not at~$t_2$, and the same rows give $M_{t_1a}x_a+M_{t_1b}x_b=0$ and
$x_{t_2}+M_{t_2a}x_a+M_{t_2b}x_b=0$. The vectors $(x_a,x_b)$ and $(y_a,y_b)$ are annihilated by the nonzero row
$(M_{t_1a},M_{t_1b})$, so they are proportional; hence $M_{t_2a}x_a+M_{t_2b}x_b=0$ and $x_{t_2}=0$ (\lean{two_rows_zero}), a
contradiction.
\end{outline}

Since $\mathrm{NC}(X)\subseteq H$, \cref{lem:SH-OS} with $t=\card H$ says that in a minimal counterexample no run with at least
one move reaches an odd state other than~$V$.

\begin{outline}[Outline of the formal proof of \cref{thm:SH}]
\emph{The start.} The three cases below make up \lean{MinCE.false_SH}. By \cref{lem:SH-axioms}(c), $\card{S_v}\ge2$ for
all~$v$; light vertices have $\card{S_v}\le3$ and heavy ones $\card{S_v}\in\{3,4\}$.
\emph{Case 1: some $v_0$ has $\card{S_{v_0}}=2$.} The move $(v_0,S_{v_0})$ is neutral and gives the odd state $T_{v_0}$. If
$T_{v_0}=V$, immediate closure gives a contradiction; otherwise \cref{lem:SH-OS} does.
\emph{Case 2: no $\card{S_v}=2$, and some heavy~$h$ has $\card{S_h}=3$.} The move $(h,S_h)$ is neutral at $\card X=1$ and gives
the even state $X_0=T_h$ of a run from~$h$, with $h$ closed. This is Phase~E with $Z=\{h\}$ and $P=S_h$: the $t_i$ are light
by~(S2), or $\card{\mathrm{NC}(X_0)}\le\card H-1\le1$; (F1) holds by \cref{lem:SH-OS} with $t=\card{\mathrm{NC}(X_0)}$, and
(F2) and~(F3) give a contradiction.
\emph{Case 3: every light~$v$ has $\card{S_v}=3$ and every heavy~$v$ has $\card{S_v}=4$.} By $(\star\star)$ and $(\star4)$ every
in-neighbour of~$v$ lies in~$S_v$, for every~$v$, so the relation is symmetric. Its graph~$\Gamma$ is triangle-free, with light
vertices of degree~$3$ and heavy ones of degree~$4$. The vertices of odd degree are even in number and $N$ is odd, so $\card H$
is odd and $H\ne\emptyset$; if $\card H\le2$, then $\card H=1$, where (S1) and (S2) hold. Pick $h\in H$ and $v_0\in S_h$. Then
$v_0$ is light by~(S2), $S_{v_0}=\{h,s_1,s_2\}$, and $s_1,s_2$ are light by~(S1), since $v_0$ points to~$h$. The move
$(v_0,S_{v_0})$ is neutral at $\card X=1$, and $X_0=\{v_0,h,s_1,s_2\}$ is an even state. (F1) and~(F2) hold at the even states
reached from~$X_0$, with the same proofs, using \cref{lem:SH-OS} with $t=\card{\mathrm{NC}(X_0)}$. The vertex~$h$, its four neighbours
and $s_1,s_2$ are seven distinct vertices, since $\Gamma$ is triangle-free. If $N=7$, each neighbour of~$h$ has its two further
neighbours among $s_1,s_2$, so $s_1$ has the four neighbours of~$h$ as neighbours, although it is light. So $N\ge9$. Since
$\Gamma$ is triangle-free, $S_{s_i}\cap X_0=\{v_0\}$, so $P_{X_0}(s_i)$ has size~$2$, and by~(F2) both pending sets equal one
pair~$p_1$, as $\card{V\setminus X_0}\ge5$. Make the move $(s_1,p_1)$; $X_1=X_0\cup p_1$ is an even state with
$\card{V\setminus X_1}\ge3$. By~(S1) some $a\in p_1$ is light. As $a\in S_{s_1}\cap S_{s_2}$, symmetry gives $S_a=\{s_1,s_2,x\}$,
so $P_{X_1}(a)\subseteq\{x\}$; by immediate closure it is nonempty, and $a$ has a single pending vertex, contradicting~(F1).

All cases are contradictory, which proves the half bound on every set of~$\mathcal H$ (\lean{finrank_subCode_le_half_SH}); the
case $W=V$ is \cref{thm:SH} (\lean{theoremSH}). For the graph form (\lean{theoremSH_graph}), a triangle-free support gives no
recovery triangle (\lean{MotionDepthCodes.thmSHyp_of_cliqueFree}), maximum degree at most~$4$ gives $(\star4)$
(\lean{star4_of_degree_le_four}), heavy vertices at pairwise distance at least~$3$ satisfy (S1) and~(S2)
(\lean{sep_of_far3}), and the rank is $N$ minus the nullity.
\end{outline}

\begin{remark}[Where the argument stops]\label{rem:SH-stop}
Within the odd-state lemma and Case~2 of the start, (S1) is used only at the end of~(F2) and (S2) only to make the members of
a triple block light, and there, without them, the argument fails exactly at two configurations: two vertices whose pending
pairs are disjoint and consist of four heavy vertices, and a triple block of a heavy vertex that contains a heavy vertex. Case~3
uses them also to make $v_0$, $s_1$, $s_2$ and a member of~$p_1$ light. By \cref{thm:SH}, a minimal counterexample to the
general degree-four question has heavy vertices at distance at most two, and the contradiction tree of the argument must break
at a cluster of heavy vertices.
\end{remark}

\begin{remark}[Numerical: small cases]\label{rem:SH-small}
Exact searches over~$\F_2$ for all digraphs satisfying the hypotheses of \cref{thm:SH} with at most nine vertices, with a
positive control that allows recovery triangles, found no counterexample; such small cases cannot test the asymptotic
question.
\end{remark}

\section{What a linear-time transposer would need}\label{sec:construction}

Our aim is to decide~(T), and we would prefer a construction. This section collects what the results above require of a
linear-time transposer, oblivious unless said otherwise. It contains no new theorems: every claim is one of the results above
or of earlier papers of the series, or is labelled as a heuristic or as resting on a proof sketch.

\paragraph{The interface.} For an oblivious transposer~$M$ with $k$ work tapes and running time $C(L+1)$, and every
$\eps\in(0,1]$ and $D\ge0$, Theorem~D says: either the dependency DAG has no genie of cost at most $\eps L$ and full depth at
most~$D$, or the revealed word of such a genie is an index code with at most $\eps L$ bits for the machine's own cone side
information, and some fibre is a binary storage code of rate at least $1-\eps$ on the cone relation
\cite[Theorem~4.10 and Remark~4.9(b)]{FourthPaper2026}. After a paid trimming of a small fixed fraction of coordinates, that
relation is paired-distance separated~\cite[Theorem~3.3 and Corollary~3.8]{DirectedRecovery2026}. So a transposer must either
defeat cheap depth reduction or carry, on a paired-distance separated relation of bounded degree, a code of rate at
least~$1-\eps$, a counterexample to Question~9.1 of~\cite{DirectedRecovery2026} far beyond rate~$\frac12$.

\paragraph{Requirements on head motion.} In addition to the earlier requirements, $\Omega(\log d)$ passes
\cite[Theorems~5.9 and~5.12]{FirstPaper2026}, curved macroscopic motion~\cite[Theorems~3.10 and~3.12]{FourthPaper2026}, and
heads that are not all home-based unless the crux fails~\cite[Corollary~4.23]{FourthPaper2026}, an oblivious transposer must
satisfy the following; (R1) assumes the crux at the depths of Corollary~F and of \cref{cor:M-machine}.
\begin{enumerate}[label=(R\arabic*),leftmargin=3em]
\item \emph{Unbounded nesting.} Some work heads must have unboundedly many simultaneously open nested structures
with distinct turning points, hence cells entered unboundedly often (\cref{lem:owner}, Corollary~F); either two such heads, or
one such head together with heads whose families occupy unboundedly many separated scales (\cref{cor:M-machine}). Unboundedly
many reversals alone do not suffice.
\item \emph{A deep part of vanishing capacity.} Whatever depth-robustness the schedule has lives on the cells
entered more than $\lceil1/\eps\rceil$ times, fewer than $k\eps T$ of them; the rest of the walk DAG is cut to depth
$2^{O(k^3/\eps^4)}$ by $\eps T$ edges (\cref{cor:capacity}). The deep part may relay information between shallow stretches,
and must be allowed to (\cref{prop:relays}). On the deep-nesting family, under Conjecture~A$'$ (\cref{cor:wrap}), its size must be
of order $T/\log D$ when the rest is cut at depth~$D$ (\cref{prop:dnp-capacity}), and depth-robust schedules carry a mean
time-cut capacity that grows with the depth they force (\cref{rem:abp}, proof sketch only).
\item \emph{Not a turn counter.} Families of walks without pauses with $(\mathcal N/T)\cdot H$ bounded, where $\mathcal N$ is the
number of turns and $H$ the maximal number of hops spanning a time, are depth-reducible (\cref{prop:counter-criterion}); so the
deep part must combine many turns with many hops open at once, as the deep-nesting pairs do.
\end{enumerate}

\paragraph{Requirements on the codes.} The fibre code on the cone relation must avoid every class on which rate~$\frac12$ is
proved, on all but a small fraction of its coordinates: triangle-free side information of maximum degree three and the further
directed classes of the companion paper~\cite{SeventhPaper2026}, symmetric side information of circular chromatic number below
three~\cite[Corollary~3.2]{CircularMemo2026}, hyperfinite incidence graphs and strongly connected components with at most two
internal inputs~\cite[Corollary~4.4 and Theorem~5.1]{DirectedRecovery2026}, and, for linear codes, bounded complex
coefficients~\cite[Theorem~6.17]{FourthPaper2026}. This paper adds the following for linear codes; parts of (R4) were known
before~\cite{Covers2026}.
\begin{enumerate}[label=(R\arabic*),leftmargin=3em,resume]
\item \emph{No fixed binary lift along a controlled tower.} Along a tower of covers that becomes bipartite, every code has
rate at most~$\frac12$ from that level on~\cite[Corollary~6.9]{FourthPaper2026}. A fixed small binary linear code lifted along
a tower of double covers never gains rate, and along a tower pulled back injectively from a chain of open normal subgroups
with trivial intersection of a free-by-$\Z_2$ pro-$2$ group its rate tends to its universal rate, at most~$\frac12$; for the
Clebsch code the limit is~$0$ (\cref{thm:eqcrux,prop:eqtower}; compare~\cite[Lemma~9.1 and Theorem~6.2]{Covers2026}).
\item \emph{Clustered heavy vertices at degree four.} A scalar linear code $\ker M$ beating $\lfloor N/2\rfloor$ by $\eps N$
on triangle-free side information of maximum degree four needs a matrix~$M$, and likewise~$M^\T$, with a linear number of heavy
vertices that cannot be scattered (\cref{cor:SH}).
\item \emph{Non-rigidity, for linear machines.} The transposition matrix must be the sum of a matrix of rank at most
$\eps L$ and a matrix supported on the depth-$D$ ball pattern of the schedule, for a genie under which every output is linearly
$D$-shallow (\cref{lem:rigidity}); otherwise the machine has no such genie.
\end{enumerate}

\paragraph{Coding gaps.} By the first paper, every linear-time transposer yields a causal code for complete exchange on the
block network of its run that beats multicommodity flow by a factor $\Omega(\log d)$, on all inputs if the machine is oblivious
\cite[Theorems~8.8, 9.2 and~9.9]{FirstPaper2026}, Lean-verified there; for oblivious algorithms the connection goes back to
Adler, Harvey, Jain, Kleinberg and Lehman~\cite[Theorem~6.2]{AHJKL2006}. Coding gains over routing in undirected networks
exist, and the known gains grow with the size of the network, as $(\log n)^{\eps}$ in~\cite{BravermanHe2026} and
polylogarithmically in~\cite{ZhangLiLi2026}.
Requirements (R1) and~(R2) suggest where a gain of order $\log d$ would have to be produced: in a deep part of small
capacity, visited over and over, while the shallow part only moves data. This is a heuristic. A conditional theorem of Brende and
GPT-6 Astra shows that the shape is not empty for computations in general: if $\TCP(k)$ fails for some~$k$, one fixed native
bijection with a linear-time native inverse and logarithmic work span needs genie entropy above $\delta L$ at every fixed
depth~\cite[Theorem~H]{InfoDepthMemo2026}.

\paragraph{Mechanisms that fail.} The following cannot be the whole of an oblivious transposer for which the crux holds at
the depths of Corollaries F and~B and of \cref{cor:M-machine}: home-based work heads
(Corollary~B); work heads that enter every cell boundedly often, in any coupling and at any scales (Corollary~F); one arbitrary
head together with bounded heads at boundedly many scales (\cref{cor:M-machine}). And the following cannot supply its codes on
more than a small fraction of coordinates: the classes listed above, fixed binary lifts along controlled towers for binary linear codes,
and, for scalar linear codes at degree four, scattered heavy vertices.

\section{Open problems}\label{sec:open}

For oblivious machines the series reduces $\neg$(T) to depth reduction and a coding bound (\cref{thm:D4}). This paper settles
depth reduction for bounded owner-load and for one arbitrary head plus bounded heads at boundedly many scales, and reduces the
coding side, for bounded-entry machines, to Question~9.1 of \emph{Directed Recovery Structures}. The first item below is the
question itself. For the others, by \cref{thm:D4} and the proof of Corollary~F$'$, $\TCP(k)$ for every~$k$ together with an
affirmative answer to \cref{prob:q91} would exclude oblivious linear-time transposers, while a depth-robust family of walks
(\cref{prob:tcp-cases,prob:sr3}) or a negative answer to \cref{prob:q91} would show where a construction has to work.

\begin{enumerate}[label=\arabic*.,leftmargin=2em]
\item \emph{Decide (T).} Construct a machine that transposes $d\times d$ matrices with $b(d)$-bit entries in time $O(L)$, or prove
that none exists. The first paper's frozen checker accepts either answer~\cite{FirstPaper2026}.

\item \emph{Depth reduction at unbounded owner-load.} By \cref{cor:where}, a depth-robust family of walks has two heads of
unbounded owner-load, or one such head together with bounded heads whose families occupy unboundedly many separated scales.
\begin{problem}[Two remaining cases of TCP]\label{prob:tcp-cases}
(a) Are the deep-nesting pairs depth-robust as $m\to\infty$, for some sequence of odd shifts? Conjecture~A$''$ (\cref{conj:App})
would give, by \cref{cor:wrap}, depth $2^{\Omega(m)}$ at fixed cost and a negative answer to $\TCP(2)$; the counter gives
depth $2^{m+O(1)}/\eps$ at every shift (\cref{cor:counter-cut}). (b) Is TCP true for one arbitrary walk together with walks of
owner-load one whose families occupy unboundedly many separated scales, for instance the pair of \cref{rem:umix-candidate}?
\end{problem}
Conjecture~A$''$ for block pairs would follow from visibility of the coarse classes in a window modulo~$q$ at a fixed shift
(\cref{prop:Vfixed}). Theorem~V gives less, even on average over shifts: the integer event $F\le-Q-1$, with no lower bound
on~$F$, which implies the window only for $q$ of order $(m-\kappa)Q$, with $\kappa$ the canceller level of \cref{lem:straddle}. Two points are open: a bound on $\abs F$ on the
straddling hops, or a version of Theorem~V that counts the window directly; and, at one fixed shift, equidistribution of the apex
centres of one head relative to the canceller valleys of the other. A proof for all labels must use cost, since fine levels are blind (\cref{cor:blind}), and
packing certificates cannot do it~\cite[Proposition~5.14]{FourthPaper2026}.

\item \emph{Three stacks.} $\mathrm{SR}(2)$ holds by planarity (\cref{prop:sr2}), and $\mathrm{SR}(2k)$ implies $\TCP(k)$
(\cref{prop:stacks}).
\begin{problem}[Three-stack DAGs]\label{prob:sr3}
Is every family of $3$-stack DAGs with a spine depth-reducible, that is, does $\mathrm{SR}(3)$ hold? A depth-robust family of
$3$-stack or $4$-stack DAGs, each realized as the walk DAG of two walks, would refute $\TCP(2)$.
\end{problem}

\item \emph{The coding bound.} By Corollaries F$'$ and~F$''$, an affirmative answer to the following question excludes oblivious
bounded-entry transposers, and adaptive ones whose traces have entropy at most $\eta L$ for a fixed $\eta<\frac12$ and whose
traces of probability at least $2^{-\eta L}$ have bounded entry.
\begin{problem}[Question~9.1 of~\cite{DirectedRecovery2026}]\label{prob:q91}
Do all storage codes, over every alphabet, on recovery relations of bounded incidence degree that satisfy the paired-distance
condition have rate at most $\frac12+o(1)$? By \cref{rem:weak-crux}, Corollary~F$'$ already follows from a deficit $\rho(\Delta)$
with $\rho(\Delta)\log\log\Delta\to\infty$.
\end{problem}
The companion paper proves exact half bounds on several classes of relations without short paired paths and states the
directed degree-three case as a conjecture~\cite{SeventhPaper2026}; \emph{Directed Recovery Structures} proves the hyperfinite
case and several others~\cite{DirectedRecovery2026}.

\item \emph{Degree four.}
\begin{problem}[Triangle-free side information of degree four]\label{prob:deg4}
Do triangle-free graphs of maximum degree~$4$ have guessing number $N/2+o(N)$ along families with growing girth, so that
storage codes there have rate at most $\frac12+o(1)$? For scalar linear codes, does $\dim\ker M\le N/2+o(N)$ hold there when
the heavy vertices of \cref{thm:SH} form clusters?
\end{problem}
Without the girth condition both answers are negative: the companion paper gives a $4$-regular triangle-free graph on thirteen
vertices, of odd girth five, with a scalar linear code of dimension~$7$ over~$\F_4$~\cite[Example~7.8]{SeventhPaper2026}, so
that $2h(V)\le\sum_vh(\{v\})$ fails there for a valid polymatroid, and disjoint copies of it have rate~$\frac7{13}$. The heavy
vertices of its recovery matrix form a cluster (\cref{sec:crux-deg4}). At degree three both questions are settled without a
girth condition, for all codes by the companion paper~\cite{SeventhPaper2026} and for scalar linear codes by the fourth paper's
Theorem~S. At degree four, \cref{thm:SH} reduces the scalar question to clusters of heavy vertices, the regime in which
tree-local arguments are impossible~\cite[Proposition~6.28]{FourthPaper2026}.

\item \emph{The scratch.} If TCP fails at unbounded owner-load, the route of Theorem~D needs input specific to transposition:
fixed-depth genies, even charged by entropy, then fail for some easy oblivious bijections~\cite[Theorem~H]{InfoDepthMemo2026}.
\begin{problem}[The scratch question]\label{prob:scratch}
Can an oblivious linear-time transposer use a set of fewer than $\eps L$ work-tape cells, each visited unboundedly often, to supply
the global information that the crux denies to the shallow part of its schedule (\cref{cor:capacity})?
\end{problem}

\item \emph{Rigidity.} For linear machines, linear depth-TCP together with a linear crux excludes oblivious linear-time
transposers (\cref{thm:Dlin}), and the linear crux can be stated exactly as the following question.
\begin{problem}[Relative rigidity of transposition]\label{prob:rigidity}
Is there $\rho>0$ such that, for every linear-time schedule and every fixed~$D$, the transposition matrix is not the sum of a matrix
of rank below $\rho L$ and a matrix supported on the depth-$D$ ball pattern of the schedule, for all large~$d$?
\end{problem}

\item \emph{Towers.} Do binary codes lifted along normal $2$-adic analytic towers with trivial intersection that stay
non-bipartite have asymptotic rate at most~$\frac12$? Without trivial intersection the answer is negative: two Clebsch graphs
joined by two edges, the code $\ker((I+A)\oplus(I+A))$, which ignores the joining edges, and a voltage on one joining edge
give connected non-bipartite covers of rate~$\frac58$ at every level. For the Clebsch adjacency code, connected binary-lift
towers have dimension at most $N/2+2$~\cite[Theorem~9.2]{Covers2026}, and along a normal tower with trivial intersection a
positive limit would contradict the mod-$2$ case of the Lück approximation conjecture in positive
characteristic~\cite[Section~8]{JaikinSouza2026} (\cref{sec:crux-towers}).
\end{enumerate}


\appendix
\section{Index of Lean declarations}\label{app:index}

Every declaration of this paper's library cited in the text, grouped by the Lean module that contains it, relative to the namespace
\texttt{BoundedNesting}. All of them depend only on the axioms \texttt{propext}, \texttt{Classical.choice} and
\texttt{Quot.sound} (\cref{app:verify}).
{\raggedright The declarations of other libraries named in the text are checked in those libraries: \lean{Transpose.LinearTimeTransposable} in the first paper's library~\cite{FirstPaper2026}; \lean{MotionDepthCodes.ThmSHyp}, \lean{MotionDepthCodes.card_descents_eq}, \lean{MotionDepthCodes.finrank_subCode_biUnion_le}, \lean{MotionDepthCodes.lemma_Zq}, \lean{MotionDepthCodes.start_facts}, \lean{MotionDepthCodes.subCode}, \lean{MotionDepthCodes.subCode_inf_zeroOn_le}, \lean{MotionDepthCodes.thmSHyp_of_cliqueFree} in the fourth paper's library~\cite{FourthPaper2026}; and \lean{SimpleGraph.exists_card_add_sum_card_div_two_eq_matchingNumber} in the shared library \texttt{lean-machine}.\par}

\paragraph{\texttt{BoundedNesting.Clock}}
{\small\raggedright \texttt{Fam}, \texttt{Fam.Valid}, \texttt{Fam.Valid.band\_at\_end}, \texttt{Fam.Valid.band\_of\_mem}, \texttt{Fam.Valid.overlap}, \texttt{Fam.Valid.sum\_d\_all}, \texttt{Fam.Valid.sum\_d\_le}, \texttt{Fam.Valid.unique\_span}, \texttt{Fam.band}, \texttt{Fam.card\_wrap\_le}, \texttt{Fam.clock}, \texttt{Fam.clock\_le\_of\_no\_wrap}, \texttt{Fam.clock\_lt\_of\_no\_wrap}, \texttt{Fam.cnt}, \texttt{Fam.depth\_le\_of\_lex}, \texttt{Fam.digit}, \texttt{Fam.group}, \texttt{Fam.occ}, \texttt{Fam.rk}, \texttt{flr}, \texttt{grp}, \texttt{lam}.\par}

\paragraph{\texttt{BoundedNesting.ClockCost}}
{\small\raggedright \texttt{Fam.Valid.card\_gchange\_le}, \texttt{Fam.Valid.card\_higher\_le}, \texttt{Fam.Valid.card\_incons\_le}, \texttt{Fam.Valid.card\_viol\_calm\_le}, \texttt{Fam.Valid.cost\_sum\_le}, \texttt{Fam.Valid.exists\_charge}, \texttt{Fam.Valid.exists\_cheap}, \texttt{Fam.Valid.sum\_calm\_growth\_le}, \texttt{Fam.Valid.sum\_gjumps}, \texttt{Fam.Valid.sum\_higher\_le}, \texttt{Fam.Valid.sum\_jumps\_lam}, \texttt{Fam.Valid.sum\_jumps\_pow}, \texttt{Fam.Valid.sum\_old\_le}, \texttt{Fam.Valid.sum\_viol\_calm\_le}, \texttt{Fam.calm}, \texttt{Fam.card\_desc\_still\_le}, \texttt{Fam.card\_wrap\_union}, \texttt{Fam.descents\_subset}, \texttt{Fam.digit\_lt}, \texttt{Fam.gchange}, \texttt{Fam.gjumps}, \texttt{Fam.higher}, \texttt{Fam.incons}, \texttt{Fam.still}, \texttt{Fam.sum\_desc\_still\_le}, \texttt{Fam.violated\_subset}, \texttt{card\_grp\_ne\_le}, \texttt{exists\_change}, \texttt{sum\_floor\_pow\_le}.\par}

\paragraph{\texttt{BoundedNesting.CorM}}
{\small\raggedright \texttt{BandsCovered}, \texttt{DM\_mono}, \texttt{DM\_one\_le}, \texttt{DepthRobust}, \texttt{FamTCP}, \texttt{FamTCPAt}, \texttt{corM1}, \texttt{corM1\_i}, \texttt{corM1\_ii}, \texttt{depthRobust\_iff}, \texttt{hopPairs\_allHops\_succAbove}, \texttt{lam\_mul\_le}, \texttt{mw\_le}, \texttt{others}, \texttt{owners}, \texttt{tcp\_iff\_famTCP}, \texttt{two\_pow\_lam\_sub\_le}, \texttt{where\_ownerLoad}, \texttt{where\_ratio}, \texttt{where\_windows}.\par}

\paragraph{\texttt{BoundedNesting.Families}}
{\small\raggedright \texttt{card\_fam\_le}, \texttt{disj\_aux}, \texttt{fam}, \texttt{famData}, \texttt{famData\_disj}, \texttt{famData\_valid}, \texttt{famDepth}, \texttt{famDepth\_le}, \texttt{famDepth\_lt}, \texttt{famP}, \texttt{famP\_ne}, \texttt{famU}, \texttt{fam\_cell\_injOn}, \texttt{fam\_cell\_mem}, \texttt{fam\_long}, \texttt{fam\_nested}, \texttt{fam\_no\_cross}, \texttt{fam\_span}, \texttt{famd}, \texttt{side}, \texttt{sw}, \texttt{sw\_above}, \texttt{vcode}, \texttt{vcode\_lt}, \texttt{vh\_eq}.\par}

\paragraph{\texttt{BoundedNesting.Labels}}
{\small\raggedright \texttt{card\_marked\_add\_le}, \texttt{depth\_le\_of\_label}.\par}

\paragraph{\texttt{BoundedNesting.MAverage}}
{\small\raggedright \texttt{card\_incT\_le}, \texttt{card\_rho\_lineAt}, \texttt{costB}, \texttt{sum\_cut\_le}, \texttt{sum\_offsets\_cut\_le}, \texttt{sum\_rho\_lineT}.\par}

\paragraph{\texttt{BoundedNesting.MCost}}
{\small\raggedright \texttt{ahop\_lt}, \texttt{chord\_lt}, \texttt{depth\_le\_of\_mixedRadix}, \texttt{mdig}, \texttt{mlab}, \texttt{mradix}, \texttt{nev}, \texttt{pos0}, \texttt{posD}, \texttt{posE}, \texttt{prod\_mradix}, \texttt{state\_le}.\par}

\paragraph{\texttt{BoundedNesting.MLevels}}
{\small\raggedright \texttt{LineAt}, \texttt{LineAt.mono}, \texttt{card\_hops\_span\_strip\_le}, \texttt{chordLevel}, \texttt{entries'}, \texttt{lcnt}, \texttt{lcnt\_chord\_eq}, \texttt{lcnt\_chord\_lt}, \texttt{lcnt\_state\_zero}, \texttt{stripAt}, \texttt{stripAt\_const}, \texttt{visits}.\par}

\paragraph{\texttt{BoundedNesting.MSum}}
{\small\raggedright \texttt{Zt}, \texttt{ahop\_subset}, \texttt{card\_ZH\_le}, \texttt{card\_ZS\_le}, \texttt{card\_cut\_le}, \texttt{card\_lineVisit\_le}, \texttt{card\_offsets\_wrap0}, \texttt{card\_offsets\_wrapD}, \texttt{card\_offsets\_wrapE}, \texttt{chord\_confined}, \texttt{chord\_subset}, \texttt{cutHM}, \texttt{cutSM}, \texttt{desc\_subset}, \texttt{offsets}, \texttt{sum\_ahop\_events\_le}, \texttt{sum\_ahop\_le}, \texttt{sum\_chord\_events\_le}, \texttt{sum\_chord\_growth\_le}, \texttt{sum\_chord\_le}, \texttt{sum\_desc\_le}, \texttt{sum\_pow\_beta\_lt}.\par}

\paragraph{\texttt{BoundedNesting.OneWalk}}
{\small\raggedright \texttt{IsHop}, \texttt{IsLine}, \texttt{card\_cutHops\_le}, \texttt{card\_cutStates\_le}, \texttt{cutHops}, \texttt{cutStates}, \texttt{exists\_line\_residue}, \texttt{exists\_offset}, \texttt{hop\_end\_unique}, \texttt{hop\_start\_unique}, \texttt{hop\_unique}, \texttt{hops}, \texttt{strip}, \texttt{stripCount}, \texttt{stripLabel}, \texttt{strip\_cut}, \texttt{strip\_cut\_depth}, \texttt{strip\_eq\_of\_adj}.\par}

\paragraph{\texttt{BoundedNesting.Owner}}
{\small\raggedright \texttt{IsOwner}, \texttt{IsOwner.unique}, \texttt{IsUpperHop}, \texttt{IsWalk}, \texttt{card\_crossings\_le}, \texttt{crossings}, \texttt{entries}, \texttt{owner\_entries}.\par}

\paragraph{\texttt{BoundedNesting.OwnerLoad}}
{\small\raggedright \texttt{IsLowerHop}, \texttt{IsLowerOwner}, \texttt{capacity\_cut}, \texttt{card\_heavyPairs\_lt}, \texttt{card\_heavyPairs\_mul\_le}, \texttt{card\_span\_heavyPairs\_le}, \texttt{entries\_le\_of\_lightHops}, \texttt{entries\_neg}, \texttt{entries\_sw}, \texttt{heavyPairs}, \texttt{isOwner\_hopOwner}, \texttt{isUpperHop\_sw}, \texttt{lightHops}, \texttt{ownerLoad\_le\_of\_entries}, \texttt{owner\_entries\_lower}, \texttt{sum\_card\_entries\_le}, \texttt{tcp\_of\_entries}.\par}

\paragraph{\texttt{BoundedNesting.Relays}}
{\small\raggedright \texttt{IsDagPath}, \texttt{IsDagPath.lt}, \texttt{IsDagPath.mono}, \texttt{NoPauses}, \texttt{RF}, \texttt{TCPgenie}, \texttt{card\_firstVisits\_le}, \texttt{card\_zig\_visits\_le}, \texttt{exists\_prev\_hop}, \texttt{hopPairs\_subset\_withZig}, \texttt{tcpGenie\_of\_rf}, \texttt{withZig}, \texttt{zig}, \texttt{zig\_ne}.\par}

\paragraph{\texttt{BoundedNesting.Rigidity}}
{\small\raggedright \texttt{LinNet}, \texttt{LinNet.C}, \texttt{LinNet.C\_eq}, \texttt{LinNet.GFree}, \texttt{LinNet.Pmat}, \texttt{LinNet.Tm}, \texttt{LinNet.dlin\_step}, \texttt{LinNet.minrank}, \texttt{LinNet.minrank\_le}, \texttt{LinNet.reach}, \texttt{LinNet.reach\_subset}, \texttt{LinNet.rigidity}, \texttt{LinNet.two\_pow\_finrank\_le}, \texttt{LinNet.val}, \texttt{LinNet.val\_eq\_add}, \texttt{LinNet.vlin\_expand}.\par}

\paragraph{\texttt{BoundedNesting.SHBase}}
{\small\raggedright \texttt{Good}, \texttt{MinCE}, \texttt{MinCE.dim}, \texttt{MinCE.good\_move}, \texttt{MinCE.hall\_count}, \texttt{MinCE.odd}, \texttt{MinCE.out\_card}, \texttt{MinCE.single}, \texttt{MinCE.single\_odd\_false}, \texttt{MinCE.two\_le\_rs}, \texttt{MinCE.ub}, \texttt{NC}, \texttt{NC\_anti}, \texttt{S1On}, \texttt{S2On}, \texttt{SHOn}, \texttt{SHOn.mono}, \texttt{St}, \texttt{St.first}, \texttt{St.move}, \texttt{St.pend\_le\_three}, \texttt{St.pend\_light}, \texttt{St.pend\_start}, \texttt{Star4On}, \texttt{StarOn}, \texttt{cdim}, \texttt{cdim\_move}, \texttt{exists\_restrict}, \texttt{heavySet}, \texttt{heavySet\_mono}, \texttt{imm\_closure}, \texttt{pend}, \texttt{rs}, \texttt{rs\_subset}, \texttt{test\_vec}.\par}

\paragraph{\texttt{BoundedNesting.SHPhase}}
{\small\raggedright \texttt{F1}, \texttt{F2}, \texttt{F2Mode}, \texttt{F3}, \texttt{OSAt}, \texttt{odd\_state}, \texttt{two\_rows\_zero}.\par}

\paragraph{\texttt{BoundedNesting.Stacks}}
{\small\raggedright \texttt{AvoidPath}, \texttt{SR}, \texttt{StackDag}, \texttt{StackDag.card\_pullA\_le}, \texttt{StackDag.card\_pullS\_le}, \texttt{StackDag.dep}, \texttt{StackDag.dep\_eq\_of\_arc}, \texttt{StackDag.dep\_lt\_of\_mem}, \texttt{StackDag.dep\_succ}, \texttt{StackDag.lift\_path}, \texttt{StackDag.pullA}, \texttt{StackDag.pullS}, \texttt{StackDag.pullback\_cut}, \texttt{StackDag.walk}, \texttt{StackDag.walk\_hop}, \texttt{StackDag.walk\_isWalk}, \texttt{StackDag.walk\_noPause}, \texttt{StackDag.walk\_upperHop}, \texttt{TCP}, \texttt{TCPat}, \texttt{TCPnp}, \texttt{arcCount}, \texttt{srAt\_of\_tcpnpAt}, \texttt{sr\_of\_tcpnp}, \texttt{tcpAt\_of\_srAt}, \texttt{tcp\_of\_sr}, \texttt{tcp\_of\_tcpnp\_two\_mul}, \texttt{tcpnp\_of\_tcp}, \texttt{upper\_no\_cross}, \texttt{walkDag\_isStack}.\par}

\paragraph{\texttt{BoundedNesting.TheoremF}}
{\small\raggedright \texttt{DF}, \texttt{OwnerLoadLE}, \texttt{ceil\_le\_two\_pow\_clog}, \texttt{cost\_le\_eps}, \texttt{hopOwner}, \texttt{hopOwner\_spec}, \texttt{hopPairs}, \texttt{theoremF}, \texttt{theoremF\_owner}.\par}

\paragraph{\texttt{BoundedNesting.TheoremM}}
{\small\raggedright \texttt{DM}, \texttt{card\_hopPairs\_at\_le}, \texttt{costB\_le}, \texttt{costM\_arith}, \texttt{lo\_le\_mwidth}, \texttt{mExp}, \texttt{mK}, \texttt{mK0}, \texttt{mw}, \texttt{mwidth}, \texttt{mwidth\_dvd}, \texttt{mwidth\_le}, \texttt{sum\_sl\_le}, \texttt{sum\_swev\_le}, \texttt{theoremM}, \texttt{theoremM\_owner}, \texttt{wexp}.\par}

\paragraph{\texttt{BoundedNesting.TheoremSH}}
{\small\raggedright \texttt{Far3}, \texttt{MinCE.false\_SH}, \texttt{exists\_rho3}, \texttt{finrank\_ker\_add\_card\_le\_matching}, \texttt{finrank\_ker\_le\_barrier\_SH}, \texttt{finrank\_ker\_le\_matching\_rho3}, \texttt{finrank\_ker\_le\_matching\_rho3\_transpose}, \texttt{finrank\_ker\_transpose}, \texttt{finrank\_subCode\_le\_half\_SH}, \texttt{recTriangleFree\_transpose}, \texttt{rho3}, \texttt{sep\_of\_far3}, \texttt{star4\_of\_degree\_le\_four}, \texttt{theoremSH}, \texttt{theoremSH\_graph}, \texttt{theoremSH\_transpose}.\par}

\paragraph{\texttt{BoundedNesting.TurnCounter}}
{\small\raggedright \texttt{turnCounter\_cost}.\par}

\paragraph{\texttt{BoundedNesting.TwistPairing}}
{\small\raggedright \texttt{blockMap\_sheetEquiv}, \texttt{excess\_lift2}, \texttt{finrank\_ker\_blockMap}, \texttt{finrank\_ker\_lift2}, \texttt{lift2}, \texttt{mem\_range\_iff\_dotProduct\_ker\_transpose}, \texttt{rate\_lift2\_le}, \texttt{sheetEquiv}, \texttt{twistPairing}, \texttt{twistPart}.\par}

\paragraph{\texttt{BoundedNesting.WalkDag}}
{\small\raggedright \texttt{allHops}, \texttt{descentsOf}, \texttt{labelCut\_depth}, \texttt{labelCut\_depth\_walkDag}, \texttt{violatedOf}.\par}

\section{Building and verifying the formalization}\label{app:verify}

The Lean project is the directory \leanfile{lean/} next to this paper. It uses Lean and Mathlib, both at version v4.34.1
(\leanfile{lean-toolchain}, \leanfile{lakefile.toml}, \leanfile{lake-manifest.json}), and requires the library of the fourth
paper~\cite{FourthPaper2026} through a relative path; that library imports the libraries of the second and first papers and the
repository's shared library \texttt{lean-machine}. Keep the repository's directory layout, so that these sibling directories stay
in place.

\paragraph{Build.} With \texttt{elan} installed, run in \leanfile{lean/}:
\begin{center}
\texttt{lake exe cache get \&\& lake build}
\end{center}
The first command downloads the compiled Mathlib; the second compiles the library and the parts of the earlier libraries that it
uses. Before the first build, fetch the second paper's pinned external dependency as its README describes
(\texttt{python3 verification/fetch\_openai\_metric.py}, run in that paper's directory).

\paragraph{Axioms.} The file \leanfile{PaperAxioms.lean}, generated from the \LaTeX{} sources of this paper by
\leanfile{tools/lean_index.py}, runs \texttt{\#print axioms} on every declaration of this library named here (\cref{app:index}):
\begin{center}
\texttt{lake env lean PaperAxioms.lean}
\end{center}
Every declaration depends at most on \texttt{propext}, \texttt{Classical.choice} and \texttt{Quot.sound}. No Lean file of this
library uses \texttt{sorry}, \texttt{admit} or \texttt{native\_decide}, or declares an \texttt{axiom}. The script regenerates both
the index of \cref{app:index} and \leanfile{PaperAxioms.lean}; with the option \texttt{--index-only} it writes only the index.

\paragraph{The computational check of \cref{thm:eqcrux}(f).} The script \leanfile{tools/clebsch_lift.py}, which needs only
Python~3,
builds the $16$-sheeted lift of the Clebsch graph described at its top and prints the rank and the nullity of the lift of $I+A$
over~$\F_2$:
\begin{center}
\texttt{python3 tools/clebsch\_lift.py}
\end{center}
It reports rank $246$ and nullity~$10$. The computation is exact elimination over~$\F_2$ and is not checked in Lean.

\paragraph{Where things are.} \leanfile{Owner.lean} (walks, upper hops, owners, entries and crossings; the owner lemma),
\leanfile{Labels.lean} (cuts are labels), \leanfile{OneWalk.lean} (hops, lines, strips, the strip label and Lemma~W),
\leanfile{TurnCounter.lean} (the label arithmetic of the turn counter, from the fourth paper's Lemma~Zq),
\leanfile{Clock.lean} (abstract family data, the family clock and the comparison of clocks), \leanfile{ClockCost.lean} (the
cost of the family clock, summed over offsets: the bounds of the proof of Theorem~F), \leanfile{Families.lean} (sides,
families, spans, heights, depths and virtual heads of the hops of walks; the structure of families), and
\leanfile{TheoremF.lean} (owner-load, owners, the bound $\DF$ and Theorem~F).


\bibliographystyle{plainurl}
\bibliography{references}

\end{document}
